\section{Analysis of the Pauli basis test}
\label{sec:qld-analysis}
In this appendix we give a proof of \Cref{thm:pauli}. For
convenience, we repeat the statement of the theorem here.

\begin{theorem}
  \label{thm:pauli-appendix}
There exists a function 
\[\delta_{\qld}(\eps,m,d,q) = a(md)^a(\eps^b + q^{-b} + 2^{-bmd})\]
 for universal constants $a \geq 1$ and $0 < b < 1$ such that the following holds. For all admissible parameter tuples $\qldparams = (q,m,d)$ and for all strategies $\strategy = (\ket{\psi},A,B)$ for the game $\game^\pauli_\qldparams$ that succeed with probability at least $1 - \eps$, there exist local isometries $\phi_\alice: \mH_\alice \to \mH_{\alice'} \otimes \mH_{\alice''},\phi_\bob: \mH_\bob \to \mH_{\bob'} \otimes \mH_{\bob''}$ (where $\ket{\psi} \in \mH_\alice \otimes \mH_\bob$ and $\mH_{\alice''},\mH_{\bob''} \cong (\C^q)^{\otimes M}$ with $M = 2^m$) and a state $\ket{\aux} \in \mH_{\alice'} \otimes \mH_{\bob'}$ such that 
\begin{enumerate}
    \item $\left \| \phi_\alice \otimes \phi_B \ket{\psi} - \ket{\aux} \otimes \ket{\EPR_q}^{\otimes M}  \right \| \leq \delta_{\qld}(\eps,m,d,q)$,
    \item Letting $\tilde{A}^x_a = \phi_A\, A^x_a \, \phi_A^\dagger$ and $\tilde{B}^y_b = \phi_B B^y_b \phi_B^\dagger$, we have for $W \in \{X,Z\}$
    \begin{gather*}
        \tilde{A}^{(\pauli,W)}_u \otimes I_{\bob'\bob''} \approx_{\delta_{\qld}} (\tau^W_u)_{\alice''} \otimes I_{\alice' \bob' \bob''} \\
        I_{\alice'\alice''} \otimes \tilde{B}^{(\pauli,W)}_u    \approx_{\delta_{\qld}} I_{\alice' \alice'' \bob'} \otimes (\tau^W_u)_{\bob''}\;,
    \end{gather*}
    where the $\approx_{\delta_{\qld}}$ statement holds with respect to the state $\ket{\aux}_{\alice'\bob'} \otimes \ket{\EPR_q}^{\otimes M}_{\alice''\bob''}$ and the answer summation is over $u \in \F_q^M$.
\end{enumerate}
\end{theorem}
We also recall the meaning of the parameters: $q$ is the field size,
$m$ the number of variables, $d$ the degree of the low-degree code,
and $\nqubits=2^m$ is the number of $\ket{\EPR_q}$ states being tested for.

This appendix is structured as follows. First, we establish some preliminary facts and definitions 
in \Cref{sec:qld-prelim}. In the remainder of the appendix, we
show \Cref{thm:pauli-appendix} by showing how to construct a representation of
the Pauli group on the provers' Hilbert space, starting from a
strategy $\strategy$ that succeeds in $\game^\pauli_\qldparams$ with
high probability. This is done in several stages: in
\Cref{sec:commutation}, we use the strategy measurements to define a
set of observables that correspond to a subset of the Pauli $X$ and
$Z$ observables, and show that these
approximately satisfy the associated Pauli group relations. In \Cref{sec:expanding},
we adjoin ancillary registers containing EPR states to the provers'
state $\ket{\psi}$, and define a new set of approximately-commuting
observables acting on the expanded state. In \Cref{sec:combining}, we
combine the approximately-commuting observables into a strategy for
the \emph{classical} low-degree test, and in \Cref{sec:apply-ldt}, we
apply the soundness theorem of the classical low-degree test to
construct a single projective measurement that simultaneously measures all of the
approximately-commuting observables. Finally, in
\Cref{sec:separating}, we use this single simultaneous measurement to
construct an exact representation of the Pauli group, show that it
is close to the original strategy observables on the expanded state, and deduce that the game $\game^\pauli_\qldparams$ is a self-test for the strategy $\strategy^\pauli$.

\subsection{Preliminaries}
\label{sec:qld-prelim}


\begin{definition}
Given integer $d,m$ and a prime power $q$ we denote by $\ideg_{d,m}(\F_q)$ the set of polynomials in $m$ variables
  with \emph{individual} degree at most $d$ in each variable. When the field is clear from context we write $\ideg_{d,m}$. 	
	When
  the number of variables is clear from context we
  write $\ideg_d$. We denote by $\deg_{d}$ the set of polynomials with
  \emph{total} degree at most $d$.
\end{definition}

Recall the low-degree code defined in \Cref{sec:ld-encoding}, in which vectors $h \in \F_q^\nqubits$ get mapped to polynomials $g_h \in \ideg_{d,m}(\F_q)$ such that for every $u \in \{0,1\}^m$, $g_h(u) = h_u$ where the coordinates of $h$ are indexed by binary strings of length $m$. For $x\in \F_q^m$, recall the definition of the vector $\ind_m(x) \in \F_q^\nqubits$ from Section~\ref{sec:ld-encoding}. Recall also the finite field trace $\tr(\cdot)=\tr_{q\to 2}(\cdot):\F_q\to\F_2$, where $q=2^k$ for odd $k$, defined in Section~\ref{sec:subfields}. 



Call a tuple $\omega = (u_\xpt,u_\zpt,r_\xpt,r_\zpt) \in (\F_q^m)^2
\times (\F_q)^2$ \emph{anticommuting} if the quantity 
 \[\gamma\,=\,\tr ((r_\xpt \cdot
\ind_m(u_\xpt)) \cdot (r_\zpt \cdot \ind_m(u_\zpt)))\]
introduced in~\Cref{eq:gamma-value} is not $0$, and \emph{commuting} if $\gamma=0$.

\begin{fact}
\label{fact:omega-anticomm-prob}
	The probability that a tuple $\omega= (u_\xpt,u_\zpt,r_\xpt,r_\zpt)$ sampled uniformly at random from $(\F_q^m)^2 \times (\F_q)^2$
	is anti-commuting is at least 
	\[
		(1 - q^{-m}) \cdot (1 - q^{-1}) \cdot \Big (1 - \frac{md}{q} \Big )  \cdot \frac{1}{2} \geq \Big (1 - \frac{3md}{q} \Big ) \cdot \frac{1}{2}.
	\]
	The probability that $\omega$ is commuting is also at least $\Big(1 - \frac{3md}{q} \Big ) \cdot \frac{1}{2}$.
\end{fact}
\begin{proof}
	First, note that $\ind_m(u_\xpt) = 0$ if and only if $u_\xpt = 0$. If $u_\xpt=0$ then $\omega$ is commuting. Otherwise, if $u_\xpt \neq 0$, then observe that for all $r_\xpt, r_\zpt, u_\zpt$,
	\begin{align*}
		\tr( (r_\xpt \cdot \ind_m(u_\xpt)) \cdot (r_\zpt \cdot
                \ind_m(u_\zpt))) &= \tr( r_\xpt r_\zpt (\ind_m(u_\xpt) \cdot \ind_m(u_\zpt))) \\
								&= \tr( (r_\xpt r_\zpt) \cdot g_{\ind_m(u_\xpt)}(u_\zpt))\;,
	\end{align*}
	where $g_{\ind_m(u_\xpt)}$ is a nonzero polynomial in $\ideg_{d,m}(\F_q)$.
	
	If $r_\xpt = 0$, then $\omega$ is commuting. Condition on $r_\xpt \neq 0$.  Since $g_{\ind_m(u_\xpt)}\neq 0$ has total degree at most $md$, by the Schwartz-Zippel lemma the probability that $g_{\ind_m(u_\xpt)}(u_\zpt)$ is zero over a uniformly random choice of $u_\zpt$ is at most $md/q$. Conditioned on $g_{\ind_m(u_\xpt)}(u_\zpt) \neq 0$, the product $(r_\xpt r_\zpt) \cdot g_{\ind_m(u_\xpt)}(u_\zpt)$ is a uniformly random element of $\F_q$ when $r_\zpt$ is chosen uniformly at random. 
	
	The probability that $u_\xpt \neq 0, r_\xpt \neq 0, g_{\ind_m(u_\xpt)}(u_\zpt) \neq 0$ is at least $(1 - q^{-m}) \cdot (1 - q^{-1}) \cdot (1 - md/q)$ since $u_\xpt, r_\xpt, u_\zpt$ are chosen independently.
	
	Since $q = 2^t$ is an admissible field size, we have that for all $a \in \F_q$, the trace of $a$ is equal to $\sum_i a_i$ where $(a_1,\ldots,a_t)=\kappa(a) \in \F_2^t$ are defined in Section~\ref{sec:subfields}. Thus the probability that a uniformly random element of $\F_q$ has trace $0$ is exactly $1/2$. 
\end{proof}

In this appendix we often use the ``$\approx_\delta$'' notation to specify closeness of general operators (not necessarily POVMs) that are indexed by questions but not answers: given sets of operators $\{A^x \}$ and $\{B^x\}$ indexed by questions $x \in \cal{X}$, we write $A^x \approx_\delta B^x$ to denote \[
	\E_{x \sim \mu} \bra{\psi} (A^x - B^x)^\dagger (A^x - B^x)
        \ket{\psi}\,\leq\, O(\delta)
      \]
where $\mu$ is a question distribution and $\ket{\psi}$ is a state that has been fixed by context. This will be useful for expressing closeness of observables, for example. The following two lemmas relate the closeness between sets of operators and weighted sums of those operators. 
\begin{lemma}
\label{lem:avg-closeness}
Let $\{A^x \}$ and $\{B^x\}$ be operators indexed by some finite set $\mathcal{X}$, and let $\mu$ be a probability distribution over $\mathcal{X}$. Let $\{\alpha_x \}$ be a set of complex numbers such that $|\alpha_x| \leq 1$. Then if $A^x \approx_\delta B^x$ on average over $x$ drawn from $\mu$, then $A \approx_\delta B$, where $A = \E_x \alpha_x A^x$ and $B= \E_x \alpha_x B^x$.
\end{lemma}
\begin{proof}
Expand
\begin{align*}
	\Big\| \E_x \alpha_x \big ( A^x - B^x \big) \ket{\psi} \Big\|^2 &\leq \Big( \E_x \norm{ (A^x -B^x) \ket{\psi}} \Big)^2 \\
	&\leq \E_x \norm{(A^x - B^x) \ket{\psi}}^2 \\
	&\leq \delta\;.
\end{align*}
The first inequality is the triangle inequality and the second inequality is Jensen's inequality. 
\end{proof}

\begin{lemma}
\label{lem:povm-to-obs}
Let $\mathcal{A}$ be a finite set. Let $\{A^x_a \}_{a \in \mathcal{A}}$ and $\{B^x_a\}_{a \in \mathcal{A}}$ be POVMs and let $\{ \alpha_a \}_{a \in \mathcal{A}}$ be a collection of complex numbers on the unit circle. Let $A^x = \sum_a \alpha_a A^x_a$ and $B^x = \sum_a  \alpha_a B^x_a$. Then 
\[
	A^x_a \approx_\delta B^x_a \quad \text{implies} \quad A^x \approx_{\delta'} B^x
\]
for $\delta' = |\mathcal{A}| \cdot \delta$.
\end{lemma}
\begin{proof}
Expand
	\begin{align*}
		\E_x \norm{(A^x - B^x) \ket{\psi}}^2 &= \E_x \Big\|\sum_a \alpha_a (A^x_a - B^x_a) \ket{\psi}\Big\|^2 \\
		&\leq \E_x \Big( \sum_a \norm{ (A^x_a - B^x_a) \ket{\psi}} \Big)^2 \\
		&\leq \E_x |\mathcal{A}| \cdot \sum_a \norm{ (A^x_a - B^x_a) \ket{\psi}}^2 \\
		&\leq |\mathcal{A}| \cdot \delta\;.
	\end{align*}
The first inequality follows from the triangle inequality, the second inequality follows from Cauchy-Schwarz, and the last inequality follows from the assumption that $A^x_a \approx_\delta B^x_a$. 
\end{proof}


The following lemma shows that POVMs that are approximately projective (in the sense of being self-consistent) are close to exact projective measurements. The lemma first appears in~\cite{KV11}; for a self-contained proof, see~\cite{ML20}.

\begin{lemma}[Orthonormalization lemma]\label{lem:ortho}
Let $0\leq \delta\leq 1$. There exists a function $\eta_{ortho}(\delta)\,=\, O\big(\delta^{1/4}\big)$ such that the following holds.  Let $\ket{\psi}$ be a state on $\mH_\alice \otimes \mH_\bob$ where $\mH_\alice,\mH_\bob$ are finite dimensional. Let $\{Q_a\}$ and $\{R_a\}$ be POVM on $\mH_\alice$ and $\mH_\bob$ respectively such that 
\begin{equation}\label{eq:cons-cond} 
Q_a \abc R_a \,.
\end{equation}
Then there exists a projective measurement $\{P_a\}$ such that
\begin{equation}\label{eq:cons-ccl}
  P_a \otimes \Id_\bob \,\approx_{\eta_{ortho}} Q_a \otimes \Id_\bob \;.
\end{equation}
\end{lemma}









 \newcommand{\acom}{\typestyle{Anticom}}

\subsection{Strategies}
\label{sec:commutation}

  Let $\qldparams = (q,m,d)$ be an admissible parameter tuple (Definition~\ref{def:admissible}) and let
  $\strategy = (\ket{\psi},M)$ be a strategy for
  $\game^\pauli_\qldparams$ that succeeds with probability at least $1 - \eps$, where $\ket{\psi}$ is a bipartite state on registers $\alice$ and $\bob$, with associated Hilbert spaces $\mH_\alice$ and $\mH_\bob$ respectively. For notational convenience we use $M$ to denote the operators for both players; it will be clear from context which player's operators we are referring to (for example, we write $M^{(\Point,W),u}_a \otimes \Id_\bob$ to indicate that $M^{(\Point,W),u}_a$ is viewed as an operator acting on register $\alice$).
  
  Using Naimark's
  theorem (as formulated in~\cite[Theorem 5.1]{ML20}), at the cost of
  extending the state $\ket{\psi}$ by adding sufficiently many qubits
  initialized in the $\ket{0}$ state we may assume
  that the measurements in $\strategy$ are projective. Let
  \[ \ket{\psi'} = (\ket{\psi} \otimes \ket{0 \cdots 0} \otimes \ket{0
      \cdots 0})_{\alice \bob} \]
  be the extended state. In the remainder of this section, we
  will work with the projective strategy on the extended state, which we continue labeling $\ket{\psi}$ (instead of $\ket{\psi'}$) for notational convenience.

We introduce several additional notational shorthands. By definition the strategy $\strategy$ contains measurement
  operators for each of the possible question types and content
  summarized in \cref{fig:decider_pauli}, such as $\{M^{(\Point,
    W),y}_a\}_{a\in\F_q}$ for all $y\in\F_q^m$, etc. 
For ``line'' questions (of type $\ALine$ or $\DLine$) we will use
 $\Line$ as a formal placeholder for either type $\ALine$ or $\DLine$;
  which type is meant will be clear from the nature of the
  line. Moreover, we will use $\ell$ to denote the description of
  either an axis-parallel line (given by a pair $\ell=(u_0, s)$) or a
  diagonal line (given by a triple $\ell=(u_0, s, v)$). We also introduce the following collection of observables constructed from the
  measurement operators in the strategy. 	
	For every $u \in
  \F_q^m$, $r \in \F_q$, and $W \in \{X, Z\}$, define
  \begin{equation}
    W^r(u) =\sum_{a\in\F_q} (-1)^{\tr(ar)}  M^{(\Point, W), u}_a\;. \label{eq:qld-strat-obs}
  \end{equation}
	The observable $W^r(u)$ acts on register $\alice$ or $\bob$ depending on whether $M^{(\Point, W), u}_a$ is viewed as being an operator acting on register $\alice$ or $\bob$; this will be made clear from context.
	



Recall the definitions of (anti-)commuting tuple given at the start of Section~\ref{sec:qld-prelim}. 

\begin{lemma}
\label{lem:qld-win-implications}
Assume that $\strategy$ succeeds in $\game^\pauli_\qldparams$ with probability at least $1-\eps$, for some $\eps\geq 0$. Then the following hold: 
	\begin{enumerate}
		\item \label{enu:qld-cons} (\textbf{Consistency check}) On average over $(\tvar,x)$ sampled from the question distribution of $\game^\pauli_\qldparams$, we have that $M^{\tvar,x}_a \otimes \Id_\bob \simeq_\eps \Id_\alice \otimes M^{\tvar,x}_a$.
		\item \label{enu:qld-ld} (\textbf{Low-degree check})
                  For all $W \in \{X,Z\}$, on average over $(\line,u)$
                  drawn from the line-point distribution (see
                  Definition~\ref{def:line-point-dist}), we
                  have $M^{(\Line,W),\line}_{[\eval_u(\cdot) = a]} \otimes \Id_\bob \simeq_\eps \Id_\alice \otimes M^{(\Point,W),u}_a$, where the answer summation is over $a \in \F_q$. 
		
	\item \label{enu:qld-pauli-basis-cons} (\textbf{Pauli basis consistency check}) For all $W \in \{X,Z\}$, on average over $u \in \F_q^m$, we have $M^{(\Point,W),u}_a \otimes \Id_\bob \simeq_\eps \Id_\alice \otimes M^{(\Pauli,W)}_{[g_h(u)=a]}$ where $M^{(\Pauli,W)}_{[g_h(u)=a]} = \sum_{h \in \F_q^{\nqubits} : g_h(u) = a} M^{(\Pauli,W)}_h$ and $g_h$ is the low-degree encoding of $h$. 
	
		\item \label{enu:qld-comm} (\textbf{Commutation check}) For all $W \in \{X,Z\}$, on average over commuting $\omega = (u_\xpt,u_\zpt,r_\xpt,r_\zpt)$, we have $M^{(\Pair,W),\omega}_a \otimes \Id_\bob \simeq_\eps \Id_\alice \otimes M^{\Pair,\omega}_{[\beta_W = a]}$, where the answer summation is over $a \in \F_2$.
		\item \label{enu:qld-comm-cons} (\textbf{Commutation consistency check}) For all $W \in \{X,Z\}$, on average over commuting $\omega = (u_\xpt,u_\zpt,r_\xpt,r_\zpt)$, we have $M^{(\Point,W),u_W}_{[\tr(\cdot r_W) = a]} \otimes \Id_\bob \simeq_\eps \Id_\alice \otimes M^{(\Pair,W),\omega}_{a}$, where the answer summation is over $a\in \F_2$.
		
		\item\label{enu:qld-ms} (\textbf{Magic square check}) For any anticommuting $\omega = (u_\xpt,u_\zpt,r_\xpt,r_\zpt)$ let $\Lambda_\omega$ be the value of the strategy $\Big (\ket{\psi}, \{M^{\Constraint_i,\omega}_{\alpha_1,\alpha_2,\alpha_3}\} \cup \{M^{\Variable_j,\omega}_a \} \Big )$ in the game $\game^\ms$. Then 
		\[
			\E_{\omega \, : \, \tr((r_\xpt \cdot
                          u_\xpt) \cdot (r_\zpt \cdot u_\zpt)) \neq 0} \Lambda_\omega = 1 - O(\eps)\;.
		\]

	\item \label{enu:qld-ms-cons} (\textbf{Magic square consistency check}) For all $W \in \{X,Z\}$, on average over anticommuting $\omega = (u_\xpt,u_\zpt,r_\xpt,r_\zpt)$,
	\begin{gather}
		M^{(\Point,X),u_\xpt}_{[\tr(\cdot r_\xpt) = a]} \otimes \Id_\bob \simeq_\eps \Id_\alice \otimes M^{\Variable_1,\omega}_{a} \;,\\
		M^{(\Point,Z),u_\zpt}_{[\tr(\cdot r_\zpt) = a]} \otimes \Id_\bob \simeq_\eps \Id_\alice \otimes M^{\Variable_5,\omega}_{a}\;,
	\end{gather}	
	where the answer summations are over $\F_2$.
	
	\end{enumerate}
	Furthermore, all of the approximations above hold with ``$\approx_\eps$'' instead of ``$\simeq_\eps$'', and they also hold with the tensor factors interchanged.
\end{lemma}

\begin{proof}
	Since the question types are sampled according to the type graph $G^\Pauli$ of
  finite size in \cref{fig:type-graph-pauli} in the game
  $\game^\pauli_\qldparams$, the probability of any of the subtests of
  $\decider^\pauli_\qldparams$ (described in \cref{fig:decider_pauli})
  being executed is at least some universal constant.
  Thus the consistency conditions expressed in Items
  \ref{enu:qld-cons},~\ref{enu:qld-ld},~\ref{enu:qld-pauli-basis-cons} of the
  lemma follow directly from the corresponding tests in
  \cref{fig:decider_pauli} (where \Cref{enu:qld-ld} follows from the
  consistency test that is a subtest of $\decider^\ld_\ldparams$) and the fact
  that the strategy $\strategy$ is assumed to succeed with probability at least
  $1 - \eps$ in $\game^\pauli_\qldparams$.
				
  For the remaining items, recall that by Fact~\ref{fact:omega-anticomm-prob}
  the probability that a tuple $\omega=(u_\xpt,u_\zpt,r_\xpt,r_\zpt)$ sampled as
  question content for a question of type $\Pair$, $(\Pair,W)$, $\Constraint_i$
  or $\Variable_j$ is commuting is at least $(1 - 3md/q) \cdot (1/2)$, which is
  at least $1/4$ assuming $6md \leq q$. We claim this assumption holds
  without loss of generality. Indeed, observe that the function
  $\delta$ in \Cref{thm:pauli-appendix} is monotonically increasing in $a$, so without loss of generality we may assume that $a
  \geq 6$.  Then it follows that we may assume that $6md \leq q$
  without loss of generality, since otherwise, $\delta \geq a md /q
  \geq 1$ and the conclusion of the theorem holds trivially.
  The same lower bound holds for the probability that $\omega$ is anticommuting.
  Therefore, the consistency conditions in
  Items~\ref{enu:qld-comm},~\ref{enu:qld-comm-cons} and~\ref{enu:qld-ms-cons} of
  the lemma also follow directly from the corresponding tests in
  \cref{fig:decider_pauli}.
  Finally, \Cref{enu:qld-ms} in the lemma follows from the fact that the Magic
  Square check in \cref{fig:decider_pauli} is executed with constant probability
  over the choice of a pair of questions.
				
That all of the approximations in the Lemma statement also hold with
$\simeq_\eps$ replaced by $\approx_\eps$ is immediate from Fact~\ref{fact:agreement}.
\end{proof}

\begin{lemma}
\label{len:qld-win-implications-obs}The following hold for the observables defined in~\eqref{eq:qld-strat-obs}. 
	\begin{itemize}
		\item For all $W \in \{X,Z\}$, $r \in \F_q$, on average over uniformly random $u \in \F_q^m$ we have
		\begin{equation}
		\label{eq:pts-obs-consistency}
			W^r(u) \otimes \Id_\bob \approx_\eps \Id_\alice \otimes W^r(u).
		\end{equation}
		\item On average over uniformly random $\omega = (u_\xpt,u_\zpt,r_\xpt,r_\zpt)$, 
		\begin{align}
			X^{r_\xpt}(u_\xpt)
                        Z^{r_\zpt}(u_\zpt) \otimes \Id_\bob\notag\\
												&       \approx_{\sqrt{\eps}} (-1)^{\tr ((r_\xpt \cdot
                          \ind_m(u_\xpt)) \cdot (r_\zpt \cdot \ind_m(u_\zpt)))} \, Z^{r_\zpt}(u_\zpt) X^{r_\xpt}(u_\xpt) \otimes \Id_\bob	\;.	\label{eq:pts-obs-commutation}
		\end{align}
	\end{itemize}
	Furthermore, all approximations above hold with the tensor factors interchanged.
\end{lemma}
\begin{proof}
We first establish~\eqref{eq:pts-obs-consistency}. Fix a $W \in \{X,Z\}$ and $r \in \F_q$. From item~\ref{enu:qld-cons} of Lemma~\ref{lem:qld-win-implications} and the data processing inequality (Fact~\ref{fact:data-processing}), using that the question type $(\Point,W)$ has constant probability of being sampled by $\sampler^\pauli$, we get that on average over $u \in \F_q$,
\[
	M^{(\Point,W),u}_{[\tr(\cdot r) = a]} \otimes \Id_\bob \approx_\eps \Id_\alice \otimes M^{(\Point,W),u}_{[\tr(\cdot r) = a]}\;,
\]
where the answer summation is over $a \in \F_2$. Equation~\eqref{eq:pts-obs-consistency} follows directly from~\Cref{lem:povm-to-obs}. 

We now establish~\eqref{eq:pts-obs-commutation}. We have that on average over {commuting} $\omega = (u_\xpt,u_\zpt,r_\xpt,r_\zpt)$ each of the approximate consistency relations hold, for any $W\in\{X,Z\}$:
\begin{align*}
	&& M^{(\Point,W),u_W}_{[\tr(\cdot r_W)=a]} \otimes \Id_\bob &\simeq_\eps \Id_\alice \otimes M^{(\Pair,W),\omega}_a && \text{(Item~\ref{enu:qld-comm-cons} of Lemma~\ref{lem:qld-win-implications})}\\
	&& &\simeq_\eps M^{\Pair,\omega}_{[\beta_W = a]} \otimes \Id_\bob && \text{(Item~\ref{enu:qld-comm})} \\
	&& &\simeq_\eps \Id_\alice \otimes M^{\Pair,\omega}_{[\beta_W = a]}\;. && \text{(Item~\ref{enu:qld-cons} and Fact~\ref{fact:data-processing})}
\end{align*}
Fact~\ref{fact:agreement} then implies that on average over commuting $\omega$, 
\begin{equation}\label{eq:lc-11}
	M^{(\Point,W),u_W}_{[\tr(\cdot r_W)=a]} \otimes \Id_\bob \approx_\eps \Id_\alice \otimes M^{\Pair,\omega}_{[\beta_W = a]}\;.
\end{equation}
Since
$\{M^{\Pair,\omega}_{\beta_X,\beta_Z} \}$ is projective,
\Cref{eq:lc-11} puts us in
a position to apply Lemma~\ref{lem:commutation-analysis}, setting
$A^{x}_{a,b} = M^{(\Point, X), u_X}_{[\tr(\cdot r_X) = b]}$,
$C^{x}_{a,c} = M^{(\Point, Z), u_Z}_{[\tr(\cdot r_Z) = c]}$, and
$B^{x}_{a,b,c} = M^{\Pair, \omega}_{b,c}$. It follows that
\begin{equation}
\label{eq:qld-obs-comm}
	 M^{(\Point,X),u_\xpt}_{[\tr(\cdot r_\xpt)=b]} M^{(\Point,Z),u_\zpt}_{[\tr(\cdot r_\zpt)=c]} \otimes \Id_\bob \approx_\eps  M^{(\Point,Z),u_\zpt}_{[\tr(\cdot r_\zpt)=c]} M^{(\Point,X),u_\xpt}_{[\tr(\cdot r_\xpt)=b]}  \otimes \Id_\bob
\end{equation}
on average over commuting $\omega$. Then for every $\omega = (u_\xpt,u_\zpt,r_\xpt,r_\zpt)$, since the measurements $\{M^{\tvar,x}_a\}$ are projective we get that
\begin{align}
	X^{r_\xpt}(u_\xpt) Z^{r_\zpt}(u_\zpt) &= \Big ( \Id - 2 M^{(\Point,X),u_\xpt}_{[\tr(\cdot r_\xpt)=1]} \Big ) \Big ( \Id - 2 M^{(\Point,Z),u_\zpt}_{[\tr(\cdot r_\zpt)=1]}  \Big ) \notag\\
	&= I - 2 M^{(\Point,X),u_\xpt}_{[\tr(\cdot r_\xpt)=1]} - 2 M^{(\Point,Z),u_\zpt}_{[\tr(\cdot r_\zpt)=1]} + 4 M^{(\Point,X),u_\xpt}_{[\tr(\cdot r_\xpt)=1]} M^{(\Point,Z),u_\zpt}_{[\tr(\cdot r_\zpt)=1]}\;.\label{eq:lc-12}
\end{align}
Applying~\eqref{eq:qld-obs-comm} with $b = c = 1$ to commute the
measurement operators in the the fourth term
in~\eqref{eq:lc-12}, and performing the same steps in reverse, we deduce that on average over commuting $\omega$,
\[
	X^{r_\xpt}(u_\xpt) Z^{r_\zpt}(u_\zpt) \approx_\eps  Z^{r_\zpt}(u_\zpt) X^{r_\xpt}(u_\xpt).
\]
This shows the ``commuting'' part of~\eqref{eq:pts-obs-commutation}.

Now we consider the case that $\omega$ is anticommuting. Item~\ref{enu:qld-ms-cons} of Lemma~\ref{lem:qld-win-implications}, combined with \Cref{lem:povm-to-obs}, implies that the following two approximations hold on average over anticommuting $\omega$:
\begin{gather}
	X^{r_\xpt}(u_\xpt) \otimes \Id_\bob \approx_\eps \Id_\alice \otimes M^{\Variable_1,\omega} \;,\label{eq:lc-11a}\\
	Z^{r_\zpt}(u_\zpt) \otimes \Id_\bob \approx_\eps \Id_\alice \otimes M^{\Variable_5,\omega}\;,\label{eq:lc-11b}
\end{gather}
where for $j\in\{1,5\}$, $M^{\Variable_j,\omega} = M^{\Variable_j,\omega}_0-M^{\Variable_j,\omega}_1$.

For all anticommuting $\omega$, let $\eps_\omega = 1 -
\Lambda_\omega$, where $\Lambda_\omega$ is defined in
item~\ref{enu:qld-ms} of Lemma~\ref{lem:qld-win-implications}. From
item~\ref{enu:qld-ms} of Lemma~\ref{lem:qld-win-implications} we get
that $\E_{\omega \text{ anticomm.}} \eps_\omega \leq
O(\eps)$. For each anticommuting $\omega$,
Theorem~\ref{thm:ms-rigidity} implies that
\[ \Id_\alice \otimes M^{\Variable_1,\omega} M^{\Variable_5,\omega}
  \approx_{\sqrt{\eps_\omega}} -\Id_\alice \otimes M^{\Variable_5,\omega}
  M^{\Variable_1,\omega}\;. \]
Thus, averaging over anticommuting $\omega$, and applying Jensen's
inequality to the definition of $\approx_{\eps}$, we obtain the same
anticommutation relation with error $ \E_{\omega \text{anticomm.}}
\eps_{\omega}^{1/2} \leq (\E_{\omega \text{anticomm.}}
\eps_\omega)^{1/2} = O(\eps^{1/2})$. In summary, on average over anticommuting $\omega$, 
\begin{equation}
\label{eq:qld-implication-ms-anticomm}
\Id_\alice \otimes M^{\Variable_1,\omega} M^{\Variable_5,\omega} \approx_{\sqrt{\eps}} -\Id_\alice \otimes M^{\Variable_5,\omega} M^{\Variable_1,\omega}\;.
\end{equation}
Using \ref{fact:add-a-proj} twice we get from~\eqref{eq:lc-11a} and~\eqref{eq:lc-11b} that on average over anticommuting $\omega$,
\begin{gather}
	X^{r_\xpt}(u_\xpt) Z^{r_\zpt}(u_\zpt) \otimes \Id_\bob \approx_\eps \Id_\alice \otimes M^{\Variable_1,\omega}  M^{\Variable_5,\omega}\;, \\
	Z^{r_\xpt}(u_\xpt) X^{r_\zpt}(u_\zpt) \otimes \Id_\bob \approx_\eps \Id_\alice \otimes M^{\Variable_5,\omega}  M^{\Variable_1,\omega}\;.
\end{gather}
Using~\eqref{eq:qld-implication-ms-anticomm} we get that on average over anticommuting $\omega$, we have
\[
X^{r_\xpt}(u_\xpt) Z^{r_\zpt}(u_\zpt) \otimes \Id_\bob \approx_{\sqrt{\eps}} - Z^{r_\zpt}(u_\zpt)X^{r_\xpt}(u_\xpt)  \otimes \Id_\bob\;.
\]
This shows the ``anticommuting'' part of~\eqref{eq:pts-obs-commutation}.
\end{proof}

\subsection{Expanding the Hilbert space and defining commuting observables}
\label{sec:expanding}

\paragraph{Expanded state.} We introduce registers $\alice', \alice'', \bob', \bob''$ that are each isomorphic to $(\C^q)^{\otimes \nqubits}$. Define the state \begin{equation}\label{eq:def-psihat}
	\ket{\hat{\psi}}_{\alice \alice' \alice'' \bob \bob' \bob''} = (\ket{\psi})_{\alice \bob} \otimes \ket{\EPR_q}^{\otimes \nqubits}_{\alice' \alice''} \otimes \ket{\EPR_q}^{\otimes \nqubits}_{\bob' \bob''}\;,
\end{equation}
where we have added maximally entangled states in registers
$\alice' \alice''$ and $\bob' \bob''$; recall also that the state $\ket{\psi}$ is already implicitly  appended with sufficiently many ancilla qubits initialized in the $\ket{0}$ state. 
For the remainder of the Appendix all approximations will be evaluated with respect to this state.

\paragraph{Expanded observables.} For all $W \in \{X,Z\}$, $r \in \F_q$, and $u \in \F_q^m$ define the following observable $\hat{W}^r(u)$: \begin{equation}\label{eq:lc-23}
	\hat{W}^r(u) = W^r(u) \otimes \tau^W(r \cdot \ind_m(u))\;,
\end{equation}
where $W^r(u)$ is the observable defined in \Cref{eq:qld-strat-obs} and $\tau^W(r \cdot \ind_m(u))$ is the generalized Pauli observable (see \Cref{sec:generalized-pauli}) acting on $(\C^q)^{\otimes \nqubits}$. 

For $a\in \F_q$ define corresponding projections
\[
	\hat{M}^{(\Point,W),u}_a = \E_{r \in \F_q} (-1)^{\tr(ar)} \, \hat{W}^r(u)\;.
\]
Note that $\hat{M}^{(\Point,W),u}_a$ can be equivalently
written as a sum of projectors by applying Fourier
identities. Specifically, for $W \in \{X, Z\}$, $u \in \F_q^{m}$, and
$a \in \F_q$, define the projector
\begin{equation} \tau^{W, u}_{a} =
  \tau^{W}_{[ g_{\cdot}(u) = a]} = \sum_{h \in \F_q^{\nqubits}: h
    \cdot \ind_m(u) = a} \tau^{W}_{h}\;. \label{eq:qld-point-obs-def} \end{equation}
Then we can write $\hat{M}^{(\Point, W), u}_a$ as follows:
\begin{align*}
  \hat{M}^{(\Point,W),u}_a
  &= \E_{r \in \F_q} (-1)^{\tr(ar)} \,
    W^r(u) \otimes \tau^W(r
    \cdot \ind_m(u)) \\
    &= \E_{r \in \F_q}
    (-1)^{\tr(ar)} \, \Big(\sum_{a' \in \F_q}(-1)^{\tr (a' r)} M^{(\Point,W),u}_{a'}
    \Big) \otimes \Big(
    \sum_{h \in \F_q^n}
    (-1)^{\tr( r( h \cdot \ind_m(u)))} \tau^{W}_{h} \Big) \\
  &= \E_{r \in \F_q}
    (-1)^{\tr(ar)} \, \Big(\sum_{a' \in \F_q}(-1)^{\tr (a' r)} M^{(\Point,W),u}_{a'}
    \Big) \otimes \Big(
    \sum_{a'' \in \F_q}
    (-1)^{\tr(a'' r)} \tau^{W,u}_{a''} \Big) \\
  &=	\sum_{a',a'' \in \F_q} \E_{r \in \F_q} (-1)^{\tr((a + a' + a'')r)} \, M^{(\Point,W),u}_{a'}  \otimes \tau^{W,u}_{a''} \\
  &= \sum_{\substack{a',a'' \in \F_q \\ a' + a'' = a}}  \, M^{(\Point,W),u}_{a'}  \otimes \tau^{W,u}_{a''}\;,
\end{align*}
where in the second line we have used \Cref{eq:qld-strat-obs} on the
first tensor factor and \Cref{eq:pauli-obs-proj} on the second, and in the third line we
have used \Cref{eq:qld-point-obs-def}.
For a fixed $u\in\F_q^m$, $r \in \F_q$ and $b \in \F_2$, define a projection
\begin{equation}
	\hat{M}^{(\Point,W),u,r}_b = \sum_{a \in \F_q : \tr(ar) = b}
        \hat{M}^{(\Point,W),u}_a\;, \label{eq:qld-def-mptur}
\end{equation}
which can be equivalently written as
\begin{equation}\label{eq:lc-22}
\hat{M}^{(\Point,W),u,r}_b = \frac{1}{2} \Big ( \Id + (-1)^b \hat{W}^r(u) \Big)\;.
\end{equation}

The registers that the operators $\hat{W}^r(u)$,
$\hat{M}^{(\Point,W),u}_a$, and $\hat{M}^{(\Point,W),u,r}_b$ act on
depend on context. For example if $W^r(u)$ is viewed as an operator
acting on register $\alice$, then we will view $\hat{W}^r(u)$ as
acting on registers $\alice \alice'$ (with the $\tau^W$ observable
acting on register $\alice'$). We generally indicate via subscripts
which registers the operators act on (for example we will write
$(\hat{W}^r(u))_{\alice \alice'}$ or $(\hat{W}^r(u))_{\bob
  \alice''}$).



\paragraph{Partitioning the registers, and symmetries}
In the remainder of this section as well as \Cref{sec:combining,sec:apply-ldt}, we will partition the six registers
$\alice \alice' \alice'' \bob \bob' \bob''$ into
two ``parties'' in two different ways. The first is to consider
$\alice \alice'$ together to form  a single party, and $\bob \alice''$
to form the second party. The second is to consider $\bob \bob'$ to
form the first party, and $\alice \bob''$ to form the second
party. Thanks to the symmetry of the Pauli basis test under exchanging the
two players, as well as the symmetry of the ancilla states in
registers $\alice' \alice''$ and $\bob' \bob''$, every ``bipartite'' consistency relation we derive between two
operators $M, N$ will hold for both of the two bipartitioning schemes,
and with the first and second party interchanged. That is, whenever we
derive a bipartite relation of the form
\[ M_{\alice \alice'} \approx_\delta N_{\bob \alice''}, \]
the same steps, with the registers appropriately changed, will also
yield
\begin{align*}
  N_{\alice \alice'} &\approx_{\delta} M_{\bob \alice''}, \\
  M_{\bob \bob'} &\approx_{\delta} N_{\alice \bob''}, \\
  N_{\bob \bob'} &\approx_{\delta} M_{\alice \bob''}.
\end{align*}

Similarly, whenever we derive a single-party relation of the form
$M_{\alice \alice'} \approx N_{\alice \alice'}$, the same steps with
the registers appropriately changed will also yield the same relation
on $\bob\alice''$, $\bob \bob'$ and $\alice \bob''$. We refer to these
relations as \emph{symmetric equivalents}.

\begin{lemma}
\label{lem:qld-comm-cons}
	There exists a function $\delta_{\acom}(\eps) = \poly(\eps)$ such that the following holds. 


\begin{enumerate}
	\item (\textbf{Self-consistency}) For all $W \in \{X,Z\}$, on average over $r \in \F_q$ and $u\in \F_q^m$, we have 
	\[
		(\hat{M}^{(\Point,W),u}_a)_{\alice \alice'} \approx_\eps (\hat{M}^{(\Point,W),u}_a)_{\bob \alice''}\,.
		\]
	\item (\textbf{Approximate commutation}) On average over
          uniformly random $\omega = (u_\xpt,u_\zpt,r_\xpt,r_\zpt)$,
          we have 
	\[
		(\hat{M}^{(\Point,\xpt),u_\xpt,r_\xpt}_{b}
                \hat{M}^{(\Point,\zpt),u_\zpt,r_\zpt}_{b'})_{\alice
                  \alice'} \approx_{\delta_{\acom}}
                (\hat{M}^{(\Point,\zpt),u_\zpt,r_\zpt}_{b'}
                \hat{M}^{(\Point,\xpt),u_\xpt,r_\xpt}_{b})_{\alice
                  \alice'} \;.
	\]
      \end{enumerate}
      Furthermore, all symmetric equivalents of these approximations
      also hold.
\end{lemma}

\begin{proof}
Item 1, the self-consistency of $\{\hat{M}^{(\Point,W),u}_a \}$, follows from the fact that the original points measurements $\{ M^{(\Point,W),u}_a \}$ are approximately self-consistent between registers $\alice$ and $\bob$ in $\ket{\hat{\psi}}$ (item~\ref{enu:qld-cons} in Lemma~\ref{lem:qld-win-implications}), the Pauli projections $\{\tau^W_a(\ind_m(u))\}$ are perfectly self-consistent between registers $\alice'$ and $\alice''$, and~\Cref{fact:data-processing}. 
	
	We now establish the approximate commutation relations. On average over uniformly random $\omega = (u_\xpt,u_\zpt,r_\xpt,r_\zpt)$, we have
	\begin{align*}
	&\hat{M}^{(\Point,\xpt),u_\xpt,r_\xpt}_{b} \hat{M}^{(\Point,\zpt),u_\zpt,r_\zpt}_{b'}  \label{eq:qld-m-x-m-z} \\
	&= \frac{1}{4} \Big ( \Id + (-1)^b \hat{X}^{r_\xpt}(u_\xpt) \Big ) \Big ( \Id + (-1)^{b'} \hat{Z}^{r_\zpt}(u_\zpt) \Big ) \\
	&= \frac{1}{4} \Big ( \Id + (-1)^b \hat{X}^{r_\xpt}(u_\xpt) + (-1)^{b'} \hat{Z}^{r_\zpt}(u_\zpt) + (-1)^{b + b'} \hat{X}^{r_\xpt}(u_\xpt) \hat{Z}^{r_\zpt}(u_\zpt) \Big)\;,
	\end{align*}
	where the first equality follows from~\eqref{eq:lc-22}.	
	Note that
	\begin{align*}
	\hat{X}^{r_\xpt}(u_\xpt)
         & \hat{Z}^{r_\zpt}(u_\zpt) \\
				&=
                                             X^{r_\xpt}(u_\xpt)
                                             Z^{r_\zpt}(u_\zpt)
                                             \otimes \tau^X(r_\xpt
                                             \cdot \ind_m(u_\xpt))
                                             \tau^Z(r_\zpt \cdot \ind_m(u_\zpt)) \\
	&\approx_{\sqrt{\eps}} Z^{r_\zpt}(u_\zpt)
   X^{r_\xpt}(u_\xpt)  \otimes \tau^Z(r_\zpt\cdot \ind_m( u_\zpt))
   \tau^X(r_\xpt \cdot \ind_m(u_\xpt)) \\
	&= \hat{Z}^{r_\zpt}(u_\zpt) \hat{X}^{r_\xpt}(u_\xpt)\;, 
	\end{align*}
	where the first equality follows from~\eqref{eq:lc-23} and for the second line we used the approximate (anti-)commutation
        relation~\eqref{eq:pts-obs-commutation} for the
        $X^{r_\xpt}(u_\xpt)$ and  $Z^{r_\zpt}(u_\zpt)$
        observables, as well as the exact (anti-)commutation relations
        for the $\tau^Z(r_\zpt \cdot \ind_m(u_\zpt))$ and
        $\tau^X(r_\xpt \cdot \ind_m(u_\xpt))$ observables. This shows
        the desired approximate commutation relation, with
        $\delta_{\acom} = \sqrt{\eps}$.
\end{proof}

\begin{lemma}\label{lem:qld-comm-line-cons}
There exists a function $\delta_{\Line}(\eps) = \poly(\eps)$ such that the following holds. 
	For all $W \in \{X,Z\}$ and lines $\line$ there exists a projective measurement $\{ \hat{M}^{(\Line,W),\line}_f \}_{f \in \deg_{md}(\line)}$, acting on registers $\alice\alice'$ or $\bob\alice''$, such that:
	\begin{enumerate}
		\item  (\textbf{Self-consistency})\label{enu:qld-comm-line-self-cons}  For all $W \in \{X,Z\}$ and $r \in \F_q$, on average over a uniformly random  pair $(\line, u)$
                drawn from the line-point distribution
                (\Cref{def:line-point-dist}),
		\[
			(\hat{M}^{(\Line,W),\line}_f)_{\alice \alice'} \approx_{\delta_\Line} (\hat{M}^{(\Line,W),\line}_f)_{\bob \alice''}\;.
		\]
		
		\item (\textbf{Consistency with points measurements I}) \label{eq:qld-comm-line-pt-cons}
		On average over a uniformly random pair $(\line, u)$
                drawn from the line-point distribution, 
		\[
			(\hat{M}^{(\Line,W),\line}_f)_{\alice \alice'} \approx_{\delta_\Line} (\hat{M}^{(\Line,W),\line}_f)_{\alice \alice'} \otimes (\hat{M}^{(\Point,W),u}_{f(u)})_{\bob \alice''}\;.
		\]		
		\item (\textbf{Consistency with points measurements II}) \label{eq:qld-comm-line-pt-cons2}
		On average over a uniformly random pair $(\line, u)$
                drawn from the line-point distribution,
		\[
			(\hat{M}^{(\Line,W),\line}_{[\eval_u(\cdot) = a]})_{\alice \alice'} \approx_{\delta_\Line}  (\hat{M}^{(\Point,W),u}_{a})_{\bob \alice''}\;.
		\]
              \end{enumerate}
                    Furthermore, all symmetric equivalents of these approximations
      also hold.
\end{lemma}

\begin{proof}
	Define the operator
	\[
		\hat{M}^{(\Line,W),\line}_f = \sum_{\substack{f',f'' \in \deg_{md}(\line) : \\ f' + f'' = f}} M^{(\Line,W),\line}_{f'} \otimes \tau^{W,\line}_{f''}\;,
	\]
	where the Pauli operator $\tau^{W,\line}_{f''}$ acting on register $\alice'$ is obtained by measuring all $\nqubits$ qudits in the basis $W$ to obtain an outcome $h \in \F_q^\nqubits$ and then returning the restriction $f''$ of $g_h: \F_q^m \to \F_q$ to the line $\line$. Here, $g_h$ is the degree-$d$ low-degree extension of $h$ defined in~\eqref{eq:ld-encoding}. 
	
	\Cref{enu:qld-comm-line-self-cons},	the self-consistency of $\{\hat{M}^{(\Line,W),\line}_f \}$, follows from the fact that the original lines measurement $\{ M^{(\Line,W),\line}_f \}$ is self-consistent between registers $\alice$ and $\bob$ in $\ket{\hat{\psi}}$ (\Cref{enu:qld-cons} in Lemma~\ref{lem:qld-win-implications}), the Pauli projections $\{\tau^{W,\line}_{f''}\}$ are perfectly self-consistent on between registers $\alice'$ and $\alice''$, and~\Cref{fact:data-processing}. 

\Cref{eq:qld-comm-line-pt-cons}, the consistency with the points measurements, follows because on average over $(\line,u)$ drawn from the line-point distribution we have
\begin{align}
	&(\hat{M}^{(\Line,W),\line}_f)_{\alice \alice'} \\
	&= \sum_{\substack{f',f'' \in \deg_{md}(\line) : \\ f' + f'' =
  f}} (M^{(\Line,W),\line}_{f'} \otimes \tau^{W,\line}_{f''})_{\alice
  \alice'} \\
        &= \sum_{\substack{f', f'' \in \deg_{md}(\line): \\ f' + f'' =
  f}} ((M^{(\Line, W), \line}_{f'} \otimes \tau^{W, \line}_{f''})
  \cdot (M^{(\Line, W), \line}_{[\eval_u(\cdot) = f'(u)]} )_{\alice
  \alice'} \label{eq:qld-mhat-line-1} \\
	&\approx_\eps \sum_{\substack{f',f'' \in \deg_{md}(\line) : \\
  f' + f'' = f}} (M^{(\Line,W),\line}_{f'} \otimes
  \tau^{W,\line}_{f''})_{\alice \alice'} \otimes
  (M^{(\Point,W),u}_{f'(u)})_{\bob
  \alice''} \label{eq:qld-mhat-line-2} \\
  	&= \sum_{\substack{f',f'' \in \deg_{md}(\line) : \\ f' + f'' =
  f}} (M^{(\Line,W),\line}_{f'} \otimes \tau^{W,\line}_{f''})_{\alice
  \alice'} \otimes (M^{(\Point,W),u}_{f'(u)} \otimes
  \tau^{W,u}_{f''(u)})_{\bob \alice''} \label{eq:qld-mhat-line-3} \\
	&= \sum_{\substack{f',f'' \in \deg_{md}(\line) : \\ f' + f'' =
  f}} (M^{(\Line,W),\line}_{f'} \otimes \tau^{W,\line}_{f''})_{\alice
  \alice'} \otimes \Big (\sum_{\substack{a', a'' : \\ a' + a'' = f'(u)
  + f''(u)}}  M^{(\Point,W),u}_{a'} \otimes \tau^{W,u}_{a''} \Big
  )_{\bob \alice''} \label{eq:qld-mhat-line-4} \\
	&= (\hat{M}^{(\Line,W),\line}_f)_{\alice \alice'} \otimes
   (\hat{M}^{(\Point,W),u}_{f(u)})_{\bob \alice''}\;. \label{eq:qld-comm-line-pt-cons-eps}
\end{align}
\Cref{eq:qld-mhat-line-1} follows from the projectivity of
$\{\hat{M}^{(\Line, W), \line}_{f}\}$. \Cref{eq:qld-mhat-line-2} follows from the consistency between the line
and point measurements (\Cref{enu:qld-ld} in
Lemma~\ref{lem:qld-win-implications}) together with
\Cref{fact:add-a-proj2}. \Cref{eq:qld-mhat-line-3}
follows from the exact consistency between the $\{ \tau^{W,\line}_{f''} \}$ and $\{
\tau^{W,u}_{f''(u)} \}$ measurements on the $\alice'$ and $\alice''$ registers. \Cref{eq:qld-mhat-line-4} follows because
the consistency of the $\tau$ measurements force $a' = f'(u)$ and $a''
= f''(u)$ . This establishes \Cref{eq:qld-comm-line-pt-cons}, with
approximation error $\eps$.

To show \Cref{eq:qld-comm-line-pt-cons2} we would like to first
apply~\Cref{fact:data-processing} to \Cref{eq:qld-comm-line-pt-cons-eps},
to sum over all functions $f$ that evaluate to the same value at a given
point $u$. Unfortunately, we cannot use the Fact as written, since the
bound in \Cref{eq:qld-comm-line-pt-cons} is in terms of $\approx$, not
$\simeq$, and moreover the right-hand side does not have the form $I
\ot B$. Instead, we show the desired bound by direct calculation from \Cref{eq:qld-comm-line-pt-cons-eps}
\begin{align}
  \eps &\geq \E_{\omega} \sum_{f} \bra{\hat{\psi}} (\hat{M}^{(\Line, W),
  \line}_f )^2 \ot (I - \hat{M}^{(\Point, W), u}_{f(u)})^2 \ket{\hat{\psi}}
  \\
                 &= \E_{\omega} \sum_{f} \bra{\hat{\psi}} \hat{M}^{(\Line, W),
                   \line}_f  \ot (I - \hat{M}^{(\Point, W), u}_{f(u)})
                   \ket{\hat{\psi}} \\
                 &= \E_{\omega} \sum_{a} \sum_{f: f(u) = a} \bra{\hat{\psi}} \hat{M}^{(\Line, W),
                   \line}_f  \ot (I - \hat{M}^{(\Point, W), u}_{a})
                   \ket{\hat{\psi}} \\
                 &=\E_{\omega} \sum_{a}  \bra{\hat{\psi}} \hat{M}^{(\Line, W),
                   \line}_{[\eval_u(\cdot) = a]}  \ot (I - \hat{M}^{(\Point, W), u}_{a})
                   \ket{\hat{\psi}},
\end{align}
where in going to the second line we have used the projectivity of the
measurements $\{\hat{M}^{(\Line, W), \line}_f\}$ and
$\{\hat{M}^{(\Point, W), u}_{a}\}$.
Thus, we have deduced that
\begin{equation}
\label{eq:qld-comm-line-1}
(\hat{M}^{(\Line,W),\line}_{[\eval_u(\cdot) = a]})_{\alice \alice'} \approx_\eps (\hat{M}^{(\Line,W),\line}_{[\eval_u(\cdot) = a]})_{\alice \alice'} \otimes (\hat{M}^{(\Point,W),u}_{a})_{\bob \alice''}\;.
\end{equation}
Since both the $\{ \hat{M}^{(\Line,W),\line}_f\}$ and $\{ \hat{M}^{(\Point,W),u}_a \}$ measurements are projective, we also get 
\begin{align}
\Id &= \sum_a (\hat{M}^{(\Line,W),\line}_{[\eval_u(\cdot) =
      a]})_{\alice \alice'} \otimes I_{\bob \alice''}\notag\\
&= \sum_a \Big((\hat{M}^{(\Line,W),\line}_{[\eval_u(\cdot) =
                                                               a]})_{\alice
                                                               \alice'}
                                                               \ot
                                                               I_{\bob
                                                               \alice''}-(\hat{M}^{(\Line,W),\line}_{[\eval_u(\cdot) = a]})_{\alice \alice'} \otimes (\hat{M}^{(\Point,W),u}_{a})_{\bob \alice''}\Big) \notag\\
&\qquad\qquad + \sum_a (\hat{M}^{(\Line,W),\line}_{[\eval_u(\cdot) = a]})_{\alice \alice'} \otimes (\hat{M}^{(\Point,W),u}_{a})_{\bob \alice''}\;.\label{eq:com-1a}
\end{align}
For projective sub-measurements $\{A_i\}$ and $\{B_i\}$ and a state $\ket{\varphi}$,
\begin{align*}
\sum_i \bra{\varphi}\big(A_i-B_i\big)\ket{\varphi} &= \sum_i \bra{\varphi}A_i(A_i-B_i)\ket{\varphi} + \sum_i \bra{\varphi}(A_i-B_i)B_i\ket{\varphi}\\
&\leq \Big(\sum_i \bra{\varphi}A_i^2\ket{\varphi} \Big)^{1/2}\Big(\sum_i \bra{\varphi}(A_i-B_i)^2\ket{\varphi}\Big)^{1/2} \\
&\qquad\qquad+ \Big(\sum_i \bra{\varphi}(A_i-B_i)^2 \ket{\varphi}\Big)^{1/2}\Big(\sum_i \bra{\varphi}B_i^2\ket{\varphi}\Big)^{1/2}\\
&\leq 2\Big(\sum_i \bra{\varphi}(A_i-B_i)^2 \ket{\varphi}\Big)^{1/2}\;.
\end{align*}
Applying this fact to the projective sub-measurements $\{(\hat{M}^{(\Line,W),\line}_{[\eval_u(\cdot) = a]})_{\alice \alice'}\}_a$ and $\{(\hat{M}^{(\Line,W),\line}_{[\eval_u(\cdot) = a]})_{\alice \alice'} \otimes (\hat{M}^{(\Point,W),u}_{a})_{\bob \alice''}\}$ in~\eqref{eq:com-1a} and using~\eqref{eq:qld-comm-line-1} gives
\begin{equation}
\label{eq:qld-comm-line-2}
	\Id \approx_{\sqrt{\eps}} \sum_a (\hat{M}^{(\Line,W),\line}_{[\eval_u(\cdot) = a]})_{\alice \alice'} \otimes (\hat{M}^{(\Point,W),u}_{a})_{\bob \alice''}\;.
      \end{equation}
      To go from~\eqref{eq:qld-comm-line-2} to \Cref{eq:qld-comm-line-pt-cons2} of the lemma we use another fact about
      projective measurements. For families of projective measurements $\{A^x_i\}$
      and $\{B^x_i\}$ acting on separate registers of a bipartite state
      $\ket{\phi}$, suppose $I \approx_{\delta} \sum_i A^x_i \ot B^x_i$. Then
      \begin{align*}
        \delta &\geq \E_{x} \Big\| \Big(\Id - \sum_i A^x_i \ot B^x_i\Big) \ket{\phi}\Big\|^2  \\
               &= \E_{x} \Big(1 + \bra{\phi} \Big(\sum_i A^x_i \ot
                 B^x_i\Big)^2 \ket{\phi} - 2 \sum_i \bra{\phi} A^x_i \ot B^x_i
                 \ket{\phi} \Big) \\
               &= 1 - \E_{x} \sum_i \bra{\phi} A^x_i \ot B^x_i \ket{\phi}\;,
\end{align*}
      where we used the projectivity in going from the second to
      the third line, to remove the square. Thus
      \[ \E_{x} \sum_{i \neq j} \bra{\phi} A^x_i \ot B^x_j \ket{\phi} \leq
        \delta\;, \]
      and $A^x_i \ot I \simeq_{\delta} I \ot B^x_i$. By
      Item~1 of \Cref{fact:agreement}, this in turn implies that $A^x_i \ot I
      \approx_\delta I \ot B^x_i$. Applying this fact to the projective
      measurements $\{ (\hat{M}^{(\Line, W), \line}_{[\eval_u(\cdot) =
        a]})_{\alice \alice'} \}_a$ and $\{(\hat{M}^{(\Point,W),
        u}_{a})_{\bob \alice''}\}$ in~\eqref{eq:qld-comm-line-2} gives
      \[
	(\hat{M}^{(\Line,W),\line}_{[\eval_u(\cdot) = a]})_{\alice \alice'} \approx_{\sqrt{\eps}} (\hat{M}^{(\Point,W),u}_{a})_{\bob \alice''}.
      \]
      Taking the function $\delta_{\Line}$ in the lemma to be $\sqrt{\eps}$,
      this yields the conclusion of \Cref{eq:qld-comm-line-pt-cons2}
      of the lemma.
\end{proof}

 \subsection{Combining the $X$ and $Z$ measurements}
\label{sec:combining}

In this section our goal is to create a set of combined measurements
that simultaneously measure the approximately-commuting measurements
constructed in the previous section. We do this first for the points.
Recall that $\mH_\alice$ and $\mH_\bob$ denote the Hilbert spaces corresponding to
registers $\alice$ and $\bob$, respectively.

\begin{lemma}\label{lem:qld-4-10}
There exists a function $\delta_Q(\eps) = \poly(\eps)$ such that the following holds. 
	For every $x,z \in \F_q^m$ and $\mH \in \{ \mH_\alice, \mH_\bob\}$ there are projective measurements $\{ \hat{Q}^{x,z}_{a,b} \}_{a,b \in \F_q}$ acting on $\mathcal{H} \otimes (\C^q)^{\otimes \nqubits}$ such that:
	\begin{enumerate}
		\item (\textbf{Self-consistency}) On average over uniformly random $x,z \in \F_q^m$,
		\begin{equation}
			\label{eq:qld-q-self-cons}
			(\hat{Q}^{x,z}_{a,b})_{\alice \alice'}  \approx_{\delta_Q} (\hat{Q}^{x,z}_{a,b})_{\bob \alice''}\,.
		\end{equation}
		
		\item (\textbf{Consistency with $\hat{M}$}) On average over uniformly random $x,z \in \F_q^m$,
		\begin{gather}
			\label{eq:qld-q-cons-m-hat-xz}
				(\hat{Q}^{x,z}_{a,b})_{\alice \alice'} \approx_{\delta_Q} (\hat{M}^{(\Point,X),x}_{a} \hat{M}^{(\Point,Z),z}_b)_{\bob \alice''}\;, \\
							\label{eq:qld-q-cons-m-hat-zx}
				(\hat{Q}^{x,z}_{a,b})_{\alice \alice'} \approx_{\delta_Q} (\hat{M}^{(\Point,Z),z}_b  \hat{M}^{(\Point,X),x}_{a} )_{\bob \alice''}\;.
		\end{gather}
              \end{enumerate}
                    Furthermore, all symmetric equivalents of these approximations
      also hold.
\end{lemma}

The construction of the measurements $\{ \hat{Q}^{x,z}_{a,b} \}$ whose existence is guaranteed in the lemma proceeds in three steps. In the first step, for every $(r,s)\in\F_q^2$ we combine the binary-outcome measurements $\hat{M}^{(\Point,X),x,r}_a$ and $\hat{M}^{(\Point,X),z,s}_b$; this is made possible by their approximate commutation as shown in Lemma~\ref{lem:qld-comm-cons}. In a second step, we show that the measurements thus constructed have a particular ``approximate linearity'' property, when seen as a function of the pair $(r,s)$. This allows us, in the third step, to combine them into a single measurement $\{ \hat{Q}^{x,z}_{a,b} \}$ with outcomes in $\F_q^2\cong \F_2^{2t}$ by applying the following  result from~\cite{NV17}:
	
	\begin{theorem}[Quantum linearity test]
	\label{thm:linearity}
		Let $0 \leq \delta \leq 1$, let $t$ be a positive integer, and let $\ket{\psi}$ be a state in a Hilbert space $\mathcal{H}_\alice \otimes \mathcal{H}_\bob$. Let $\{\mathcal{O}^u \}$ be a set of observables indexed by $u \in \F_2^t$ that act on $\mathcal{H}_\alice$. Suppose that on average over $u, u'$ chosen uniformly from $\F_2^t$, we have the following ``approximate linearity'' holds:
		\[
(\mathcal{O}^u \mathcal{O}^{u'})_{\alice} \approx_\delta (\mathcal{O}^{u+u'})_{\alice}
		\]
		where the $\approx$ statement is taken with respect to the state $\ket{\psi}$. 
Then there exists an extended state $\ket{\psi'}_{\alice \alice' \bob} = \ket{\psi}_{\alice \bob} \otimes \ket{0}_{\alice'}$ in an extended space $\mathcal{H}_{\alice} \otimes \mathcal{H}_\bob \otimes \mathcal{H}_{\alice'}$ and observables $\{ \mathcal{L}^u \}$ acting on $\mathcal{H}_\alice \otimes \mathcal{H}_{\alice'}$ such that for all $u, u' \in \F_2^t$, 
		\[
			\mathcal{L}^u \mathcal{L}^{u'} = \mathcal{L}^{u + u'} \qquad \text{and} \qquad (\mathcal{L}^u)_{\alice \alice'} \approx_\delta (\mathcal{O}^u)_\alice \,.
		\]
		Furthermore, the observable $\mathcal{L}^u = \Id$ when $u = 0$.
	\end{theorem}
		
		We now give the proof of Lemma~\ref{lem:qld-4-10}.
		
\begin{proof}[Proof of Lemma~\ref{lem:qld-4-10}]
For convenience in the proof we slightly abuse the notation $\approx_\delta$ to mean $\approx_{\poly(\eps)}$ for some polynomial function $\delta=\poly(\eps)$ that may differ each time the notation is used. 

\textbf{First step: construction of projective binary measurements $\{ \hat{R}^\omega_{a,b} \}$.}
	For every $\omega = (x,z,r,s)$ and $a,b \in \F_2$, define a measurement operator
	\[
		\hat{R}^\omega_{a,b} = \hat{M}^{(\Point,\zpt),z,s}_{b} \cdot \hat{M}^{(\Point,\xpt),x,r}_{a} \cdot \hat{M}^{(\Point,\zpt),z,s}_{b} \,.
	\]
	It is clear that $\{ \hat{R}^\omega_{a,b} \}_{a,b\in\F_2}$ is a POVM. Observe that on average over uniformly random $\omega$, 
	\begin{align}
	\label{eq:qld-r-2}
	\hat{R}^\omega_{a,b} \approx_{\sqrt{\delta_{\acom}}} \hat{M}^{(\Point,\zpt),z,s}_{b} \cdot \hat{M}^{(\Point,\xpt),x,r}_{a}\;,
	\end{align}
	which follows from the approximate commutation relation of Lemma~\ref{lem:qld-comm-cons} and the projectivity of $\{ \hat{M}^{(\Point,\zpt),z,s}_{b} \}$, together with Fact~\ref{fact:add-a-proj}. Next we show that $\{ \hat{R}^\omega_{a,b} \}$ is approximately self-consistent:  we aim to show that on average over $\omega$, it holds that $(\hat{R}^\omega_{a,b})_{\alice \alice'} \simeq_\delta (\hat{R}^\omega_{a,b})_{\bob \alice''}$ for some polynomial function $\delta = \poly(\eps)$. To show this, we perform the following calculation: (recall the state $\ket{\hat{\psi}}$ defined in~\eqref{eq:def-psihat})
	\begin{align}
		&\E_\omega \sum_{a,b} \bra{\hat{\psi}} (\hat{R}^\omega_{a,b})_{\alice \alice'} \otimes (\hat{R}^\omega_{a,b})_{\bob \alice''} \ket{\hat{\psi}} \notag \\
		&\approx_{\delta_{\acom}^{1/4}} \E_\omega \sum_{a,b} \bra{\hat{\psi}} (\hat{M}^{(\Point,\zpt),z,s}_{b} \cdot \hat{M}^{(\Point,\xpt),x,r}_{a})_{\alice \alice'} \otimes (\hat{R}^\omega_{a,b})_{\bob \alice''} \ket{\hat{\psi}} \label{eq:qld-rw-self-cons-1} \\
		&\approx_{\delta_{\acom}^{1/4}} \E_\omega \sum_{a,b} \bra{\hat{\psi}} (\hat{M}^{(\Point,\zpt),z,s}_{b} \cdot \hat{M}^{(\Point,\xpt),x,r}_{a})_{\alice \alice'} \otimes (\hat{M}^{(\Point,\zpt),z,s}_{b} \cdot \hat{M}^{(\Point,\xpt),x,r}_{a})_{\bob \alice''} \ket{\hat{\psi}} \label{eq:qld-rw-self-cons-2} \\		
		&\approx_{\sqrt{\eps}} \E_\omega \sum_{a,b} \bra{\hat{\psi}} (\hat{M}^{(\Point,\zpt),z,s}_{b})_{\alice \alice'} \otimes (\hat{M}^{(\Point,\zpt),z,s}_{b} \cdot \hat{M}^{(\Point,\xpt),x,r}_{a})_{\bob \alice''} \ket{\hat{\psi}} \label{eq:qld-rw-self-cons-3} \\
		&= \E_\omega \sum_{b} \bra{\hat{\psi}} (\hat{M}^{(\Point,\zpt),z,s}_{b})_{\alice \alice'} \otimes (\hat{M}^{(\Point,\zpt),z,s}_{b})_{\bob \alice''} \ket{\hat{\psi}} \notag \\
		&\approx_{\eps} \E_\omega \sum_{b} \bra{\hat{\psi}} (\hat{M}^{(\Point,\zpt),z,s}_{b})_{\alice \alice'} \ket{\hat{\psi}} \label{eq:qld-rw-self-cons-4} \\
		&= 1\;. \notag
	\end{align}
   	The approximation in~\eqref{eq:qld-rw-self-cons-1} is derived by bounding the magnitude of the difference:
	\begin{align*}
		&\left | \E_\omega \sum_{a,b} \bra{\hat{\psi}} (\hat{R}^\omega_{a,b} - \hat{M}^{(\Point,\zpt),z,s}_{b} \cdot \hat{M}^{(\Point,\xpt),x,r}_{a})_{\alice \alice'} \otimes (\hat{R}^\omega_{a,b})_{\bob \alice''} \ket{\hat{\psi}} \right | \\
		&\leq \sqrt{\E_\omega \sum_{a,b} \bra{\hat{\psi}} (\hat{R}^\omega_{a,b} - \hat{M}^{(\Point,\zpt),z,s}_{b} \cdot \hat{M}^{(\Point,\xpt),x,r}_{a})_{\alice \alice'}^2 \ket{\hat{\psi}} } \cdot \sqrt{\E_\omega \sum_{a,b} \bra{\psi} (\hat{R}^\omega_{a,b})^2_{\bob \alice''} \ket{\psi} } \\
		&\leq \sqrt{\delta_{\acom}^{1/2}} \cdot \sqrt{1}\;,
	\end{align*}
	where the second line follows from Cauchy-Schwarz, and the third line follows from~\eqref{eq:qld-r-2} and the fact that $\sum_{a,b} (\hat{R}^\omega_{a,b})^2 \leq \Id$. The approximation in~\eqref{eq:qld-rw-self-cons-2} follows from a similar calculation. 
	The approximation in~\eqref{eq:qld-rw-self-cons-3} is derived by bounding the magnitude of the difference:
	\begin{align}
		&\left | \E_\omega \sum_{a,b} \bra{\hat{\psi}} (\hat{M}^{(\Point,\zpt),z,s}_{b} \cdot (\Id - \hat{M}^{(\Point,\xpt),x,r}_{a}))_{\alice \alice'}  \otimes (\hat{M}^{(\Point,\zpt),z,s}_{b} \cdot \hat{M}^{(\Point,\xpt),x,r}_{a})_{\bob \alice''} \ket{\hat{\psi}} \right | \notag \\
		&= \left | \E_\omega  \bra{\hat{\psi}} \Big(\sum_b \hat{M}^{(\Point,\zpt),z,s}_{b} \otimes \hat{M}^{(\Point,\zpt),z,s}_{b}\Big) \cdot \Big(\sum_{a} (\Id - \hat{M}^{(\Point,\xpt),x,r}_{a}) \otimes \hat{M}^{(\Point,\xpt),x,r}_{a} \Big) \ket{\hat{\psi}} \right | \\
		&\leq \sqrt{\E_\omega \sum_{a} \bra{\hat{\psi}} (\Id - \hat{M}^{(\Point,\xpt),x,r}_{a})_{\alice \alice'}^2 \otimes (\hat{M}^{(\Point,\xpt),x,r}_{a})_{\bob \alice''}^2 \ket{\hat{\psi}}} \label{eq:qld-rw-self-cons-3-1}\\
                &= \sqrt{\E_\omega \sum_{a} \bra{\hat{\psi}} (\Id - \hat{M}^{(\Point,\xpt),x,r}_{a})_{\alice \alice'} \otimes (\hat{M}^{(\Point,\xpt),x,r}_{a})_{\bob \alice''} \ket{\hat{\psi}}} \label{eq:qld-rw-self-cons-3-2}\\
                &= \sqrt{ 1 - \E_{\omega} \sum_a \bra{\hat{\psi}} (\hat{M}^{(\Point,\xpt),x,r}_{a})_{\alice \alice'} \ot  (\hat{M}^{(\Point,\xpt),x,r}_{a})_{\bob \alice''} \ket{\hat{\psi}}} \label{eq:qld-rw-self-cons-3-3}\\
                &=\sqrt{ \frac{1}{2} \E_{\omega} \sum_a \bra{\hat{\psi}} ( (\hat{M}^{(\Point,\xpt),x,r}_{a})_{\alice \alice'} \ot I_{\bob \alice''} - I_{\alice \alice'} \ot  (\hat{M}^{(\Point,\xpt),x,r}_{a})_{\bob \alice''})^2 \ket{\hat{\psi}}} \label{eq:qld-rw-self-cons-3-4}\\
		&\leq \sqrt{\eps}\;. \label{eq:qld-rw-self-cons-3-5}
	\end{align}
	Here \Cref{eq:qld-rw-self-cons-3-1} follows from the Cauchy-Schwarz inequality, the projectivity of $\hat{M}^{(\Point,\xpt),x,r}_{a}$, and the fact that $\sum_b \hat{M}^{(\Point,\zpt),z,s}_{b} \otimes \hat{M}^{(\Point,\zpt),z,s}_{b} \leq \Id$. \Cref{eq:qld-rw-self-cons-3-2,eq:qld-rw-self-cons-3-3,eq:qld-rw-self-cons-3-4} follow from projectivity and completeness of $\hat{M}^{(\Point, \xpt), x, r}_{a}$. \Cref{eq:qld-rw-self-cons-3-5} follows from the self-consistency of $\hat{M}^{(\Point,\xpt),x,r}_{a}$, established in \Cref{lem:qld-comm-cons}. 

        The approximation in~\eqref{eq:qld-rw-self-cons-4} follows from the self-consistency of $\hat{M}^{(\Point, \xpt), x, r}_a$, again from \Cref{lem:qld-comm-cons}, together with projectivity and \Cref{fact:agreement}.

        
	
	This shows that $(\hat{R}^\omega_{a,b})_{\alice \alice'} \simeq_\delta (\hat{R}^\omega_{a,b})_{\bob \alice''}$, as desired. Combined with \Cref{fact:agreement}, we have
	\begin{equation}
		\label{eq:rw-sc}
		(\hat{R}^\omega_{a,b})_{\alice \alice'} \approx_\delta (\hat{R}^\omega_{a,b})_{\bob \alice''}\;.
	\end{equation}


		Finally we apply the orthonormalization Lemma (Lemma~\ref{lem:ortho}) to $\{ \hat{R}^\omega_{a,b} \}$ in order to obtain projective measurements $\{ \hat{L}^\omega_{a,b} \}$ such that on average over $\omega$,
	\begin{equation}\label{eq:rl-close}
		\hat{L}^\omega_{a,b} \approx_{\delta} \hat{R}^\omega_{a,b}\;.
	\end{equation}
	
	\textbf{Second step: approximate linearity relations for $\{ \hat{R}^\omega_{a,b} \}$.}
	We proceed in three steps. First we establish linearity relations for $\{ \hat{M}^{(\Point,W),u,r}_b \}$. For all $W \in \{X,Z\}$, $u \in \F_q^m$, $r, r' \in \F_q$ and $c \in \F_2$, 
	\begin{align}
		\sum_{\substack{b,b' \in \F_2 : \\ b + b' = c }} \hat{M}^{(\Point,W),u,r}_b \, \hat{M}^{(\Point,W),u,r'}_{b'} &= \sum_{\substack{a,a' \in \F_q: \\ \tr(ar + a'r') = c}} \hat{M}^{(\Point,W),u}_a \cdot \hat{M}^{(\Point,W),u}_{a'} \notag\\
		&= \sum_{a : \tr(a (r +r')) = c } \hat{M}^{(\Point,W),u}_a \notag\\
		&= \hat{M}^{(\Point,W),u,r+r'}_c\;,\label{eq:m-lin}
	\end{align}
	where the second line follows from projectivity of $\{\hat{M}^{(\Point,W),u}_a\}$. 
	
	Second we show approximate linearity of the $\{ \hat{R}^\omega_{a,b} \}$. On average over $x,z \in \F_q^m$ and $r,r',s,s' \in \F_q$, 
	\begin{align}
	& \sum_{\substack{a,b,a',b' : \\ a + a' = c \\ b + b' = d}} \Big (\hat{R}^{x,z,r,s}_{a,b} \cdot \hat{R}^{x,z,r',s'}_{a',b'} \Big)_{\alice \alice'} \notag\\
	&\approx_\delta \sum_{\substack{a,b,a',b' : \\ a + a' = c \\ b + b' = d}} \Big( \hat{R}^{x,z,r,s}_{a,b} \Big)_{\alice \alice'} \otimes \Big( \hat{R}^{x,z,r',s'}_{a',b'} \Big)_{\bob \alice''} \label{eq:qld-r-52} \\
		&\approx_\delta \sum_{\substack{a,b,a',b': \\ a + a' = c \\ b + b' = d}} \Big( \hat{M}^{(\Point,\zpt),z,s}_{b} \cdot \hat{M}^{(\Point,\xpt),x,r}_{a}
 \Big)_{\alice \alice'} \otimes \Big(\hat{M}^{(\Point,\zpt),z,s'}_{b'} \cdot \hat{M}^{(\Point,\xpt),x,r'}_{a'} \Big)_{\bob \alice''} \label{eq:qld-r-53} \\
 	&\approx_\delta \sum_{\substack{b,b': \\ b + b' = d}} \Big( \hat{M}^{(\Point,\zpt),z,s}_{b}
 \Big)_{\alice \alice'} \otimes \Big(\hat{M}^{(\Point,\zpt),z,s'}_{b'} \cdot \hat{M}^{(\Point,\xpt),x,r+r'}_{c} \Big)_{\bob \alice''} \label{eq:qld-r-54} \\
 &\approx_\delta \sum_{\substack{b,b' : \\ b + b' = d}} \Big( \hat{M}^{(\Point,\zpt),z,s}_{b}
 \Big)_{\alice \alice'} \otimes \Big(\hat{M}^{(\Point,\xpt),x,r+r'}_{c} \cdot \hat{M}^{(\Point,\zpt),z,s'}_{b'}  \Big)_{\bob \alice''} \;.\label{eq:qld-r-55}
 \end{align}
	In each approximation, the answer summation is over $c,d \in \F_2$. Each approximation follows by applying a previously derived approximation to part of the expression, together with \Cref{fact:add-a-proj}.\anote{Eventual TODO: split up all these steps} Specifically, line~\eqref{eq:qld-r-52} uses the self-consistency of $\{ \hat{R}^\omega_{a,b} \}$ shown in~\eqref{eq:rw-sc}. Line~\eqref{eq:qld-r-53} follows from using~\eqref{eq:qld-r-2} twice. Line~\eqref{eq:qld-r-54} follows from the self-consistency and linearity properties of $\{ \hat{M}^{(\Point,X),x,r}_a \}$ shown in Lemma~\ref{lem:qld-comm-cons} and~\eqref{eq:m-lin} respectively. Line~\eqref{eq:qld-r-55} follows from the approximate commutativity relation shown in Lemma~\ref{lem:qld-comm-cons}.  Continuing,
 \begin{align}
\text{\eqref{eq:qld-r-55}} &\approx_\delta \Big(\hat{M}^{(\Point,\xpt),x,r+r'}_{c} \cdot \hat{M}^{(\Point,\zpt),z,s+s'}_{d}  \Big)_{\bob \alice''} \label{eq:qld-r-56} \\
 &\approx_\delta \Big ( \hat{R}^{x,z,r+r',s+s'}_{c,d} \Big)_{\bob \alice''} \label{eq:qld-r-57} \\
 &\approx_\delta \Big ( \hat{R}^{x,z,r+r',s+s'}_{c,d} \Big)_{\alice \alice'} \label{eq:qld-r-58}.
	\end{align}	
Line~\eqref{eq:qld-r-56} follows from self-consistency and linearity properties of $\{ \hat{M}^{(\Point,Z),z,s}_b \}$. Line~\eqref{eq:qld-r-57} follows from~\eqref{eq:qld-r-2} and \Cref{lem:qld-comm-cons}, and line~\eqref{eq:qld-r-58} follows from the self-consistency of $\{ \hat{R}^\omega_{a,b} \}$ from Eq.~\eqref{eq:rw-sc}. 
	
Third we deduce approximate linearity of the $\{ \hat{L}^\omega_{a,b} \}$. On average over $x,z \in \F_q^m$, and $r,r',s,s' \in \F_q$, we have
	\begin{align}
	\sum_{\substack{a,b,a',b': \\ a + a' = c \\ b + b' = d}}  \Big ( \hat{L}^{x,z,r,s}_{a,b} \, \hat{L}^{x,z,r',s'}_{a',b'} \Big)_{\alice \alice'} &\approx_\delta  \sum_{\substack{a,b,a',b': \\ a + a' = c \\ b + b' = d}} \Big( \hat{L}^{x,z,r,s}_{a,b} \Big)_{\alice \alice'} \otimes \Big (\hat{L}^{x,z,r',s'}_{a',b'} \Big)_{\bob \alice''}\notag \\
		&\approx_\delta  \sum_{\substack{a,b,a',b' : \\ a + a' = c \\ b + b' = d}}  \Big( \hat{R}^{x,z,r,s}_{a,b} \Big)_{\alice \alice'} \otimes \Big( \hat{R}^{x,z,r',s'}_{a',b'} \Big)_{\bob \alice''} \notag\\
		&\approx_\delta \Big ( \hat{R}^{x,z,r+r',s+s'}_{c,d} \Big)_{\alice \alice'} \notag\\
		&\approx_\delta \Big ( \hat{L}^{x,z,r+r',s+s'}_{c,d} \Big)_{\alice \alice'}\;.\label{eq:l-al}
	\end{align}
	The first two approximations follow from the closeness of the $\{ \hat{R}^\omega_{a,b} \}$ and $\{ \hat{L}^\omega_{a,b} \}$  guaranteed by~\eqref{eq:rl-close} as well as the self-consistency of $\{ \hat{R}^\omega_{a,b} \}$ shown in~\eqref{eq:rw-sc}. The third approximation follows from the linearity properties of $\{ \hat{R}^\omega_{a,b} \}$ in~\eqref{eq:qld-r-55}. 
	
	\textbf{Third step: construction of $\{ \hat{Q}^{x,z}_{a,b} \}$.}
	For each $\omega$ define an observable $\mathcal{O}^\omega = \sum_{a,b \in \F_2} (-1)^{a+ b} \hat{L}^\omega_{a,b}$.  The self-consistency of $\mathcal{O}^\omega$ between registers $\alice \alice'$ and $\bob \alice''$ follows from the self-consistency of the $\{ \hat{L}^\omega_{a,b} \}$, which is obtained by combining~\eqref{eq:rw-sc} and~\eqref{eq:rl-close}. We now verify approximate linearity. On average over $x,z \in \F_q^m$, and $r,r',s,s' \in \F_q$, 
	\begin{align*}
		\mathcal{O}^{x,z,r,s} \, \mathcal{O}^{x,z,r',s'} &= \sum_{a,b,a',b'} (-1)^{a + b + a' + b'} \hat{L}^{x,z,r,s}_{a,b} \, \hat{L}^{x,z,r',s'}_{a',b'} \\
		&\approx_\delta \sum_{c,c'} (-1)^{c+c'} \hat{L}^{x,z,r+r',s+s'}_{c,c'} \\
		&= \mathcal{O}^{x,z,r+r',s+s'}\,
	\end{align*}
	where the approximation follows from the approximate linearity properties of $\{ \hat{L}^\omega_{a,b} \}$ shown in~\eqref{eq:l-al}.
	Identify the field $\F_q$ with the vector space $\F_2^t$, and thus we treat pairs $(r,s)$ as vectors in $\F_2^{2t}$. For every $x,z \in \F_q^m$, we can apply Theorem~\ref{thm:linearity} to the collection of observables $\{ \mathcal{O}^{x,z,r,s} \}_{r,s \in \F_q}$ to obtain the existence of ``{exactly linear}'' observables $\{ \mathcal{L}^{x,z,r,s} \}$ that are $\delta$-close to $\{ \mathcal{O}^{x,z,r,s} \}$, on average over the choice of $x,z,r,s$. By assuming that $\ket{\hat{\psi}}$ has sufficiently many ancilla zero qubits, we do not need to extend the state further. 
	
	We are ready to define the measurements $\{\hat{Q}^{x,z}_{a,b} \}$ whose existence is claimed in the statement of the lemma. For all $x,z\in \F_q^m$ and $a,b \in \F_q$, define
	\[
		\hat{Q}^{x,z}_{a,b} = \E_{r,s} (-1)^{\tr(ra + sb)} \mathcal{L}^{x,z,r,s}\;.
	\]
	This operator is a projection (and thus positive semidefinite):
	\begin{align*}
	(\hat{Q}^{x,z}_{a,b})^2 &= \E_{r,s,r',s'} (-1)^{\tr((r+r')a + (s+s')b)} \mathcal{L}^{x,z,r,s} \, \mathcal{L}^{x,z,r',s'}\\ 
	&= \E_{r,s,r',s'} (-1)^{\tr((r+r')a + (s+s')b)} \mathcal{L}^{x,z,r+r',s+s'} \\
	&= \E_{r,s} (-1)^{\tr(ra + sb)}  \mathcal{L}^{x,z,r,s}\;.
	\end{align*}
	The second equality follows from the exact linearity of the $\{ \mathcal{L}^{x,z,r,s} \}$. Next, the operators sum to identity:
	\begin{align*}
		\sum_{a,b} \hat{Q}^{x,z}_{a,b} &= \E_{r,s} \sum_{a,b} (-1)^{\tr(ra + sb)} \mathcal{L}^{x,z,r,s} \\
		&= \mathcal{L}^{x,z,0,0} \\
		&= \Id\;.
	\end{align*}
	Thus $\{\hat{Q}^{x,z}_{a,b}\}$ forms a projective measurement.
	
	We establish self-consistency:
	\begin{align*}
		\Big( \hat{Q}^{x,z}_{a,b} \Big)_{\alice \alice'} &= \Big( \E_{r,s} (-1)^{\tr(ra + sb)} \mathcal{L}^{x,z,r,s} \Big)_{\alice \alice'} \\
		&\approx_\delta \Big( \E_{r,s} (-1)^{\tr(ra + sb)} \mathcal{O}^{x,z,r,s} \Big)_{\alice \alice'} \\
		&\approx_\delta \Big( \E_{r,s} (-1)^{\tr(ra + sb)} \mathcal{O}^{x,z,r,s} \Big)_{\bob \alice''} \\
		&\approx_\delta \Big( \hat{Q}^{x,z}_{a,b} \Big)_{\bob \alice''}\,.
	\end{align*}
	The first approximation is due to~\Cref{lem:avg-closeness}. The second approximation is due to~\Cref{lem:avg-closeness} and the self-consistency of the $\{\mathcal{O}^{x,z,r,s}\}$ observables. The third approximation follows from~\Cref{lem:avg-closeness} again. 
	
	Finally we show consistency with the $\{ \mathcal{M}^{(\Point,W),u}_a\}$. 
	\begin{align*}
		\Big( \hat{Q}^{x,z}_{a,b} \Big)_{\alice \alice'}&\approx_\delta \Big( \E_{r,s} (-1)^{\tr(ra + sb)} \mathcal{O}^{x,z,r,s} \Big)_{\bob \alice''} \\
		&= \Big( \E_{r,s} \sum_{c,d} (-1)^{\tr(ra + sb) + c + d} \, \hat{L}^{x,z,r,s}_{c,d}  \Big)_{\bob \alice''} \\
		&\approx_\delta \Big( \E_{r,s} \sum_{c,d} (-1)^{\tr(ra + sb) + c + d} \, \hat{R}^{x,z,r,s}_{c,d}  \Big)_{\bob \alice''} \\
		&\approx_\delta \Big( \E_{r,s} \sum_{c,d} (-1)^{\tr(ra + sb) + c + d} \, \hat{M}^{(\Point,\zpt),z,s}_d \, \hat{M}^{(\Point,\xpt),x,r}_c \Big)_{\bob \alice''} \\
                                                                &= \Big ( \E_s \sum_d (-1)^{\tr(sb) + d} \hat{M}^{(\Point,\zpt),z,s}_d \Big) \cdot \Big (\E_r \sum_c (-1)^{\tr(ra) + c} \hat{M}^{(\Point,\xpt),x,r}_c  \Big)_{\bob \alice''} \\
                                                                &= \Big( \E_s \sum_d \sum_{b', a'}  (-1)^{\tr(sb) + \tr(sb')} \hat{M}^{(\Point, \zpt), z}_{b'} \Big)_{\bob \alice''} \cdot \Big( \E_r \sum_c (-1)^{\tr(ra) + \tr(r a')} \hat{M}^{(\Point, \xpt), x}_{b'} \Big)_{\bob \alice''} \\
		&= \Big ( \hat{M}^{(\Point,\zpt),z}_b \cdot \hat{M}^{(\Point,\xpt),x}_a \Big)_{\bob \alice''}\;.
	\end{align*}
	The second line follows from the definition of $\{ \mathcal{O}^{x,z,r,s} \}$. The third line follows from the closeness of $\{\hat{L}^\omega_{c,d}\}$ and $\{\hat{R}^\omega_{c,d} \}$. The fourth line follows from~\eqref{eq:qld-r-2}.  The fifth line follows from~\eqref{eq:qld-def-mptur}, and the last line from \Cref{lem:cancellation}.

        The above calculation establishes \Cref{eq:qld-q-cons-m-hat-zx} of the Lemma. To obtain \Cref{eq:qld-q-cons-m-hat-xz}, which is the same  consistency relation with $X$ and $Z$ swapped, we may use \Cref{lem:qld-comm-cons} immediately after the fourth line to swap the two operators in the product, and then proceed as above.
	
	An analogous derivation shows the same properties of $\{ \hat{Q}^{x,z}_{a,b} \}$ with the tensor factors switched to $\bob \bob'$ and $\alice \bob''$, respectively. This concludes the proof of the lemma.
	


\end{proof}


\begin{lemma}\label{lem:qld-xz-lines}
  There exists a function $\delta_P(\eps,m,d,q) = \poly(\eps,
  md/q)$ such that the following holds. For $\mH \in \{\mH_\alice,\mH_\bob\}$ and for every pair of lines $\line_X$ and $\line_Z$ there exists a POVM $\{T^{\line_X, \line_Z}_{f_X, f_Z}\}$ acting on
  $\mathcal{H} \otimes (\C^q)^{\otimes \nqubits}$ with answers
  consisting of pairs of polynomials $f_X, f_Z$ in $\deg_{md}(\line_X)
  \times \deg_{md}(\line_Z)$
  such that on average over pairs $(\line_X, x)$ and $(\line_Z, z)$
  independently sampled from the line-point distribution,
  \begin{align*}
    (T^{\line_X, \line_Z}_{[ \eval_x(\cdot) =a, \eval_z(\cdot) = b]})_{\alice \alice'} &\simeq_{\delta_P} (\hat{Q}^{x,z}_{a, b})_{\bob \alice''}, \\
\end{align*}
  as well as all symmetric equivalents of this approximation. Moreover, if $\line_W$ is axis-parallel, then $f_W \in \deg_{d}(\line_W)$.
\end{lemma}

\begin{proof}
  We construct the required measurements by combining the measurements
  $\hat{M}^{(\Line,W),\line_W}_{f_W}$ for $W \in \{X, Z\}$. For convenience in the proof we slightly abuse the notation $\approx_\delta$ to mean $\approx_{\poly(\eps)}$ for some polynomial function $\delta=\poly(\eps)$ that may differ each time the notation is used.

  	We first argue that the measurements $\{ \hat{Q}^{x,z}_{a,b} \}$ obtained in Lemma~\ref{lem:qld-4-10} are consistent with the $\{\hat{M}^{(\Line,W),\line_W}_{f_W} \}$. 

	Applying Fact~\ref{fact:add-a-proj} to~\eqref{eq:qld-q-cons-m-hat-zx}, we get that on average over uniformly random $x,z$,
	\begin{align}
		\Big( \hat{Q}^{x,z}_{a,b} \Big)_{\alice \alice'} \cdot \Big( \hat{M}^{(\Point,\zpt),z}_{b} \Big)_{\bob \alice''}
&=	\Big( \hat{M}^{(\Point,\zpt),z}_{b} \Big)_{\bob \alice''}\cdot	\Big( \hat{Q}^{x,z}_{a,b} \Big)_{\alice \alice'} \notag\\
& \approx_\delta \Big (\hat{M}^{(\Point,\zpt),z}_{b} \hat{M}^{(\Point,\xpt),x}_{a} \Big)_{\bob \alice''} \notag\\
&\approx_\delta \Big( \hat{Q}^{x,z}_{a,b} \Big)_{\alice \alice'}\;.\notag
	\end{align}
	Since $\{\hat{Q}^{x,z}_{a,b}\}$ is a projective measurement, applying \Cref{lem:cool-closeness-fact} with $(\hat{Q}^{x,z}_{a,b})_{\alice, \alice'}$ playing the role of $A^{x}_a$ and $(\hat{Q}^{x,z}_{a,b} )_{\alice, \alice'} \cdot (\hat{M}^{(\Point, \zpt), z}_{b})_{\bob \alice''}$ playing the role of $B^{x}_a$, we obtain that
	\begin{equation}
	\Big( \hat{Q}^{x,z}_{a} \Big)_{\alice \alice'} \approx_\delta \sum_b \Big( \hat{Q}^{x,z}_{a,b} \Big)_{\alice \alice'} \cdot \Big( \hat{M}^{(\Point,\zpt),z}_{b} \Big)_{\bob \alice''}\;.\label{eq:qqm}
	\end{equation}
	We now show that on average over uniformly random $x,z$, 
	\begin{equation}
	\label{eq:qld-qxz-close-to-point}
	\Big ( \hat{Q}^{x,z}_{a} \Big)_{\alice \alice'} \approx_{\delta_Q} \Big (\hat{M}^{(\Point,X),x}_{a} \Big)_{\bob \alice''} \;.
	\end{equation}
	This follows from~\eqref{eq:qqm} and the following calculation:
	\begin{align*}
		&\E_{x,z} \sum_a \norm{ \sum_b  \Big( \hat{Q}^{x,z}_{a,b} \Big)_{\alice \alice'} \cdot \Big ( \hat{M}^{(\Point,\zpt),z}_{b} \Big)_{\bob \alice''} - \Big ( \hat{M}^{(\Point,\xpt),x}_{a} \Big)_{\bob \alice''}  \ket{\hat{\psi}} }^2  \\
		&=\E_{x,z} \sum_a \norm{ \sum_b \Big ( \hat{M}^{(\Point,\zpt),z}_{b} \Big)_{\bob \alice''} \cdot  \Big( \Big( \hat{Q}^{x,z}_{a,b} \Big)_{\alice \alice'} - \Big ( \hat{M}^{(\Point,\xpt),x}_{a} \Big)_{\bob \alice''}  \Big) \ket{\hat{\psi}} }^2  \\
		&= \E_{x,z} \sum_{a,b}  \norm{  \Big ( \hat{M}^{(\Point,\zpt),z}_{b} \Big)_{\bob \alice''} \cdot  \Big( \Big( \hat{Q}^{x,z}_{a,b} \Big)_{\alice \alice'} - \Big ( \hat{M}^{(\Point,\xpt),x}_{a} \Big)_{\bob \alice''}  \Big) \ket{\hat{\psi}} }^2  \\
                &= \E_{x,z} \sum_{a,b} \norm{ \Big(  \hat{M}^{(\Point,\zpt),z}_{b} \Big)_{\bob \alice''} \cdot  \Big( \Big( \hat{Q}^{x,z}_{a,b} \Big)_{\alice \alice'} - \Big( \hat{M}^{(\Point, \zpt),z}_{b} \cdot \hat{M}^{(\Point,\xpt),x}_{a} \Big)_{\bob \alice''}  \Big) \ket{\hat{\psi}} }^2  \\
                &\leq \E_{x,z} \sum_{a,b} \norm{ \Big( \Big( \hat{Q}^{x,z}_{a,b} \Big)_{\alice \alice'} - \Big( \hat{M}^{(\Point, \zpt),z}_{b} \cdot \hat{M}^{(\Point,\xpt),x}_{a} \Big)_{\bob \alice''}  \Big) \ket{\hat{\psi}} }^2  \\
		&= O(\delta)\;,
	\end{align*} 
	where the second line uses the property $\sum_b \hat{M}^{(\Point,\zpt),z}_b = I$, the third and fourth lines uses the projectivity of $\{ \hat{M}^{(\Point,\zpt),z}_{b} \}$, the fifth line uses $\hat{M}^{(\Point, \zpt), z}_{b} \leq I$, and the last line uses \Cref{eq:qld-q-cons-m-hat-zx}. This shows~\eqref{eq:qld-qxz-close-to-point}. A similar argument shows that
	\begin{equation}\label{eq:qld-qxz-close-to-point-2}
	\Big ( \hat{Q}^{x,z}_{b} \Big)_{\alice \alice'} \approx_{\delta_Q} \Big (\hat{M}^{(\Point,Z),z}_{b} \Big)_{\bob \alice''}\;.
	\end{equation}
	Next, from Lemma~\ref{lem:qld-comm-line-cons} we get that on average over $(\line_X,x)$ and $(\line_Z,z)$ sampled independently from the line-point distribution (\Cref{def:line-point-dist}),
	\begin{gather}
	(\hat{M}^{(\Line,X),\line_\xpt}_{[\eval_x(\cdot) = a]})_{\bob \alice''} \approx_\delta  (\hat{M}^{(\Point,\xpt),x}_{a})_{\bob \alice''} \;,\\
	( \hat{M}^{(\Line,Z),\line_\zpt}_{[\eval_z(\cdot) = b]})_{\bob \alice''} \approx_\delta  (\hat{M}^{(\Point,\zpt),z}_{b})_{\bob \alice''}\;.
	\end{gather}
	where we used both the self-consistency of the line measurements (item 1 in Lemma~\ref{lem:qld-comm-line-cons}) as well as their consistency with the points measurements (item 3). To deduce both of these relations, we used \Cref{fact:agreement} to convert between $\approx$ and $\simeq$ distance, and \Cref{fact:data-processing} to sum over outcomes to the lines measurement. Since the $\{\hat{M}^{(\Line,W),\line}_f \}$ and $\{ \hat{Q}^{x,z}_{a,b} \}$ are projective, together with~\eqref{eq:qld-qxz-close-to-point} and~\eqref{eq:qld-qxz-close-to-point-2}
	by Fact~\ref{fact:agreement} we get
	\begin{align}\label{eq:pasting-q1}
		\Big ( \hat{Q}^{x,z}_{a} \Big)_{\alice \alice'} \simeq_{\delta_Q} \Big (\hat{M}^{(\Line,\xpt),\line_\xpt}_{[\eval_x(\cdot) = a]} \Big)_{\bob \alice''} \qquad \text{and} \qquad \Big ( \hat{Q}^{x,z}_{b} \Big)_{\alice \alice'} \simeq_{\delta_Q} \Big (\hat{M}^{(\Line,\zpt),\line_\zpt}_{[\eval_z(\cdot) = b]} \Big )_{\bob \alice''} \;.
	\end{align}

	For all lines $\line_\xpt,\line_\zpt$ and $(f_\xpt,f_\zpt) \in \deg_{md}(\line_\xpt) \times \deg_{md}(\line_\zpt)$ define the operator
	\[
		T^{\line_\xpt,\line_\zpt}_{f_\xpt,f_\zpt} = \hat{M}^{(\Line,\xpt),\line_\xpt}_{f_\xpt} \cdot \hat{M}^{(\Line,\zpt),\line_\zpt}_{f_\zpt}  \cdot \hat{M}^{(\Line,\xpt),\line_\xpt}_{f_\xpt} \;.
	\]
	Then the collection $\{ T^{\line_\xpt,\line_\zpt}_{f_\xpt,f_\zpt} \}$ forms a valid POVM. Moreover, by the fact that $\{\hat{M}^{(\Line, W), \line_W}_{f_W}\}$ comes from a valid strategy for the low-degree test, it holds that $f_W \in \deg_{d}(\line_W)$ whenver $\line_W$ is axis-parallel, and thus property holds for $T$ as well. 

        It only remains to bound the consistency of $\{T^{\line_\xpt, \line_\zpt}_{f_\xpt, f_\zpt}\}$ with $\{\hat{Q}^{x,z}_{a,b}\}$. Using the following choices of measurements,
	\begin{equation*}
		``A^{x,y_1,y_2}_{a_1,a_2}" : \hat{Q}^{x,z}_{a,b} \;, \quad
		``(G_1)^x_g" : \hat{M}^{(\Line,\xpt),\line_\xpt}_{f_\xpt} \;, \quad
		``(G_2)^x_g" : \hat{M}^{(\Line,\zpt),\line_\zpt}_{f_\zpt} \;, \quad
		``J^x_{g_1,g_2}": T^{\line_\xpt,\line_\zpt}_{f_\xpt,f_\zpt} \;,
              \end{equation*}
	Then the hypotheses of \Cref{lem:pasting} hold, with $\delta$ in the lemma set to $\delta_Q$, and $\eta$ in the lemma set to $md/q$ (this follows from the Schwartz-Zippel lemma applied to the polynomials $f_X, f_Z$, which are of total degree at most $md$). Applying the lemma, we conclude that
	\begin{equation}
	\label{eq:qld-4-13-1}
		\Big( \hat{Q}^{x,z}_{a,b} \Big)_{\alice \alice'} \simeq_{\delta_P} \Big( T^{\line_\xpt,\line_\zpt}_{[\eval_x(\cdot) = a, \eval_z(\cdot) = b]} \Big)_{\bob \alice''}\;,
	\end{equation}
	where $\delta_P = \poly(\eps,md/q)$ is the $\delta_{pasting}$ as given by Lemma~\ref{lem:pasting}, conditions~\eqref{eq:pasting-1} and~\eqref{eq:pasting-2} in the lemma are given by~\eqref{eq:pasting-q1} and~\eqref{eq:qld-q-self-cons} respectively, 
	and the distribution over $(\ell_\xpt,\ell_\zpt)$ are two independently and uniformly chosen lines. This is a symmetric equivalent of the consistency relation in the conclusion of Lemma. Thus, by applying an identical argument with the appropriate modifications to the registers, all symmetric equivalents of this relation hold, establishing the conclusion of the Lemma.



\end{proof}


      \subsection{Applying the classical low-degree test}
      \label{sec:apply-ldt}
      In the previous section we showed how to combine the
      approximately commuting $X$ and $Z$ basis point and line
      measurements into a joint point measurement and a joint line
      measurement. In this section we show how these measurements can be combined
			in a single measurement
      that returns a pair of global polynomials encoding the $X$-basis and
      $Z$-basis information, respectively, by applying the lemma that states quantum soundness of the classical low-degree test, Theorem~\ref{lem:ld-soundness}.
			
      Introduce a ``combining'' map that
      takes two polynomials $f, g \in \ideg_{d,m}(\F_q)$ and returns a new polynomial $\combine_{f,g} \in \ideg_{d,
        2m+2}(\F_q)$ defined by
      \[ \combine_{f,g}(\underbrace{x}_{\in \F_q^{m}},
        \underbrace{z}_{\in \F_q^m}, \underbrace{\alpha}_{\in \F_q},
        \underbrace{\beta}_{\in \F_q}) = \alpha f(x) + \beta g(z)\;. \]
      This combining map can also be defined for polynomials restricted to lines. Let $\line$ be a line in $\F_q^{2m+2}$ with intercept $u = (u_X, u_Z, u_\alpha, u_\beta)$ and slope $v = (v_X, v_Z, v_\alpha, v_\beta)$. Let $\line_X$ and $\line_Z$ be lines in $\F_q^{m}$ such that for any point $(x,z,\alpha,\beta) \in \line$, it holds that $x \in \line_X$ and $z \in \line_Z$. Then for any two polynomials $f \in \deg_{dm}(\line_X)$ and $g \in \deg_{dm}(\line_Z)$, we define $\combine_{f,g} \in \deg_{dm + 1}(\line)$ by
      \begin{equation}
        \combine_{f,g}(u + t \cdot v) = (u_\alpha + t\cdot v_\alpha) f(u_X + t \cdot v_X) + (u_\beta + t \cdot v_\beta) g(u_Z + t \cdot v_Z)\;. \label{eq:combine-lines}
      \end{equation}
      

      In the following two lemma we define combined point and line
      measurements according to this definition of $\combine_{f,g}$.

\begin{lemma}\label{lem:qld-4-12}
	For all $x,z \in \F_q^m$ and $\alpha,\beta \in \F_q$ and $\mH \in \{ \mH_\alice, \mH_\bob\}$ there exists a projective measurement $\{ \hat{Q}^{x,z,\alpha,\beta}_c \}_{c \in \F_q}$ acting on $\mathcal{H} \otimes (\C^q)^{\otimes \nqubits}$ such that the following holds:
	\begin{enumerate}
		\item (\textbf{Self-consistency}) On average over uniformly random $(x,z,\alpha,\beta) \in \F_q^{2m + 2}$:
		\begin{equation}
			\label{eq:qld-4-12-self-cons}
			(\hat{Q}^{x,z,\alpha,\beta}_{c})_{\alice \alice'} \approx_{\delta_Q}  (\hat{Q}^{x,z,\alpha,\beta}_{c})_{\bob \alice''}\,.
		\end{equation}
\item (\textbf{Consistency with $\hat{M}$}) On average over uniformly random $(x,z,\alpha,\beta) \in \F_q^{2m + 2}$:
		\begin{gather}
			\label{eq:qld-4-12-cons-m-hat}
			(\hat{Q}^{x,z,\alpha,\beta}_{c})_{\alice \alice'} \approx_{\delta_Q} \Big ( \sum_{\substack{a,b: \\ \alpha a + \beta b = c}} \hat{M}^{(\Point,X),x}_{a} \hat{M}^{(\Point,Z),z}_{b} \Big )_{\bob \alice''} \;, \\
			\label{eq:qld-4-12-cons-m-hat2}
			(\hat{Q}^{x,z,\alpha,\beta}_{c})_{\alice \alice'} \approx_{\delta_Q}  \Big ( \sum_{\substack{a,b: \\ \alpha a + \beta b = c}} \hat{M}^{(\Point,Z),z}_{b} \, \hat{M}^{(\Point,X),x}_{a}  \Big )_{\bob \alice''}  \,.			
		\end{gather}
              \end{enumerate}
                    Furthermore, all symmetric equivalents of these approximations
      also hold.
\end{lemma}
\begin{proof}
	Define 
	\[
		\hat{Q}^{x,z,\alpha,\beta}_{c} = \sum_{\substack{a,b:\\ \alpha a + \beta b = c}} \hat{Q}^{x,z}_{a,b}.
	\]
	The self-consistency and consistency properties of $\{ \hat{Q}^{x,z,\alpha,\beta}_{c} \}$ follow from the consistency properties of $\{ \hat{Q}^{x,z}_{a,b} \}$ established by Lemma~\ref{lem:qld-4-10}, and Fact~\ref{fact:data-processing}. 
\end{proof}

\begin{lemma}\label{lem:qld-4-13}
There exists a function $\delta_{\combine}(\eps,m,d,q) = \poly(m^2\eps,md/q)$ such that the following holds. 
For every line $\line$ in
$\F_q^{2m + 2}$ and $\mH \in \{ \mH_\alice,\mH_\bob\}$ there exists a POVM $\{ \hat{Q}^\line_f \}$ acting on
$\mathcal{H} \otimes (\C^q)^{\otimes \nqubits}$ with outcomes $f\in
\deg_{md+1}(\line)$ such that on average over a uniformly random
pair $(\line, (x,z,\alpha, \beta))$ drawn from the line-point
distribution over $\F_q^{2m+2}$,
\begin{gather}
\label{eq:qld-4-13}
(\hat{Q}^\line_{[\eval_{(x,z,\alpha,\beta)}(\cdot) = a]})_{\alice \alice'} \simeq_{\delta_{\combine}} (\hat{Q}^{x,z,\alpha,\beta}_{a} )_{\bob \alice''}\;, \\
(\hat{Q}^\line_{[\eval_{(x,z,\alpha,\beta)}(\cdot) = a]})_{\bob \bob'} \simeq_{\delta_{\combine}} ( \hat{Q}^{x,z,\alpha,\beta}_{a} )_{\alice \bob''}\;,
\end{gather}
where the answer summation is over $a \in \F_q$. Moreover, if $\line$ is axis-parallel, then the outcome $f \in \deg_{d}(\line)$.
\end{lemma}

Before proving the lemma we give definitions regarding distributions over lines and points.

\begin{definition}\label{def:ith-restricted-line}
  For any $i \in \{1, \dots,
  m\}$ the \emph{$i$-th restricted axis parallel line distribution $D_{\ALine,
    i}$} is the restriction of $D_{\ALine}$ to pairs  $(\line, u)$
  where the line $\line$ is along the direction $e_i$ (i.e. is of the
  form $\line = (u_0, e_i)$. Similarly, the \emph{$i$-th restricted
  diagonal line distribution $D_{\DLine, i}$} is the restriction of
  $D_{\DLine}$ to pairs $(\line, u)$, where the line $\line = (u_0,
  s, v)$ is such that $\chi(s) = i$, and coordinates $1, \dots, i-1$ of
  $v$ are $0$. (Recall that $\chi$ is defined in~\eqref{eq:chi-func}.)
\end{definition}
  Observe that $D_{\ALine}$ is a uniform mixture over the
  distributions $D_{\ALine, i}$ for $i \in \{1, \dots, m\}$, and
  likewise for $D_{\DLine}$.
  As a consequence, any approximation that
  holds on average over the line-point distribution $D_{\Line}$ with
  error $\delta$ holds over each of the $i$-th restricted distributions
  $D_{\ALine, i}$ and $D_{\DLine, i}$ with error $2m\delta$. Similarly,
  any approximation that holds over $D_{\Line} \times D_{\Line}$ holds
  over any product of the $i$-th restricted distributions with error
  $4m^2 \delta$. For example, by \Cref{lem:qld-xz-lines} for any $\tvar_1, \tvar_2 \in \{\ALine, \DLine\}$ and $i, j \in \{1, \dots, m\}$, we have
  \begin{equation}
    \E_{(\line_X, x) \sim D_{\tvar_1, i}} \E_{(\line_Z, z) \sim D_{\tvar_2, j}} \sum_{f_X, f_Z} \bra{\hat{\psi}} (T^{\line_X, \line_Z}_{f_X, f_Z})_{\alice \alice'} \ot (I - (\hat{Q}^{x,z}_{f_X(x), f_Z(z)}))_{\bob \alice''} \ket{\hat{\psi}} = O(m^2 \delta_P)\;. \label{eq:qld-xz-lines-restricted}
    \end{equation}

    \begin{lemma}\label{lem:qld-sublines}
      There exists a distribution $D$ over tuples $(\line, \line_X, \line_Z)$ with the following properties:
      \begin{enumerate}
      \item The marginal distribution over $\line$ is the same as the marginal of the line-point distribution over $\F_q^{2m+2}$.
\item For any point of the form $u = (u_X, u_Z, \alpha, \beta)$ lying on $\line$, it holds that $u_X \in \line_X$ and $u_Z \in \line_Z$. Moreover, for $u$ chosen uniformly at random in $\line$ the marginal distribution over $(\line_X,\line_Z, u_X)$ (resp.\ over $(\line_X,\line_Z, u_Z)$) is such that the marginal over $(\ell_X,\ell_Z)$ is a mixture of $D_{\ALine, i}\times D_{\ALine, j}$ and $D_{\DLine, i}\times D_{\DLine, j}$ for $i,j\in\{1,\ldots,m\}$ and moreover, conditioned on $(\ell_X,\ell_Z)$, $u_X$ is uniformly random on $\line_X$ (resp. $u_Z$ uniformly random on $\line_Z$). 
\end{enumerate}
    \end{lemma}
		
    \begin{proof}
      We describe $D$ by giving a procedure to sample from it. We start by sampling $(\line, u)$ from the line-point distribution over $\F_q^{2m+2}$. Write $u = (u_X, u_Z, \alpha, \beta)$. 
	There are two cases to consider: the case where $\line$ is
        axis-parallel and the case where $\line$ is diagonal. In each case, we sample lines $\line_\xpt$ and $\line_\zpt$ such that the tuples  $(\line_\xpt,
        \line_\zpt,u_\xpt)$ and $(\line_\xpt, \line_\zpt,u_\zpt)$ satisfy Property~2 of the lemma.
        \begin{enumerate}
          \item Suppose $\line$ is an  axis-parallel line $(v,e_j)$ where
            $v = (v_\xpt,v_\zpt,\alpha_0,\beta_0)$. Then we sample $(v'_\xpt, i_\xpt), (v'_\zpt, i_\zpt) \in \F_q^m \times \{1,2, \ldots, m\}$ as described below, and set $\line_\xpt =
            (v'_\xpt,e_{i_\xpt}),\line_\zpt = (v'_\zpt, e_{i_\zpt})$.
            \begin{enumerate}
              \item If $1 \leq j \leq m$, then set $i_\xpt = j$ and
                choose $i_\zpt$ to be uniformly random in $\{1, \dots,
                m\}$.
              \item If  $m+1 \leq j \leq 2m$, then set $i_\xpt$ to be
                uniformly random in $\{1, \dots, m\}$ and $i_\zpt = j - m$.
              \item If $2m + 1 \leq j \leq 2m+2$, then set $i_\xpt$
                and $i_\zpt$ to both be uniformly random in $\{1, \dots, m\}$.
              \item Let $v'_\xpt$ be the canonical representative
                $L_{e_{i_\xpt}}^\lnf(v_\xpt)$, and
                likewise let $v'_\zpt$ be the canonical
                representative $L_{e_{i_\zpt}}^\lnf(v_{\zpt})$ (see
                \Cref{def:line-representative} for the definition of
                the canonical representative of a line). Note that
                $v_\xpt$ lies on the line $\line_\xpt = (v'_\xpt,
                e_{i_\xpt})$, and similarly for $v_\zpt$ and
                $\line_{\zpt} = (v'_\zpt, e_{i_\zpt})$.
              \end{enumerate}
              We claim that this construction satisfies Property~2 of the lemma. We show this first $(\line_X,\line_Z)$ and $u_X$, the other case being symmetric. If we are in case (a), then $u_X$ is a uniformly random point on the line $\line_X=(v_X, e_j)$ and $\line_Z=(v_Z,e_{i_Z})$, which is independent from $\line_X$. Thus, the marginal distribution over $(\line_X,u_X)$ conditioned on $j$ in this case is equal to $D_{\ALine, j}$ and furthermore even conditioned on $(\line_X,u_X)$, $\line_Z$ is distributed according to $D_{\ALine,i_Z}$ for a uniformly random $i_Z$. Cases (b) and (c) are analogous except that in this case even conditioned on $j$, $i_X$ is uniform. 
      \item Suppose $\ell$ is a random diagonal line $(v, s, w)$,
        where $v = (v_\xpt, v_\zpt, \alpha, \beta)$, and
        where $w = (w_\xpt, w_\zpt, \gamma, \delta) $ is a direction vector whose first nonzero coordinate is
         $j = \chi(s)$ (where $\chi$ is defined in \Cref{eq:chi-func}). We sample $\line_\xpt =  (v'_\xpt, s_\xpt,
        w'_\xpt)$ and $\line_\zpt = (v'_\zpt, s_\zpt, w'_\zpt)$ in the following
        manner:
        \begin{enumerate}
          \item If $1 \leq j \leq m$, then $w'_\xpt = w_\xpt$ and $w'_\zpt = w_\zpt$. We choose $s_\xpt$ and $s_\zpt$ such that $\chi(s_\xpt) = j$ and $\chi(s_\zpt) = 1$.
          \item If $m+1 \leq j \leq 2m$, then $w'_\xpt$ is a random
            diagonal line direction (and $s_{\xpt}$ is chosen accordingly), and $w'_\zpt = w_\zpt$, with $s_\zpt$ such that $\chi(s_\zpt) = j - m$.
          \item If $2m+1 \leq j \leq 2m+2$ then both $w_\xpt$, and
            $w_\zpt$ are uniformly random, with $s_\xpt$ and $s_\zpt$ chosen accordingly.
          \item Set $v'_\xpt = L_{w'_{\xpt}}^\lnf(v_\xpt)$ and $v'_\zpt = L_{w'_{\zpt}}^\lnf(v_\zpt)$.
          \end{enumerate}
          We claim that this construction satisfies Property~2 of the Lemma. We show this first $(\line_X,\line_Z)$ and $u_X$, the other case being symmetric.  If we are in case (a), then $u_X$ is a uniformly random point on the diagonal line $(v_X, s_X, w_X)$ which is equal to $\line_X$. Thus, the marginal distribution over $(\line_X, u_X)$ conditioned on $j$ in this case is equal to $D_{\DLine, j}$. Furthermore, conditioned on $(\line_X,u_X)$, $\line_Z$ is distributed according to $D_{\DLine,1}$. The other cases can be analyed analogously, with the only difference being how the indices $i,j$ of $\line_X$ and $\line_Z$ distributions are chosen. 
					    \end{enumerate}
	Let $D$ be the distribution over tuples $(\line, \line_X, \line_Z)$ induced by this procedure. This distribution $D$ satisfies Property 1 in the lemma by  construction, and Property 2 as argued above.
    \end{proof}

\begin{proof}[Proof of \Cref{lem:qld-4-13}]
We construct the measurements $\hat{Q}^\line_f$
out of the lines measurements $T^{\line_X, \line_Z}_{f_X, f_Z}$
guaranteed by \Cref{lem:qld-xz-lines}. To do so we first show how
to map any line-point pair $(\line, u)$ in the expanded space
$\F_q^{2m+2}$ to two pairs  $(\line_X, u_X)$ and  $(\line_Z, u_Z)$ in $\F_q^m$. We will 
define $\hat{Q}^\line$ in terms of $T^{\line_X, \line_Z}_{f_X, f_Z}$ and
show that the resulting measurement is consistent with the combined
points measurement $\hat{Q}^{x,z,\alpha, \beta}$ from \Cref{lem:qld-4-12}.

To define $\hat{Q}^\line$ we make use of the distribution $D$ from \Cref{lem:qld-sublines}. Let $D_{|\line}$ denote the distribution derived from $D$ by conditioning on the first element being $\line$. Then for every line $\line$ in $\F_q^{2m+2}$ and function $f \in \deg_{(md+1)}(\line)$, define
\[ \hat{Q}^\line_f = \E_{(\line_X, \line_Z) \sim D_{|\line}} \sum_{\substack{f_X \in \deg_{md}(\line_X), f_Z \in \deg_{md}(\line_Z):\\ f = \combine_{f_X, f_Z}}} T^{\line_X, \line_Z}_{f_X, f_Z}\;, \]
where $\combine_{f_X, f_Z}$ is defined in~\eqref{eq:combine-lines}. Note that $\{\hat{Q}^\line_f\}$ forms a valid POVM for every choice of line $\line$, and moreover that $\combine_{f_X, f_Z}$ is well-defined for the lines $\line, \line_X, \line_Z$ because of Property 2 of \Cref{lem:qld-sublines}. 

Observe that by Property~3 of \Cref{lem:qld-sublines}, whenever $\line$ is axis-parallel, $\line_X$ and $\line_Z$ are axis-parallel also. Thus, in this case, by \Cref{lem:qld-xz-lines} it holds that the outcomes $f_X$ and $f_Z$ both have degree $d$. By definition of $\combine_{f_X, f_Z}$ it holds that the outcome $f$ of $\{\hat{Q}^\line_f\}$ has degree $d$ as well.

We argue the following bound:
\begin{align}
  \E_{(\line, \line_X, \line_Z) \sim D} \E_{(x,z,\alpha, \beta) \in \line} &\sum_{f_X, f_Z} \bra{\hat{\psi}} T^{\line_X, \line_Z}_{f_X, f_Z} \ot (I - \hat{Q}^{x, z }_{f_X(x),f_Z(z)}) \ket{\hat{\psi}}  \,=\, O\big( m \cdot (\eps^{1/4} + \delta_P^{1/4} + \delta_Q^{1/4} + \delta_{\Line}^{1/2})\big)\;. \label{eq:qld-combined-lines-consistency}
\end{align}
This bound will follow from chaining the following three claims. 

\begin{claim}\label{claim:17-1}
\begin{align}
  \E_{(\line, \line_X, \line_Z) \sim D} \E_{(x,z,\alpha, \beta) \in \line} &\sum_{f_X, f_Z} \bra{\hat{\psi}} T^{\line_X, \line_Z}_{f_X, f_Z} \ot \hat{Q}^{x, z }_{f_X(x),f_Z(z)} \ket{\hat{\psi}}  \notag\\
       &\approx_{\delta_Q^{1/2}} \E_{(\line, \line_X, \line_Z) \sim D}  \E_{(x,z,\alpha, \beta) \in \line}  \sum_{f_X, f_Z} \bra{\hat{\psi}} T^{\line_X, \line_Z}_{f_X, f_Z} \ot \hat{M}^{(\Point,Z),z}_{f_Z(z)} \hat{M}^{(\Point,X),x}_{f_X(x)}  \ket{\hat{\psi}}\;. 
\end{align}
\end{claim}

\begin{proof}
Form the difference and apply the Cauchy Schwarz inequality to obtain
\begin{align*}
&\Big|  \E_{(\line, \line_X, \line_Z) \sim D} \E_{(x,z,\alpha, \beta) \in \line} \sum_{f_X, f_Z}\bra{\hat{\psi}} T^{\line_X, \line_Z}_{f_X, f_Z} \ot \big(\hat{Q}^{x, z }_{f_X(x),f_Z(z)}-\hat{M}^{(\Point,Z),z}_{f_Z(z)}\hat{M}^{(\Point,X),x}_{f_X(x)} \big) \ket{\hat{\psi}}\Big|\\
&\leq \Big( \E_{(\line, \line_X, \line_Z) \sim D} \E_{(x,z,\alpha, \beta) \in \line} \sum_{f_X, f_Z} \big\|\sqrt{T^{\line_X, \line_Z}_{f_X, f_Z} }\ot \Id \ket{\hat{\psi}} \big\|^2 \Big)^{1/2}\\
&\qquad\cdot \Big( \E_{(\line, \line_X, \line_Z) \sim D} \E_{(x,z,\alpha, \beta) \in \line} \sum_{f_X, f_Z}\big\| \sqrt{T^{\line_X, \line_Z}_{f_X, f_Z}} \ot \big(\hat{Q}^{x, z }_{f_X(x),f_Z(z)}-\hat{M}^{(\Point,Z),z}_{f_Z(z)}\hat{M}^{(\Point,X),x}_{f_X(x)} \big) \ket{\hat{\psi}}\big\|^2 \Big)^{1/2}\;.
\end{align*}
The first term above is at most $1$. For the second term, writing 
\[W^{x,z}_{a,b} = \big(\hat{Q}^{x, z }_{a,b}- \hat{M}^{(\Point,Z),z}_{b}\hat{M}^{(\Point,X),x}_{a}\big)\]
for short
we bound it as 
\begin{align*}
\E_{(\line, \line_X, \line_Z) \sim D} \E_{(x,z,\alpha, \beta) \in \line}  \sum_{f_X, f_Z}\big\| \sqrt{T^{\line_X, \line_Z}_{f_X, f_Z}} \ot W^{x,z}_{f_X(x),f_Z(z)} \ket{\hat{\psi}}\big\|^2 
&\leq \E_{(\line, \line_X, \line_Z) \sim D} \E_{(x,z,\alpha, \beta) \in \line} \sum_{a,b}\bra{\hat{\psi}} \Id \ot (W^{x,z}_{a,b})^2 \ket{\hat{\psi}} \\
&= \E_{(x,z)} \sum_{a,b}\bra{\hat{\psi}} \Id \ot (W^{x,z}_{a,b})^2 \ket{\hat{\psi}} \\
&\leq \delta_Q\;.
\end{align*}
Here the first inequality uses that for any $x,z,a,b$, $\sum_{f_X,f_X: f_X(x)=a,f_Z(z)=b} T^{\line_X, \line_Z}_{f_X, f_Z}\leq \Id$. 
In the equality in the before-last line, the expectation is over independent uniform $x,z\in \F_q^m$. The equality holds because for a line $\ell$ sampled from the line-point distribution $D$, a uniformly random point $(x,z,\alpha,\beta)$ on $\ell$ is distributed as a uniformly random point in $\F_q^{2m+2}$. The last inequality is by~\eqref{eq:qld-q-cons-m-hat-zx}.
\end{proof}

\begin{claim}\label{claim:17-2}
\begin{align}
  \E_{(\line, \line_X, \line_Z) \sim D} \E_{(x,z,\alpha, \beta) \in \line} &\sum_{f_X, f_Z} \bra{\hat{\psi}} T^{\line_X, \line_Z}_{f_X, f_Z} \ot \hat{M}^{(\Point,Z),z}_{f_Z(z)}  \hat{M}^{(\Point,X),x}_{f_X(x)} \ket{\hat{\psi}}  \notag\\
       &\approx_{m\delta_\Line^{1/2}} \E_{(\line, \line_X, \line_Z) \sim D}  \E_{(x,z,\alpha, \beta) \in \line} \sum_{f_X,f_Z} \bra{\hat{\psi}} T^{\line_X, \line_Z}_{f_X, f_Z} \ot  \hat{M}^{(\Point,Z),Z}_{f_Z(z)} \ket{\hat{\psi}}\;. 
\end{align}
\end{claim}

\begin{proof}
We form the difference and apply the Cauchy-Schwarz inequality
\begin{align*}
 \Big|\E_{(\line, \line_X, \line_Z) \sim D} \E_{(x,z,\alpha, \beta) \in \line} &\sum_{f_X, f_Z} \bra{\hat{\psi}} T^{\line_X, \line_Z}_{f_X, f_Z} \ot \hat{M}^{(\Point,Z),z}_{f_Z(z)} (\Id-\hat{M}^{(\Point,X),x}_{f_X(x)} ) \ket{\hat{\psi}}\Big|\\
&\leq \Big(\E_{(\line, \line_X, \line_Z) \sim D} \E_{(x,z,\alpha, \beta) \in \line} \sum_{f_X, f_Z} \big\| \sqrt{ T^{\line_X, \line_Z}_{f_X, f_Z}} \ot  \hat{M}^{(\Point,Z),z}_{f_Z(z)} \ket{\hat{\psi}}\big\|^2 \Big)^{1/2}\\
&\qquad\cdot \Big(\E_{(\line, \line_X, \line_Z) \sim D} \E_{(x,z,\alpha, \beta) \in \line} \sum_{f_X, f_Z} \big\| \sqrt{ T^{\line_X, \line_Z}_{f_X, f_Z}} \ot  (\Id-\hat{M}^{(\Point,X),x}_{f_X(x)}) \ket{\hat{\psi}}\big\|^2 \Big)^{1/2}\;.
\end{align*}
The first term above is at most $1$. The second term can be bounded as
\begin{align}
\E_{(\line, \line_X, \line_Z) \sim D} \E_{(x,z,\alpha, \beta) \in \line} &\sum_{f_X, f_Z} \big\| \sqrt{ T^{\line_X, \line_Z}_{f_X, f_Z}} \ot  (\Id-\hat{M}^{(\Point,X),x}_{f_X(x)}) \ket{\hat{\psi}}\big\|^2 \notag\\
&= \E_{(\line, \line_X, \line_Z) \sim D} \E_{(x,z,\alpha, \beta) \in \line} \sum_{a} \bra{\hat{\psi}}  T^{\line_X, \line_Z,x}_{a} \ot  (\Id-\hat{M}^{(\Point,X),x}_{a})^2 \ket{\hat{\psi}}\notag\\
&= 1-\E_{(\line, \line_X, \line_Z) \sim D} \E_{(x,z,\alpha, \beta) \in \line} \sum_{a} \bra{\hat{\psi}}  T^{\line_X, \line_Z,x}_{a} \ot \hat{M}^{(\Point,X),x}_{a} \ket{\hat{\psi}}\;,\label{eq:17-4}
\end{align}
where we used the shorthand $T^{\line_X, \line_Z,x}_{a} = \sum_{f_X,f_Z: f_X(x)=a} T^{\line_X, \line_Z}_{f_X, f_Z}$ and the second equality uses that $\{\hat{M}^{(\Point,X),x}_{a}\}$ is a projective measurement. Using the definition 
\[	T^{\line_\xpt,\line_\zpt}_{f_\xpt,f_\zpt} = \hat{M}^{(\Line,\xpt),\line_\xpt}_{f_\xpt} \cdot \hat{M}^{(\Line,\zpt),\line_\zpt}_{f_\zpt}  \cdot \hat{M}^{(\Line,\xpt),\line_\xpt}_{f_\xpt} \]
and the fact that $\{\hat{M}^{(\Line,\xpt),\line_\xpt}_{f_\xpt} \}$ is projective we get
\begin{align}
1-\eqref{eq:17-4}&=\E_{(\line, \line_X, \line_Z) \sim D} \E_{(x,z,\alpha, \beta) \in \line} \sum_{a} \bra{\hat{\psi}}  \hat{M}^{(\Line,\xpt),\line_\xpt}_{[\eval_x(\cdot)=a]} \ot \hat{M}^{(\Point,X),x}_{a} \ket{\hat{\psi}}\;.\label{eq:a17-8}
\end{align}
By Property~2 of \Cref{lem:qld-sublines} the right-hand side of~\eqref{eq:a17-8}
 is a mixture of terms of the form
\begin{equation}\label{eq:a17-9}
 \E_{(\line_X, x) \sim D_{\tvar, i}}  \sum_{a} \bra{\hat{\psi}}  \hat{M}^{(\Line,\xpt),\line_\xpt}_{[\eval_x(\cdot)=a]} \ot \hat{M}^{(\Point,X),x}_{a} \ket{\hat{\psi}}\;, 
\end{equation}
for $\tvar \in \{\ALine, \DLine\}$ and $i \in \{1, \dots, m\}$. 
Hence by the discussion following Definition~\ref{def:ith-restricted-line} and  item 3 of Lemma~\ref{lem:qld-comm-line-cons} it follows that
\[ \eqref{eq:a17-9} \,=\, O(m^2 \delta_\Line)\;. \]
\end{proof}

\begin{claim}\label{claim:17-3}
\begin{align}
  \E_{(\line, \line_X, \line_Z) \sim D}\E_{(x,z,\alpha, \beta) \in \line}\sum_{f_X,f_Z} \bra{\hat{\psi}} T^{\line_X, \line_Z}_{f_X, f_Z} \ot  \hat{M}^{(\Point,Z),z}_{f_Z(z)} \ket{\hat{\psi}} 
       &\approx_{\sqrt{m}\cdot(\delta_P^{1/4}+\delta_Q^{1/4}+\eps^{1/4})}1 \;. 
\end{align}
\end{claim}

\begin{proof}
At first we proceed similarly to the proof of Claim~\ref{claim:17-2}, bounding the difference using Cauchy-Schwarz as
\begin{align*}
&\Big| \E_{(\line, \line_X, \line_Z) \sim D} \E_{(x,z,\alpha, \beta) \in \line} \sum_{f_X,f_Z} \bra{\hat{\psi}} T^{\line_X, \line_Z}_{f_X, f_Z} \ot  (\Id-\hat{M}^{(\Point,Z),z}_{f_Z(z)}) \ket{\hat{\psi}}\Big|\\
&\leq\Big(\E_{(\line, \line_X, \line_Z) \sim D} \E_{(x,z,\alpha, \beta) \in \line} \sum_{f_X, f_Z} \big\| \sqrt{ T^{\line_X, \line_Z}_{f_X, f_Z}} \ot  \Id \ket{\hat{\psi}}\big\|^2 \Big)^{1/2}\\
&\qquad\cdot \Big(\E_{(\line, \line_X, \line_Z) \sim D} \E_{(x,z,\alpha, \beta) \in \line} \sum_{f_X, f_Z} \big\| \sqrt{ T^{\line_X, \line_Z}_{f_X, f_Z}} \ot  (\Id-\hat{M}^{(\Point,Z),z}_{f_Z(z)}) \ket{\hat{\psi}}\big\|^2 \Big)^{1/2}\;.
\end{align*}
The first term is at most $1$. The second, as in the proof of Claim~\ref{claim:17-2}, equals
\begin{align}
1-\E_{(\line, \line_X, \line_Z) \sim D} &\E_{(x,z,\alpha, \beta) \in \line} \sum_{b} \bra{\hat{\psi}}  T^{\line_X, \line_Z, z}_{b} \ot \hat{M}^{(\Point,Z),z}_{b} \ket{\hat{\psi}}\notag\\
&=1-\E_{(\line_X,\line_Z)\sim D_{\tvar,i}\times D_{\tvar,j}}
\E_{z\in \line_Z } \sum_{b} \bra{\hat{\psi}}  T^{\line_X, \line_Z, z}_{b}  \ot \hat{M}^{(\Point,Z),z}_{b} \ket{\hat{\psi}}\label{eq:a17-5}\\
&=1-\E_{(\line_X,\line_Z)\sim D_{\tvar,i}\times D_{\tvar,j}}
\E_{z\in \line_Z }\E_{x\in \line_X} \sum_{a,b} \bra{\hat{\psi}}  T^{\line_X, \line_Z}_{[\eval_x(\cdot)=a,\eval_z(\cdot)=b]}  \ot \hat{M}^{(\Point,Z),z}_{b} \ket{\hat{\psi}}\label{eq:a17-6}
\end{align}
where we use the shorthand $T^{\line_X, \line_Z, z}_{b} = \sum_{f_X} T^{\line_X, \line_Z}_{f_X,[\eval_z(\cdot)=b]}$ and $D_{\tvar,i}\times D_{\tvar,j}$ to denote a mixture over all such distributions where $\tvar\in\{\ALine,\DLine\}$ and $i,j\in\{1,\ldots,m\}$. Equality holds in~\eqref{eq:a17-5} because by item 2 of Lemma~\ref{lem:qld-sublines} for $(\ell,\ell_X,\ell_Z)\sim D$ the marginal distribution on $(\ell_X,\ell_Z)$ is a mixture of the product distributions  $D_{\ALine, i}\times D_{\ALine, j}$ and $D_{\DLine, i}\times D_{\DLine,j}$, for $i,j\in \{1,\ldots,m\}$. Moreover, conditioned on $(\ell_X,\ell_Z)$, $z$ is a uniformly random point on $\ell_X$. The last equality in~\eqref{eq:a17-6} holds by definition of 
\begin{align*}
T^{\line_X, \line_Z, z}_{b}&= \sum_{f_X} T^{\line_X, \line_Z}_{f_X,[\eval_z(\cdot)=b]}\\
&= \Es{x \in \line_X} \sum_a \sum_{f_X: f_X(x)=a} T^{\line_X, \line_Z}_{f_X,[\eval_z(\cdot)=b]}\\
&= \Es{x \in \line_X} \sum_a  T^{\line_X, \line_Z}_{[\eval_x(\cdot)=a,\eval_z(\cdot)=b]}\;.
\end{align*}
We next write
\begin{align*}
1-\eqref{eq:a17-6}&\approx_{m(\sqrt{\delta_P}+\sqrt{\delta_Q})}\E_{(\line_X,\line_Z)\sim D_{\tvar,i}\times D_{\tvar,j}}\E_{x\in \line_X} \E_{z\in \line_Z } \sum_{a,b} \bra{\hat{\psi}} \hat{M}^{(\Point,X),x}_{a}\hat{M}^{(\Point,Z),z}_{b}  \ot \hat{M}^{(\Point,Z),z}_{b} \ket{\hat{\psi}} \\
&= \sum_{b} \bra{\hat{\psi}} \hat{M}^{(\Point,Z),z}_{b}  \ot \hat{M}^{(\Point,Z),z}_{b} \ket{\hat{\psi}} \\
&\geq 1-\sqrt{\eps}\;,
\end{align*}
where the first approximation follows from the discussion following Definition~\ref{def:ith-restricted-line} applied to the bounds from Lemma~\ref{lem:qld-sublines} and Lemma~\ref{lem:qld-4-10}, the second line uses that $\{\hat{M}^{(\Point,X),x}_{a}\}$ is a measurement, and the last is by self-consistency of the points measurements (Lemma~\ref{lem:qld-comm-cons}).
\end{proof}

Together, Claim~\ref{claim:17-1}, Claim~\ref{claim:17-2} and Claim~\ref{claim:17-3} show~\eqref{eq:qld-combined-lines-consistency}. To complete the proof of~\eqref{eq:qld-4-13}, observe that by Property~2 of \Cref{lem:qld-sublines} the right-hand side of \Cref{eq:qld-combined-lines-consistency} is a mixture of terms of the form
Together, Claim~\ref{claim:17-1}, Claim~\ref{claim:17-2} and Claim~\ref{claim:17-3} show {eq:qld-combined-lines-consistency}. To complete the proof of~\eqref{eq:qld-4-13}, observe that by Property~2 of \Cref{lem:qld-sublines} the right-hand side of \Cref{eq:qld-combined-lines-consistency} is a mixture of terms of the form
\[ \E_{(\line_X, x) \sim D_{\tvar_1, i}} \E_{(\line_Z, z) \sim D_{\tvar_2, j}} \sum_{f_X, f_Z} \bra{\hat{\psi}} (T^{\line_X, \line_Z}_{f_X, f_Z})_{\alice \alice'} \ot (I - (\hat{Q}^{x,z}_{f_X(x), f_Z(z)})) \ket{\hat{\psi}} \]
for $\tvar_1, \tvar_2 \in \{\ALine, \DLine\}$ and $i, j \in \{1, \dots, m\}$. Hence using~\Cref{eq:qld-xz-lines-restricted} it follows that
\[ \eqref{eq:qld-combined-lines-consistency} \,=\, O(m^2 \delta_P)\;. \]
This establishes the first relation~\eqref{eq:qld-4-13}. An analogous argument shows the second relation.
















	
\end{proof}


\begin{lemma}\label{lem:qld-4-7}
There exists a function $\delta_S(\eps,m,d,q) = a(md)^a (\eps^b + q^{-b} + 2^{-bmd})$ for some universal constants $a > 1, 0 < b < 1$ such that the following holds. For $\mH \in \{\mH_\alice,\mH_\bob\}$ there exists a projective measurement $\{ \hat{S}_{g_\xpt,g_\zpt} \}$ acting on $\mathcal{H} \otimes (\C^q)^{\otimes \nqubits}$ with outcomes consisting of pairs $(g_\xpt,g_\zpt)$ of polynomials each in $\ideg_{d,m}(\F_q)$ such that for $W \in \{X,Z\}$, on average over uniformly random $u \in \F_q^m$,
\begin{align}
  (\hat{S}_{[\eval_{u}(\cdot_{W}) = a]})_{\alice \alice'} &\simeq_{\delta_S} (\hat{M}^{(\Point, W), u}_{a})_{\bob \alice''} \label{eq:qld-s-point-con-alice} \\ 
    (\hat{S}_{[\eval_{u}(\cdot_{W}) = a]})_{\bob \bob'} &\simeq_{\delta_S} (\hat{M}^{(\Point, W), u}_{a})_{\alice \bob''} \label{eq:qld-s-point-con-bob},
\end{align}
where the notation $\eval_{u}(\cdot_W)$ means the evaluation of the first outcome $g_X$ of $\hat{S}$ if $W = X$, or the second outcome $g_Z$ if $W = Z$, at the point $u$.
\end{lemma}


\begin{proof}
	Lemma~\ref{lem:qld-4-13} establishes the existence of a strategy $\strategy_\ld = (\hat{\psi}, \{
        \hat{Q}^\line_f \} \cup \{ \hat{Q}^{x,z,\alpha,\beta}_a \} )$
        for the classical low-degree test $\game^\ld_\ldparams$ for
        $\ldparams = (q,2m + 2,d, 1)$, with value $1 -
        \delta_{\combine}$. This strategy is not necessarily projective because the measurement $\{\hat{Q}^{\line}_f\}$ is not guaranteed to be; however, we may projectivize it by applying Naimark's theorem to this measurement. This preserves the consistency ($\simeq$) relation obtained as the conclusion of \Cref{lem:qld-4-13}. Theorem~\ref{lem:ld-soundness} applied to this projectivized strategy implies the
        existence of a POVM $\{ \hat{S}_g \}$ with
        outcomes $g \in \ideg_{d,2m+2}(\F_q)$ that is
        $\delta_\ld$-self-consistent and consistent with the
        ``points'' measurements $\{ \hat{Q}^{x,z,\alpha,\beta}_a\}$,
        where $\delta_\ld = a(md)^a(\delta_\combine^b + q^{-b} + 2^{-bmd})$ for some universal constants $a > 1, 0 < b < 1$ is given
        by Theorem~\ref{lem:ld-soundness}. We may apply Naimark's theorem once again to $\{\hat{S}_g\}$ to obtain a projective measurement with the same consistency guarantees, and henceforth, we will use $\{\hat{S}_g\}$ to refer to this projective measurement.
	
	We now argue that with high probability, the polynomial $g$ returned by the measurement $\{\hat{S}_g \}$ has the form $g(x,z,\alpha,\beta) = \alpha g_\xpt(x) + \beta g_\zpt (z)$ for polynomials $g_\xpt, g_\zpt \in \ideg_{d,m}(\F_q)$. Let $\mathcal{G}$ denote the set of such polynomials. On average over uniformly random $(x,z,\alpha,\beta)$, 
	\begin{align}
          (\hat{S}_g)_{\alice \alice'} &= (\hat{S}_g)^2_{\alice \alice'} \notag \\
                                       &\approx_{\delta_\ld} (\hat{S}_g)_{\alice \alice'} \otimes (\hat{Q}^{x,z,\alpha,\beta}_{g(x,z,\alpha,\beta)})_{\bob \alice''}\notag\\
		&\approx_{\delta_Q} (\hat{S}_g)_{\alice \alice'} (\hat{Q}^{x,z,\alpha,\beta}_{g(x,z,\alpha,\beta)})_{\alice \alice'}\notag \\
		&\approx_{\delta_Q} (\hat{S}_g)_{\alice \alice'} \otimes \Big ( \sum_{b,b'} 1_{\alpha b + \beta b' = g(x,z,\alpha,\beta)} \, \hat{M}^{(\Point,X),x}_{b} \hat{M}^{(\Point,Z),z}_{b'} \Big)_{\bob \alice''} \label{eq:qld-g-42} \\
		&\approx_{\delta_Q} (\hat{S}_g)_{\alice \alice'} \otimes \Big ( \sum_{b,b'} 1_{\alpha b + \beta b' = g(x,z,\alpha,\beta)} \, \hat{M}^{(\Point,Z),z}_{b'} \hat{M}^{(\Point,X),x}_{b} \Big)_{\bob \alice''} \;.\label{eq:qld-g-43} 		
	\end{align}
	 Here the first approximation uses the projectivity of $\{\hat{S}_g\}$, and the first approximation uses the consistency of $\{\hat{S}_g\}$ with  $\{\hat{Q}^{x,z,\alpha,\beta}_a\}$ from Theorem~\ref{lem:ld-soundness}, together with Fact~\ref{fact:agreement} to switch from $\simeq$ to $\approx$, and Fact~\ref{fact:add-a-proj2}. The second, third and fourth approximations follow from the consistency properties of $\{\hat{Q}^{x,z,\alpha,\beta}_a\}$ from Lemma~\ref{lem:qld-4-12} combined with Fact~\ref{fact:add-a-proj2}. Note that we will need both \Cref{eq:qld-g-42} and \Cref{eq:qld-g-43} for the subsequent calculations.
	
	For all $g$ define $ \ket{\hat{\psi}_g} =(\hat{S}_g)_{\alice \alice'} \ket{\hat{\psi}}$. 
	Then on average over uniformly random $(x,z,\alpha,\beta)$, 
	  \begin{align}
	  	&\E_{x,z,\alpha,\beta} \norm{ \Big ( \sum_{b,b'} 1_{\alpha b + \beta b' = g(x,z,\alpha,\beta)} \, \hat{M}^{(\Point,X),x}_{b} \hat{M}^{(\Point,Z),z}_{b'} \Big)_{\bob \alice''} \ket{\hat{\psi}_g} }^2\notag \\
		=_{\phantom{1/q}}& \E_{x,z,\alpha,\beta} \sum_{b,b',b''} 1_{\alpha b + \beta b' = g(x,z,\alpha,\beta)} 1_{\alpha b + \beta b'' = g(x,z,\alpha,\beta)} \bra{\hat{\psi}_g}  \hat{M}^{(\Point,Z),z}_{b''} \hat{M}^{(\Point,X),x}_{b} \hat{M}^{(\Point,Z),z}_{b'} \ket{\hat{\psi}_g} \notag\\
		\approx_{1/q}& \E_{x,z,\alpha,\beta} \sum_{b,b'} 1_{\alpha b + \beta b' = g(x,z,\alpha,\beta)} \bra{\hat{\psi}_g}  \hat{M}^{(\Point,Z),z}_{b'} \hat{M}^{(\Point,X),x}_{b} \hat{M}^{(\Point,Z),z}_{b'} \ket{\hat{\psi}_g}
		\label{eq:qld-g-prime}
	  \end{align}
	  where in the first equality, we used the fact that $ \{  \hat{M}^{(\Point,X),x}_{b} \}$ is projective, and in the approximation, we used that for fixed $x,z,\alpha,\beta$ with $\beta \neq 0$, if $\alpha b + \beta b' = g(x,z,\alpha,\beta)$ and $\alpha b + \beta b'' = g(x,z,\alpha,\beta)$, then $b' = b''$. The probability that $\beta = 0$ is $1/q$, and the term under the approximation has an absolute value of at most $1$, so the error incurred is at most $1/q$.

	 We will now show that the right-hand side of \Cref{eq:qld-g-prime} is small for polynomials $g$ that are not of the desired form. 
	Fix $b,b' \in \F_q$, and let $h_{b,b'}(x,z,\alpha,\beta) = g(x,z,\alpha,\beta) - (\alpha b + \beta b')$. Observe that since $g$ has individual degree $d$, so does $h_{b,b'}$. Write 
	\[
	h_{b,b'}(x,z,\alpha,\beta) = \sum_{i,j=0}^{d}\, h_{b,b',i,j}(x,z) \,  \alpha^i \beta^j 
	\]
	for some polynomials $h_{b,b',i,j}(x,z)$ of individual degree at most $d$. 
	Let $\mathcal{G}' \subset \ideg_{d,2m+2}(\F_q)$ denote the polynomials that are linear in $\alpha, \beta$, i.e., $g = \alpha g_1(x,z) + \beta g_2(x,z)$. Fix a $g \notin \mathcal{G}'$. Then $h_{b,b'}(x,z,\alpha,\beta)$ is a nonzero polynomial of total degree at most $(2m+2)d$, and in particular there must exist an $(i,j)$ such that the polynomial $h_{b,b',i,j}(x,z)$ is not the identically zero polynomial. Call a pair $(x,z) \in \F_q^2$ \emph{good} if $h_{b,b'}(x,z,\alpha,\beta)$ is a nonzero polynomial function of $\alpha,\beta$, and \emph{bad} otherwise. The probability that $(x,z)$ is bad is at most $(2m+2)d/q$ by the Schwartz-Zippel lemma. Otherwise, conditioned on a good $(x,z)$ pair, the probability that $h_{b,b'}(x,z,\alpha,\beta)=0$ over the choice of $\alpha,\beta$ is at most $(2m+2)d/q$, again by the Schwartz-Zippel lemma. Thus for $g \notin \mathcal{G}'$ we can upper-bound~\eqref{eq:qld-g-prime} by $\frac{2(2m+2)d}{q} \norm{ \ket{\hat{\psi}_g}}^2$. Combined with~\eqref{eq:qld-g-42} we get that
	\begin{equation}
	\label{eq:qld-g-prime-bound}
		\sum_{g \notin \mathcal{G}'} \norm{ (\hat{S}_g)_{\alice \alice'} \ket{\hat{\psi}}}^2 \leq O(\delta_\ld + \delta_Q + md/q).
	\end{equation}
	




	Next, we argue that not only is $g$ linear in $\alpha, \beta$
        with high probability, but also that $g = \alpha g_\xpt +
        \beta g_\zpt$ where $g_\xpt$ depends on $x$ only and $g_\zpt$
        depends on $z$ only. Fix a polynomial $g \in \mathcal{G}'$,
        i.e., $g(x,z,\alpha,\beta) = \alpha g_1(x,z) + \beta g_2(x,z)$
        for polynomials $g_1, g_2 \in \ideg_{d,2m}(\F_q)$. For all $x,z \in \F_q^m$, $\alpha, \beta, b,b' \in \F_q$,
	\begin{align*}
		1_{\alpha b + \beta b' = \alpha g_1(x,z) + \beta g_2(x,z)}
		= 1_{b = g_1(x,z)} 1_{b' = g_2(x,z)} + 1_{b \neq g_1(x,z) \, \vee \, b' \neq g_2(x,z)} \cdot 1_{\alpha b + \beta b' = \alpha g_1(x,z) + \beta g_2(x,z)}\;.
	\end{align*}
	By the Schwartz-Zippel lemma,  
	\begin{equation}
	\label{eq:qld-g-separable}
	\E_{\alpha, \beta} 1_{b \neq g_1(x,z) \, \vee \, b' \neq g_2(x,z)} \cdot 1_{\alpha b + \beta b' = \alpha g_1(x,z) + \beta g_2(x,z)} \leq 1/q\;.
	\end{equation}
	Let $\mathcal{G}' = \mathcal{G}'_\xpt \cup \mathcal{G}'_\zpt \cup \mathcal{G}$, where $\mathcal{G}'_\xpt$ (resp. $\mathcal{G'}_\zpt$) is the set of polynomials $g \in \mathcal{G}'$ where $g_1(x,z)$ depends nontrivially on $z$ (resp. $g_2(x,z)$ depends nontrivially on $x$). Fix $g \in \mathcal{G}'_\zpt$. We can upper-bound~\eqref{eq:qld-g-prime} by
	\begin{align}
		&\frac{1}{q} \norm{ \ket{\hat{\psi}_g}}^2 + \E_{x,z} \sum_{b,b'} 1_{b = g_1(x,z)} 1_{b' = g_2(x,z)} \bra{\hat{\psi}_g}   \hat{M}^{(\Point,Z),z}_{b'} \hat{M}^{(\Point,X),x}_{b} \hat{M}^{(\Point,Z),z}_{b'}  \ket{\hat{\psi}_g} \label{eq:qld-g-48} \\
		&\leq \frac{1}{q} \norm{ \ket{\hat{\psi}_g}}^2 + \E_{z} \sum_{b'} \Big(\E_x 1_{b' = g_2(x,z)} \Big) \, \bra{\hat{\psi}_g}  \hat{M}^{(\Point,Z),z}_{b'}  \ket{\hat{\psi}_g}
		\label{eq:qld-g-2}
\end{align}
	where the indicator $1_{b = g_1(x,z)}$ is removed using
        positivity in the inequality. To bound the second term
      in~\eqref{eq:qld-g-2}, we use an argument similar to the one we used to
      derive~\eqref{eq:qld-g-prime-bound}. That is, say that a value
      of $z$ is \emph{good} if $g_2(\cdot ,z)$ is not a constant
      function. The probability that $z$ is good is at least $1 -
      (2m+2)d/q$ by the Schwartz-Zippel lemma and the assumption that $g \in \mathcal{G}'_{\zpt}$ (and so $g_2$ depends nontrivially on $x$). Moreover, for each good $z$ and
      for each $b'$ the expectation $\E_x 1_{b' = g_2(x,z)}$ is at
      most $(2m+2)d/q$. Hence, we conclude that if $g_2(x,z)$ depends on $x$, then~\eqref{eq:qld-g-2} is at most $ \frac{2(2m+2)d + 1}{q} \norm {\ket{\hat{\psi}_g}}^2$, and thus 
	\begin{equation}
	\label{eq:qld-g-prime-xpt-bound}
		\sum_{g \in \mathcal{G}'_\zpt} \norm{ (\hat{S}_g)_{\alice \alice'} \ket{\hat{\psi}}}^2 \leq O(\delta_\ld + \delta_Q + md/q).
	\end{equation}	
	
	By starting with~\eqref{eq:qld-g-43} instead (i.e., the order of $M^{(\Point,Z),z}_{b'}$ and $M^{(\Point,X),x}_b$ are switched), we can perform similar reasoning to deduce that
	\begin{equation}
	\label{eq:qld-g-prime-zpt-bound}
		\sum_{g \in \mathcal{G}'_\xpt} \norm{ (\hat{S}_g)_{\alice \alice'} \ket{\hat{\psi}}}^2 \leq O(\delta_\ld + \delta_Q + md/q).
	\end{equation}	
	We thus obtain
	\begin{align}
		\sum_{g \notin \mathcal{G}} \norm{ (\hat{S}_g)_{\alice \alice'} \ket{\hat{\psi}}}^2 &\leq 
		\sum_{g \notin \mathcal{G}'}\norm{ (\hat{S}_g)_{\alice \alice'} \ket{\hat{\psi}}}^2 +  \sum_{g \in \mathcal{G}'_\xpt}\norm{ (\hat{S}_g)_{\alice \alice'} \ket{\hat{\psi}}}^2 + \sum_{g \in \mathcal{G}'_\zpt}\norm{ (\hat{S}_g)_{\alice \alice'} \ket{\hat{\psi}}}^2 \\
		&\leq O(\delta_\ld + \delta_Q + md/q). 
		\label{eq:qld-g-non-separable}
	\end{align}
        
	Let $\delta_{\mathcal{G}}$ be the error in the right hand side of \Cref{eq:qld-g-non-separable}. Define a projective sub-measurement $\{\hat{S}_{g_X, g_Z}\}$ with outcomes $g_X, g_Z \in \ideg_{d,m}(\F_q)$ by
        \begin{equation}
          \hat{S}_{g_X, g_Z} = \hat{S}_{\combine_{g_X, g_Z}}.
        \end{equation}
        A bound on the incompleteness of $\hat{S}_{g_X, g_Z}$ follows from~\eqref{eq:qld-g-non-separable} and the projectivity of $\hat{S}_g$:
        \begin{equation} 1 - \sum_{g_X, g_Z} \bra{\hat{\psi}} (\hat{S}_{g_X, g_Z})_{\alice \alice'} \ket{\hat{\psi}} = \sum_{g \notin \mathcal{G}} \bra{\hat{\psi}} (\hat{S}_g)_{\alice \alice'} \ket{\hat{\psi}} = \sum_{g \notin \mathcal{G}} \norm{ (\hat{S}_g)_{\alice \alice'} \ket{\hat{\psi}} }^2 \leq\delta _{\mG}. \label{eq:qld-sgg-completeness} \end{equation}
        We now bound the consistency of $\{\hat{S}_{g_X, g_Z}\}$ with the points measurements $\{\hat{M}^{(\Point, W), u}_{a}\}$.
	\begin{align}
		1 &= \sum_g\bra{\hat{\psi}}  (\hat{S}_g)_{\alice \alice'} \ket{\hat{\psi}} \label{eq:qld-s-good-and-bad} \\
		&\approx_{\delta_{\mG}} \sum_{g \in \mathcal{G}} \bra{\hat{\psi}} (\hat{S}_g)_{\alice \alice'} \ket{\hat{\psi}} \label{eq:qld-s-good-g}\\
                  &\approx_{(\delta_Q)^{1/2}} \E_{x,z,\alpha,\beta} \sum_{g_\xpt, g_\zpt}\bra{\hat{\psi}}  (\hat{S}_{g_\xpt,g_\zpt})_{\alice \alice'} \otimes \Big ( \sum_{b,b'} 1_{\alpha b + \beta b' = \alpha g_\xpt (x) + \beta g_\zpt(z)} \, \hat{M}^{(\Point,Z),z}_{b'} \hat{M}^{(\Point,X),x}_{b} \Big)_{\bob \alice''} \ket{\hat{\psi}} \label{eq:qld-s-mm}\\
                  &\approx _{(\delta_Q)^{1/2}} \E_{x,z,\alpha,\beta} \sum_{g_\xpt, g_\zpt}\bra{\hat{\psi}}  (\hat{S}_{g_\xpt,g_\zpt})_{\alice \alice'} \otimes \Big ( \sum_{b,b', b''} 1_{\alpha b + \beta b' = \alpha g_\xpt (x) + \beta g_\zpt(z) = \alpha b'' + \beta b'} \cdot\, \notag \\
          &\hspace{12em} \hat{M}^{(\Point, X), x}_{b''} \hat{M}^{(\Point,Z),z}_{b'} \hat{M}^{(\Point,X),x}_{b} \Big)_{\bob \alice''} \ket{\hat{\psi}} \label{eq:qld-s-mmm} \\
      &\approx _{1/q} \E_{x,z,\alpha,\beta} \sum_{g_\xpt, g_\zpt}\bra{\hat{\psi}}  (\hat{S}_{g_\xpt,g_\zpt})_{\alice \alice'} \otimes \Big ( \sum_{b,b'} 1_{\alpha b + \beta b' = \alpha g_\xpt (x) + \beta g_\zpt(z) } \cdot \,\notag\\
    &\hspace{12em}\hat{M}^{(\Point, X), x}_{b} \hat{M}^{(\Point,Z),z}_{b'} \hat{M}^{(\Point,X),x}_{b} \Big)_{\bob \alice''} \ket{\hat{\psi}} \label{eq:qld-s-mmm-positive} \\
		&\approx_{ 1/q}\E_{x,z}  \sum_{g_\xpt, g_\zpt} \bra{\hat{\psi}} (\hat{S}_{g_\xpt,g_\zpt})_{\alice \alice'} \otimes \Big ( \hat{M}^{(\Point, X), x}_{g_\xpt(x)} \hat{M}^{(\Point,Z),z}_{g_\zpt(z)} \hat{M}^{(\Point,X),x}_{g_\xpt(x)} \Big)_{\bob \alice''} \ket{\hat{\psi}}. \label{eq:qld-xzx-disagree}
	\end{align}
	Here \Cref{eq:qld-s-good-g} follows from~\eqref{eq:qld-sgg-completeness}. \Cref{eq:qld-s-mm} follows by approximation~\eqref{eq:qld-g-43} together with the Cauchy-Schwarz inequality and the projectivity of $\hat{S}_g$:
        \begin{align*}
          &\Big| \E_{x,z,\alpha,\beta} \sum_{g_\xpt, g_\zpt} \bra{\hat{\psi}} \Big( (\hat{S}_{g_\xpt, g_\zpt})_{\alice \alice'} \cdot (I -  \Big( \sum_{b,b'} 1_{\alpha b + \beta b' = \alpha g_\xpt (x) + \beta g_\zpt(z)} \, \hat{M}^{(\Point,Z),z}_{b'} \hat{M}^{(\Point,X),x}_{b} \Big)_{\bob \alice''}\Big) \ket{\hat{\psi}} \Big| \\
          \leq& \sqrt{\E_{x,z,\alpha,\beta} \sum_{g_\xpt, g_\zpt} \bra{\hat{\psi}} (\hat{S}_{g_\xpt, g_\zpt})_{\alice \alice'}^2 \ket{\hat{\psi}}} \\
          &\quad \cdot \sqrt{\E_{x,z,\alpha,\beta} \sum_{g_\xpt, g_\zpt} \norm{ (\hat{S}_{g_\xpt, g_\zpt})_{\alice \alice'} \ot 
\Big(I - \Big( \sum_{b,b'} 1_{\alpha b + \beta b' = \alpha g_\xpt(x) + \beta g_\zpt(z)} \, \hat{M}^{(\Point, Z),z}_{b'} \hat{M}^{(\Point, X), x}_{b} \Big)_{\bob \alice''} \Big)\ket{\hat{\psi}} }^2 } \\
          \leq& 1 \cdot \sqrt{\delta_Q}.
        \end{align*}
        \Cref{eq:qld-s-mmm} follows by a similar calculation\anote{TODO}. To pass from here to \Cref{eq:qld-s-mmm-positive}, we argue that for $\alpha \neq 0$, the indicator forces $b = b''$, and the terms where $\alpha = 0$ are bounded in absolute value by $1/q$. Finally, to reach \Cref{eq:qld-xzx-disagree}, we apply \Cref{eq:qld-g-separable} to the expectation of the indicator over $\alpha, \beta$. \Cref{eq:qld-xzx-disagree} tells us that
        \begin{equation} \big(\hat{S}_{[\eval_{x}(\cdot) = a, \eval_{z}(\cdot) = b]}\big)_{\alice \alice'} \simeq_{O(\delta_{\mG} + (\delta_Q)^{1/2} + 1/q)} \big( \hat{M}^{(\Point, X), x}_{a} \hat{M}^{(\Point, Z), z}_{b} \hat{M}^{(\Point, X), x}_{a}\big)_{\bob \alice''}. \label{eq:qld-sgg-mhat-sandwich} \end{equation}

        This is almost the desired conclusion of the lemma. The only issue is that $\{\hat{S}_{g_X, g_Z}\}$ is a sub-measurement. To address this, we complete $\hat{S}_{g_X, g_Z}$ by adding $1 - \sum_{g_X, g_Z} \hat{S}_{g_X, g_Z}$ to an arbitrary measurement element. By \Cref{eq:qld-sgg-completeness}, it follows that the consistency relation in \Cref{eq:qld-sgg-mhat-sandwich} holds for the completed measurement with an error increased by $\delta_{\mG}$. Let $\delta_S$ be this new, increased, error, and observe that $\delta_S = O(\delta_{\mG} + (\delta_Q)^{1/2} + d/q)$. Substituting in the bound on $\delta_{\mG}$ from \Cref{eq:qld-g-non-separable}, we obtain that $\delta_S = O(\delta_\ld + (\delta_Q)^{1/2} + md/q)$ has the form $a' (md)^{a'} (\eps^{b'} + q^{-b'} + 2^{-b'md})$ for some universal constants $a' > 1, 0 < b' < 1$. We thus conclude by an application of \Cref{fact:data-processing} that \Cref{eq:qld-s-point-con-alice} holds for $W = X$.

        To obtain the same result for $W = Z$, we perform exactly the same steps starting from \Cref{eq:qld-s-good-and-bad} but with $X$ and $Z$ interchanged (so we use \Cref{eq:qld-g-42} instead of \Cref{eq:qld-g-43}). Finally, entirely analogous calculations with the registers swapped yield \Cref{eq:qld-s-point-con-bob}.

        




	
	
\end{proof}

 \subsection{Pulling the $X$ and $Z$ measurements apart}
\label{sec:separating}


Recall the decoding map of the low-degree code defined in \Cref{sec:ld-encoding}. This map takes in a polynomial $g \in \ideg_{d,m}(\F_q)$ and returns a string $\coded(g) \in \F_q^{\nqubits}$ consisting of the evaluations of $g$ on the points in the hypercube $\{0,1\}^m$.
We now use this decoding map to construct a new set of measurement operators out of $\hat{S}$.
First, we introduce some notation. For all $W \in \{X,Z\}$ and all degree $d$ polynomials $g: \F_q^m \to \F_q$ define projectors
\[ \hat{S}^X_g = \sum_{g_\zpt} \hat{S}_{g,g_\zpt} \quad  \text{and}  \quad \hat{S}^Z_g = \sum_{g_\xpt} \hat{S}_{g_\xpt,g}\;,\]
where the measurement $\{\hat{S}_{g_\xpt,g_\zpt} \}$  is given by \Cref{lem:qld-4-7}.

For all $W \in \{X, Z\}$, $\tilde{u} \in \F_q^{\nqubits}$, and $a \in \F_q$, define the projector
\begin{equation}\label{eq:def-tauwu}
 \tau^{W}_{a}(\tilde{u}) = \sum_{h \in \F_q^{\nqubits} \,:\, h \cdot \tilde{u} = a} \tau^W_{h}
\end{equation}
where recall that $\tau^W_h$ is the projector corresponding to measuring $(\C^q)^{\otimes \nqubits}$ in the $W$ basis and obtaining $h \in \F_q^\nqubits$ as the outcome. For all $\tilde{u} \in \F_q^\nqubits$ define the projective measurement $\{ \tilde{M}^{W,\tilde{u}}_a \}_{a \in \F_q}$ by 
	\begin{equation}
	\label{eq:tilde_M}
	\tilde{M}^{W,\tilde{u}}_a =  \sum_{g} \hat{S}^W_g \otimes \tau_{(\coded(g) \cdot \tilde{u}) - a}^W(\tilde{u})\;.
	\end{equation}
	If $\hat{S}^W_g$ is viewed as an operator acting on registers $\alice \alice'$ (resp. $\bob \bob'$) then we view $\tilde{M}^{W,\tilde{u}}_a$ as an operator acting on registers $\alice \alice' \alice''$ (resp. $\bob \bob' \bob''$). (When necessary we will indicate via subscripts which registers the operators are acting on.)
The $\tilde{M}^{W,\tilde{u}}_a$ are projective because the $\hat{S}^W_g$ and $\tau^W$ operators are projective.


	For all $j \in \{1,\ldots,t\}$ define the observable
	\begin{equation}\label{eq:def-tildewj}
		\tilde{W}^j(\tilde{u}) = \sum_a (-1)^{\tr(e_j a)} \, \tilde{M}_a^{W,\tilde{u}} \;,
	\end{equation}
	where $\{e_j\}_{j \in \{1,\ldots,t\}}$ denotes the self-dual basis for $\F_q$ over $\F_2$ specified in \Cref{sec:subfields}. Let $\tau^W$ be the generalized Pauli observable defined in Section~\ref{sec:generalized-pauli}. Observe that
	\begin{align*}
	\Big ( \sum_{g_\xpt,g_\zpt} (-1)^{\tr(e_j (\coded(g_W) \cdot \tilde{u}))} \hat{S}_{g_\xpt,g_\zpt} \Big ) \otimes \tau^W(e_j \tilde{u})
	&= \sum_{g_\xpt,g_\zpt} \sum_h (-1)^{\tr(e_j (\coded(g_W) \cdot \tilde{u}))} (-1)^{\tr(e_j(\tilde{u} \cdot h))} \hat{S}_{g_\xpt,g_\zpt} \otimes \tau^W_h\\
	&= \sum_{g_W} \sum_a  (-1)^{\tr(e_j (\tilde{u} \cdot\coded(g_W) + a))} \hat{S}_{g_W}^W \otimes \Big( \sum_{h:\, h\cdot \tilde{u} = a}\tau^W_h\Big)\\
		&= \sum_{g_W} \sum_a  (-1)^{\tr(e_j (\tilde{u} \cdot\coded(g_W) + a))} \hat{S}_{g_W}^W \otimes \tau^W_a(\tilde{u})\\
		&= 		\tilde{W}^j(\tilde{u}) \;,
	\end{align*}
	where the first equality uses the definition of $\tau^W(e_j \tilde{u})$, the second uses the definition of $\hat{S}^W_{g_W}$ and regroups terms, the third uses the definition~\eqref{eq:def-tauwu} of $\tau^W_a(\tilde{u})$, and the last is by a change of variables. These equations in particular show that the observable $\tilde{W}^j(\tilde{u})$ is Hermitian and squares to the identity, because the measurement $\{ \hat{S}_{g_\xpt,g_\zpt} \}$ is projective. 
	Furthermore, it can be verified by direct calculation that the observables $\tilde{W}^j(\tilde{u})$ satisfy the same relations as the generalized Pauli group. Specifically, for all $j,j' \in \{1,\ldots,t\}$, $\tilde{u}, \tilde{v} \in \F_q^\nqubits$,
	\[
		\tilde{X}^j(\tilde{u}) \tilde{Z}^{j'}(\tilde{v}) = \begin{cases}
		\tilde{Z}^{j'}(\tilde{v}) \tilde{X}^j(\tilde{u})  & \qquad j \neq j' \\
		(-1)^{\tr((e_j \tilde{u}) \cdot (e_j \tilde{v}))} \tilde{Z}^{j}(\tilde{v}) \tilde{X}^j(\tilde{u}) & \qquad j = j'
		\end{cases}~.
	\]
	The next lemma shows that the measurements $\{ \tilde{M}^{W,\tilde{u}}_a \}$  and observables $\tilde{W}^j(\tilde{u})$  are consistent with the original points measurements, and are self-consistent. 
	



\begin{lemma}
  \label{lem:qld-construct-the-paulis}
Let $\delta_S$ as in Lemma~\ref{lem:qld-4-7}.
For all $W \in \{X,Z\}$, all $\tilde{u} \in \F_q^\nqubits$ and $j \in \{1,\ldots,t\}$, the measurements $\{ \tilde{M}^{W,\tilde{u}}_a \}_{a \in \F_q}$ and observables $\tilde{W}^j(\tilde{u})$ satisfy the following properties:
\begin{enumerate}
\item (\textbf{Consistency with points measurements}) For all $W \in \{X,Z\}$, for all $j \in \{1,\ldots,t\}$, and on average over uniformly random $u \in \F_q^m$,
	\[
		\Id_{\alice \alice' \alice''} \otimes \tilde{M}^{W,\ind_m(u)}_a \simeq_{\delta_S} M^{(\Point,W),u}_a \otimes \Id_{\bob \bob' \bob''}
	\]
	and
	\[
		\tilde{M}^{W,\ind_m(u)}_a \otimes \Id_{\bob \bob' \bob''} \simeq_{\delta_S} \Id_{\alice \alice' \alice''} \otimes M^{(\Point,W),u}_a \;.
	\]	
	\item (\textbf{Self-consistency}) On average over uniformly random $\tilde{u} \in \F_q^\nqubits$,
	\[
		\tilde{W}^j(\tilde{u}) \otimes \Id_{\bob \bob' \bob''} \approx_{\delta_S} \Id_{\alice \alice' \alice''} \otimes \tilde{W}^j(\tilde{u})\,.
	\]
\end{enumerate}
\end{lemma}

\begin{proof}
We show the first item of the lemma for the case when $M^{(\Point,W),u}_a$ acts on register $\alice$ and $\tilde{M}^{W,\ind_m(u)}_a$ acts on registers $\bob \bob' \bob''$ (the argument for when the $\alice$ and $\bob$ registers are interchanged proceeds analogously):
	\begin{align*}
		&\E_{u \in \F_q^m} \sum_a \bra{\hat{\psi}} M^{(\Point,W),u}_a \otimes \tilde{M}^{W,\ind_m(u)}_a \ket{\hat{\psi}} \\
		&= \E_{u \in \F_q^m} \sum_{g,a} \bra{\hat{\psi}} (M^{(\Point,W),u}_a)_{\alice} \otimes (\hat{S}^W_g)_{\bob \bob'} \otimes (\tau^W_{\coded(g) \cdot \ind_m(u) -a}(\ind_m(u)))_{\bob''} \ket{\hat{\psi}} \\
		&= \E_{u \in \F_q^m} \sum_{g} \bra{\hat{\psi}} (\hat{M}^{(\Point,W),u}_{\coded(g) \cdot \ind_m(u)})_{\alice \bob''} \otimes (\hat{S}^W_g)_{\bob \bob'} \ket{\hat{\psi}} \\
		&= \E_{u \in \F_q^m} \sum_{g} \bra{\hat{\psi}} (\hat{M}^{(\Point,W),u}_{g(u)})_{\alice \bob''} \otimes (\hat{S}^W_g)_{\bob \bob'} \ket{\hat{\psi}} \\
		&\geq 1 - \delta_S\;.
	\end{align*}
	The second line follows from expanding the definition of $\tilde{M}^{W,\ind_m(u)}_a$ using~\eqref{eq:tilde_M}, the third line follows from the definition of $\hat{M}^{(\Point,W),u}_a$ given in Section~\ref{sec:expanding}, the fourth line follows from the fact that $\coded(g) \cdot \ind_m(u) = g(u)$, and the fifth line follows from \Cref{lem:qld-4-7}. This implies the first item of the Lemma.


To show the second item we first argue that, on average over $\tilde{u}$, 
	\begin{equation}
		\label{eq:qld-pulling-cons}
		\Big ( \tilde{M}^{W,\tilde{u}}_a \Big)_{\alice \alice' \alice''} \approx_{\delta_S} \Big ( \tilde{M}^{W,\tilde{u}}_a \Big)_{\bob \bob' \bob''}.
	\end{equation}
	This follows because on average over $u \in \F_q^m$ and $\tilde{u} \in \F_q^\nqubits$ we have
	\begin{align}
		&\Big ( \tilde{M}^{W,\tilde{u}}_a \Big)_{\alice \alice' \alice''} \label{eq:qld-pulling-0} \\
		\approx_{\delta_S}& \Big ( \tilde{M}^{W,\tilde{u}}_a \Big)_{\alice \alice'\alice''} \cdot \Big ( \sum_g (\hat{S}^W_g)_{\alice \alice'} \otimes  (\hat{M}^{(\Point,W),u}_{g(u)})_{\bob \alice''} \Big) \label{eq:qld-pulling-1} \\
		=_{\phantom{\delta_S}}& \sum_g \Big ( (\hat{S}^W_g)_{\alice \alice'} \otimes (\tau^W_{(\coded(g) \cdot \tilde{u}) - a}(\tilde{u}))_{\alice''} \Big) \cdot \Big(\sum_{r,s: r + s = g(u)} (M^{(\Point,W),u}_r)_{\bob} \otimes (\tau_s^W(\ind_m(u)))_{\alice''} \Big) \label{eq:qld-pulling-2} \\
		=_{\phantom{\delta_S}}& \sum_{\substack{g,h: \\ (\coded(g)-h)\cdot \tilde{u} = a}}  (\hat{S}^W_g)_{\alice \alice'} \otimes (\tau^W_h)_{\alice''} \otimes (M^{(\Point,W),u}_{(g-g_h)(u)})_{\bob}~. \label{eq:qld-pulling-2b}
	\end{align}
	The first inequality follows from right-multiplying~\eqref{eq:qld-pulling-0} by the approximation~\eqref{eq:qld-sg-cons} in \Cref{lem:qld-constructing-the-paulis-helper} (to be proved below) and using Fact~\ref{fact:add-a-proj}. Line~\eqref{eq:qld-pulling-2} follows from expanding the definition of $\tilde{M}^{W,\tilde{u}}_a$.  To obtain Line~\eqref{eq:qld-pulling-2b}, we expanded the definitions of $\tau^W_{(\coded(g) \cdot \tilde{u}) - a}(\tilde{u})$ and $\tau^W_s(\ind_m(u))$ using~\cref{eq:def-tauwu}, and the fact that $h \cdot \ind_m(u) = g_h(u)$ (due to \cref{eq:low-degree-encoding-definition}). Thus for all $g \in \ideg_{d,m}(\F_q)$ and $s \in \F_q$ we have
	\[
		\tau^W_{(\coded(g) \cdot \tilde{u}) - a}(\tilde{u}) \cdot \tau^W_s(\ind_m(u)) = \sum_{\substack{h: \\ (\coded(g) - h) \cdot \tilde{u} = a \\ g_h(u) = s}} \tau^W_h~.
	\]


	Using the self-consistency of $\{M^{(\Point,W),u}_a\}$ shown in Lemma~\ref{lem:qld-win-implications}, we then obtain
	\begin{align}
		\text{\eqref{eq:qld-pulling-2b}} &\approx_{\eps} \sum_{\substack{g,h: \\ (\coded(g)-h)\cdot \tilde{u} = a}} (\hat{S}^W_g)_{\alice \alice'} \cdot \Big((M^{(\Point,W),u}_{(g-g_h)(u)})_{\alice} \otimes (\tau_h^W)_{\alice''} \Big) \cdot \Big ( M^{(\Point,W),u}_{(g-g_h)(u)} \Big)_{\bob} \label{eq:qld-pulling-3} \\
		&\approx_{0}\sum_{\substack{g,h: \\ (\coded(g)-h)\cdot \tilde{u} = a}} (\hat{S}^W_g)_{\alice \alice'} \cdot \Big((M^{(\Point,W),u}_{(g-g_h)(u)})_{\alice} \otimes (\tau_h^W)_{\alice'} \otimes (\tau_h^W)_{\alice''} \Big) \cdot \Big ( M^{(\Point,W),u}_{(g-g_h)(u)} \Big)_{\bob} \label{eq:qld-pulling-3a} \\
		&= \sum_{\substack{g,h: \\ (\coded(g)-h)\cdot \tilde{u} = a}} (\hat{S}^W_g)_{\alice \alice'} \cdot \Big((\hat{M}^{(\Point,W),u}_{g(u)})_{\alice \alice'} \otimes (\tau_h^W)_{\alice''} \Big) \cdot \Big ( M^{(\Point,W),u}_{(g-g_h)(u)} \Big)_{\bob} \;. \label{eq:qld-pulling-3b}
	\end{align}
	The ``$\approx_0$'' in Line~\eqref{eq:qld-pulling-3a} indicates that the left hand side is equal to the right hand side when applied to the state $\ket{\hat{\psi}}$. Line~\eqref{eq:qld-pulling-3a} follows from the self-consistency of the operators $\{\tau^W_h\}$, and Line~\eqref{eq:qld-pulling-3b} follows from the definition of $\{ \hat{M}^{(\Point,W),u}_a\}$. Using the fact that $\ket{\hat{\psi}} = \sum_{h'} (\tau^W_{h'})_{\bob'} \otimes (\tau^W_{h'})_{\bob''} \ket{\hat{\psi}}$, we get that 
	\begin{align}
	\text{\eqref{eq:qld-pulling-3b}} & \approx_0 \sum_{\substack{g,h: \\ (\coded(g)-h)\cdot \tilde{u} = a}} \Big((\hat{S}^W_g \cdot \hat{M}^{(\Point,W),u}_{g(u)})_{\alice \alice'} \otimes (\tau_h^W)_{\alice''} \Big) \cdot \Big( M^{(\Point,W),u}_{(g-g_h)(u)} \otimes \sum_{h'} \tau_{h'}^W \otimes \tau_{h'}^W\Big)_{\bob \bob' \bob''} \label{eq:qld-pulling-4} \\
		&= \sum_{\substack{g,h,h': \\ (\coded(g)-h)\cdot \tilde{u} = a}} \Big((\hat{S}^W_g \cdot \hat{M}^{(\Point,W),u}_{g(u)})_{\alice \alice'} \otimes (\tau_h^W)_{\alice''} \Big) \cdot \Big( (\hat{M}^{(\Point,W),u}_{(g-g_h+g_{h'})(u)})_{\bob \bob'} \otimes (\tau_{h'}^W)_{\bob''} \Big) \;. \label{eq:qld-pulling-5}
\end{align}
	Line~\eqref{eq:qld-pulling-5} follows from~\eqref{eq:qld-pulling-4} by expanding the definition of $\{ \hat{M}^{(\Point,W),u}_a\}$. Next, we argue that, on average over $u \in \F_q^m$,
\begin{align}
		\text{\eqref{eq:qld-pulling-5}} &\approx_{\delta_S} \sum_{\substack{g,h,h': \\ (\coded(g)-h)\cdot \tilde{u} = a}} \Big((\hat{S}^W_g)_{\alice \alice'} \otimes (\tau_h^W)_{\alice''} \Big) \otimes \Big( (\hat{M}^{(\Point,W),u}_{(g-g_h+g_{h'})(u)})_{\bob \bob'} \otimes (\tau_{h'}^W)_{\bob''} \Big)~. \label{eq:qld-pulling-7}
		\end{align}
		We show this by bounding the magnitude of the difference. Using the fact that $\hat{S}^W_g$ and $\tau^W_h$ are projective, we have:
		\begin{align}
		&\E_u \sum_a \Big \|\sum_{\substack{g,h,h': \\ (\coded(g)-h)\cdot \tilde{u} = a}} (\hat{S}^W_g \cdot (\Id - \hat{M}^{(\Point,W),u}_{g(u)}))_{\alice \alice'} \otimes (\tau_h^W)_{\alice''} \otimes (\hat{M}^{(\Point,W),u}_{(g-g_h+g_{h'})(u)})_{\bob \bob'} \otimes (\tau_{h'}^W)_{\bob''} \, \ket{\hat{\psi}} \Big \|^2 \notag \\
      =& \E_u \sum_{\substack{g,h,h'}} \bra{\hat{\psi}} \Big(((\Id - \hat{M}^{(\Point,W),u}_{g(u)}) \cdot \hat{S}^W_g \cdot (\Id - \hat{M}^{(\Point,W),u}_{g(u)}))_{\alice \alice'} \otimes (\tau_h^W)_{\alice''} \Big) \notag \\
      & \hspace{20em} \cdot \Big( (\hat{M}^{(\Point,W),u}_{(g-g_h+g_{h'})(u)})^2_{\bob \bob'} \otimes (\tau_{h'}^W)_{\bob''} \Big) \ket{\hat{\psi}} \notag \\
		\leq& \E_u \sum_{g} \bra{\hat{\psi}} ((\Id - \hat{M}^{(\Point,W),u}_{g(u)}) \cdot \hat{S}^W_g \cdot (\Id - \hat{M}^{(\Point,W),u}_{g(u)}))_{\alice \alice'} \ket{\hat{\psi}}\label{eq:qld-pulling-8} \\
		\leq& \delta_S\;. \label{eq:qld-pulling-9}
		\end{align}
		The inequality in Line~\eqref{eq:qld-pulling-8} follows from the fact that for all $g,h$, the sum $\sum_{h'} (\hat{M}^{(\Point,W),u}_{(g-g_h+g_{h'})(u)})^2_{\bob \bob'} \otimes (\tau_{h'}^W)_{\bob''}$ is at most $\Id$. The inequality in Line~\eqref{eq:qld-pulling-9} follows from the second item of \Cref{lem:qld-constructing-the-paulis-helper} (to be proved below). 
		Next, we show that on average over $u \in \F_q^m$,
		\begin{align}
		\text{\eqref{eq:qld-pulling-7}} &=\sum_{\substack{g,h,g',h': \\ (\coded(g)-h)\cdot \tilde{u} = a}} \Big((\hat{S}^W_g)_{\alice \alice'} \otimes (\tau_h^W)_{\alice''} \Big) \cdot \Big( (\hat{S}^W_{g'} \cdot \hat{M}^{(\Point,W),u}_{(g-g_h+g_{h'})(u)})_{\bob \bob'} \otimes (\tau_{h'}^W)_{\bob''} \Big) \label{eq:qld-pulling-9a} \\
		&\approx_{O(\delta_S + \sqrt{\eps})} \sum_{\substack{g,h,g',h': \\ (g' - g_{h'})(u) = (g -g_h)(u) \\ (\coded(g)-h)\cdot \tilde{u} = a}} \Big((\hat{S}^W_g)_{\alice \alice'} \otimes (\tau_h^W)_{\alice''} \Big) \cdot \Big( (\hat{S}^W_{g'} \cdot \hat{M}^{(\Point,W),u}_{g'(u)})_{\bob \bob'} \otimes (\tau_{h'}^W)_{\bob''} \Big) \label{eq:qld-pulling-10} \\
		&\approx_{\delta_S} \sum_{\substack{g,h,g',h': \\ (g' - g_{h'})(u) = (g -g_h)(u) \\ (\coded(g)-h)\cdot \tilde{u} = a}} \Big((\hat{S}^W_g)_{\alice \alice'} \otimes (\tau_h^W)_{\alice''} \Big) \cdot \Big( (\hat{S}^W_{g'})_{\bob \bob'} \otimes (\tau_{h'}^W)_{\bob''} \Big) \label{eq:qld-pulling-11} \\
		&\approx_{md/q} \sum_{\substack{g,h,g',h': \\ g' - g_{h'} = g -g_h \\ (\coded(g)-h)\cdot \tilde{u} = a}} \Big((\hat{S}^W_g)_{\alice \alice'} \otimes (\tau_h^W)_{\alice''} \Big) \cdot \Big( (\hat{S}^W_{g'})_{\bob \bob'} \otimes (\tau_{h'}^W)_{\bob''} \Big)~. \label{eq:qld-pulling-12}
		\end{align}
		Before we justify the sequence of steps to go from Line~\eqref{eq:qld-pulling-7} to Line~\eqref{eq:qld-pulling-12}, we first argue how this establishes Line~\eqref{eq:qld-pulling-cons}, and hence the second item of the Lemma statement.
		
		Absorbing the error terms $O(md/q)$ and $O(\sqrt{\eps})$ into $\delta_S$, we have shown that 
		\[
		\Big ( \tilde{M}^{W,\tilde{u}}_a \Big)_{\alice \alice' \alice''} \approx_{\delta_S} \sum_{\substack{g,h,g',h': \\ g' - g_{h'} = g -g_h \\ (\coded(g)-h)\cdot \tilde{u} = a}} \Big((\hat{S}^W_g)_{\alice \alice'} \otimes (\tau_h^W)_{\alice''} \Big) \cdot \Big( (\hat{S}^W_{g'})_{\bob \bob'} \otimes (\tau_{h'}^W)_{\bob''} \Big).
		\]
		Note that the conditions $(\coded(g)-h)\cdot \tilde{u} = a$ and $g - g_h = g' - g_{h'}$ are equivalent to the conditions $(\coded(g') - h') \cdot \tilde{u} = a$ and $g - g_h = g' - g_{h'}$. Thus the expression in~\eqref{eq:qld-pulling-12} is symmetric between $g,h$ and $g',h'$, and so an analogous derivation shows that $\Big ( \tilde{M}^{W,\tilde{u}}_a \Big)_{\bob \bob' \bob''}$ is $\delta_S$-close to~\eqref{eq:qld-pulling-12}, and thus $\delta_S$-close to $\Big ( \tilde{M}^{W,\tilde{u}}_a \Big)_{\alice \alice' \alice''}$. This shows~\eqref{eq:qld-pulling-cons}. Since $\{ \tilde{M}^{W,\tilde{u}}_a \}$ is projective, by using Item 2 of \Cref{fact:agreement}, followed by \Cref{fact:data-processing}, and then followed by Item 1 of \Cref{fact:agreement}, we get for all $j \in \{1,\ldots,t\}$, on average over $\tilde{u} \in \F_q^\nqubits$,
\[
	(\tilde{M}^{W,\tilde{u}}_{[\tr(e_j \cdot) = b]})_{\alice \alice' \alice''} \approx_{\delta_S} ( \tilde{M}^{W,\tilde{u}}_{[\tr(e_j \cdot) = b]})_{\bob \bob' \bob''}
\]
where the answer summation is over $b \in \F_2$. Finally, the second item of the Lemma follows from applying~\Cref{lem:povm-to-obs}.

		We now return to justifying the sequence of steps between Lines~\eqref{eq:qld-pulling-7} and~\eqref{eq:qld-pulling-12}. The equality in Line~\eqref{eq:qld-pulling-9a} follows by left-multiplying~\eqref{eq:qld-pulling-7} by $I_{\bob \bob'} = \sum_{g'} (\hat{S}^W_{g'})_{\bob \bob'}$. 
		
		\paragraph{The approximation in Line~\eqref{eq:qld-pulling-10}.} The approximation in Line~\eqref{eq:qld-pulling-10} follows by bounding the magnitude of the difference:	
\begin{align}
		&\E_u \sum_a \Big \| \sum_{\substack{g,h,g',h': \\ (g' - g_{h'})(u) \neq (g -g_h)(u) \\ (\coded(g)-h)\cdot \tilde{u} = a}} \Big((\hat{S}^W_g)_{\alice \alice'} \otimes (\tau_h^W)_{\alice''} \Big) \cdot \Big( (\hat{S}^W_{g'} \cdot \hat{M}^{(\Point,W),u}_{(g-g_h+g_{h'})(u)})_{\bob \bob'} \otimes (\tau_{h'}^W)_{\bob''} \Big) \ket{\hat{\psi}} \Big \|^2 \notag \\
&= \E_u \sum_{\substack{g,h,g',h': \\ (g' - g_{h'})(u) \neq (g -g_h)(u)}} \bra{\hat{\psi}} \Big(\hat{S}^W_g)_{\alice \alice'} \otimes (\tau_h^W)_{\alice''} \Big) \cdot \Big( (\hat{M}^{(\Point,W),u}_{(g-g_h+g_{h'})(u)} \cdot \hat{S}^W_{g'} \cdot \notag\\
      &\hspace{20em}\hat{M}^{(\Point,W),u}_{(g-g_h+g_{h'})(u)})_{\bob \bob'} \otimes (\tau_{h'}^W)_{\bob''} \Big) \ket{\hat{\psi}} \notag \\
		&\leq \E_u \sum_{\substack{g,h,g',h',a': \\ a' \neq g'(u)}} \bra{\hat{\psi}} \Big(\hat{S}^W_g)_{\alice \alice'} \otimes (\tau_h^W)_{\alice''} \Big) \cdot \Big( (\hat{M}^{(\Point,W),u}_{a'} \cdot \hat{S}^W_{g'} \cdot \hat{M}^{(\Point,W),u}_{a'})_{\bob \bob'} \otimes (\tau_{h'}^W)_{\bob''} \Big) \ket{\hat{\psi}} \notag \\
		&= \E_u \sum_{g,a: a \neq g(u)} \bra{\hat{\psi}} (\hat{M}^{(\Point,W),u}_{a} \cdot \hat{S}^W_{g} \cdot \hat{M}^{(\Point,W),u}_{a})_{\bob \bob'} \ket{\hat{\psi}} \label{eq:qld-pulling-13a} \\
		&\approx_{O(\sqrt{\eps})} \E_u \sum_{g} \bra{\hat{\psi}} (\Id - \hat{M}^{(\Point,W),u}_{g(u)})_{\alice \bob''} \otimes ( \hat{S}^W_{g})_{\bob \bob'} \ket{\hat{\psi}}  \label{eq:qld-pulling-13} \\
		&\leq \delta_S \;. \label{eq:qld-pulling-14}
\end{align}
The approximation in Line~\eqref{eq:qld-pulling-13} follows from the following calculation:
\begin{align}
	&\Big | \text{\eqref{eq:qld-pulling-13}} - \text{\eqref{eq:qld-pulling-13a}} \Big | \notag \\
	&\leq \Big | \E_u \sum_{g,a: a \neq g(u)} \bra{\hat{\psi}} (\hat{M}^{(\Point,W),u}_{a} \cdot \hat{S}^W_{g})_{\bob \bob'} \cdot ((\hat{M}^{(\Point,W),u}_{a})_{\alice \bob''} - (\hat{M}^{(\Point,W),u}_{a})_{\bob \bob'}) \ket{\hat{\psi}} \Big | \label{eq:qld-pulling-15} \\
	& \qquad + \Big | \E_u \sum_{g,a: a \neq g(u)} \bra{\hat{\psi}} ((\hat{M}^{(\Point,W),u}_{a})_{\alice \bob''} - (\hat{M}^{(\Point,W),u}_{a})_{\bob \bob'}) \cdot (\hat{S}^W_{g})_{\bob \bob'} \cdot (\hat{M}^{(\Point,W),u}_{a})_{\alice \bob''} \ket{\hat{\psi}} \Big | \label{eq:qld-pulling-15a}
\end{align}
Using Cauchy-Schwarz, we can bound the first term on the right hand side of the inequality by
\begin{align}
	&\Biggl( \E_u \sum_{g,a: a \neq g(u)} \bra{\hat{\psi}} (\hat{M}^{(\Point,W),u}_{a} \cdot \hat{S}^W_{g} \cdot \hat{M}^{(\Point,W),u}_{a})_{\bob \bob'} \ket{\hat{\psi}}  \Biggr)^{1/2} \cdot \notag \\
	&\Biggl( \E_u \sum_{g,a: a \neq g(u)} \bra{\hat{\psi}} ((\hat{M}^{(\Point,W),u}_{a})_{\alice \bob''} - (\hat{M}^{(\Point,W),u}_{a})_{\bob \bob'}) \cdot (\hat{S}^W_{g})_{\bob \bob'} \cdot \notag\\
  & \hspace{12em}((\hat{M}^{(\Point,W),u}_{a})_{\alice \bob''} - (\hat{M}^{(\Point,W),u}_{a})_{\bob \bob'})  \ket{\hat{\psi}}  \Biggr)^{1/2} \notag \\
	\leq & \sqrt{1} \cdot \sqrt{\eps} \notag
\end{align}
The inequality follows from the self-consistency of $\hat{M}^{(\Point,W),u}_{a}$, as established in \Cref{lem:qld-comm-cons}. Similarly, we can bound the second term in Line~\eqref{eq:qld-pulling-15a} by $\sqrt{\eps}$.
This establishes the approximation in Line~\eqref{eq:qld-pulling-13}. 

The inequality in Line~\eqref{eq:qld-pulling-14} follows from \Cref{eq:qld-s-point-con-bob} of \Cref{lem:qld-4-7}. This concludes the derivation of the approximation in Line~\eqref{eq:qld-pulling-10}.

\paragraph{The approximation in Line~\eqref{eq:qld-pulling-11}.} We establish this approximation by bounding the magnitude of the difference:
\begin{align*}
	&\E_u \sum_a \Big \| \sum_{\substack{g,h,g',h': \\ (g' - g_{h'})(u) = (g -g_h)(u) \\ (\coded(g)-h)\cdot \tilde{u} = a}} \Big((\hat{S}^W_g)_{\alice \alice'} \otimes (\tau_h^W)_{\alice''} \Big) \cdot \Big( (\hat{S}^W_{g'} \cdot (\Id - \hat{M}^{(\Point,W),u}_{g'(u)}))_{\bob \bob'} \otimes (\tau_{h'}^W)_{\bob''} \Big) \ket{\hat{\psi}} \Big \|^2 \\
	=& \E_u \sum_{\substack{g,h,g',h': \\ (g' - g_{h'})(u) \\ \quad = (g -g_h)(u)}} \bra{\hat{\psi}} (\hat{S}^W_g)_{\alice \alice'} \otimes (\tau_h^W)_{\alice''} \otimes ((\Id - \hat{M}^{(\Point,W),u}_{g'(u)}) \cdot \hat{S}^W_{g'} \cdot \notag \\
  &\hspace{20em}(\Id - \hat{M}^{(\Point,W),u}_{g'(u)}))_{\bob \bob'} \otimes (\tau_{h'}^W)_{\bob''} \ket{\hat{\psi}}  \\
	\leq& \E_u \sum_{g'} \bra{\hat{\psi}} ((\Id - \hat{M}^{(\Point,W),u}_{g'(u)}) \cdot \hat{S}^W_{g'} \cdot (\Id - \hat{M}^{(\Point,W),u}_{g'(u)}))_{\bob \bob'} \ket{\hat{\psi}} \\
	\leq& \delta_S \;.
\end{align*}
The last inequality follows from the second item of \Cref{lem:qld-constructing-the-paulis-helper} (to be proved below). 

\paragraph{The approximation in Line~\eqref{eq:qld-pulling-12}.} We establish this approximation by bounding the magnitude of the difference:
\begin{align*}
	&\E_u \sum_a \Big \| \sum_{\substack{g,h,g',h': \\ g' - g_{h'} \neq g -g_h \\ (g' - g_{h'})(u) = (g -g_h)(u) \\ (\coded(g)-h)\cdot \tilde{u} = a}} \Big((\hat{S}^W_g)_{\alice \alice'} \otimes (\tau_h^W)_{\alice''} \Big) \cdot \Big( (\hat{S}^W_{g'})_{\bob \bob'} \otimes (\tau_{h'}^W)_{\bob''} \Big) \ket{\hat{\psi}} \Big \|^2 \\
	&= \sum_{\substack{g,h,g',h': \\ g' - g_{h'} \neq g -g_h}} \bra{\hat{\psi}} \Big((\hat{S}^W_g)_{\alice \alice'} \otimes (\tau_h^W)_{\alice''} \Big) \cdot \Big( (\hat{S}^W_{g'})_{\bob \bob'} \otimes (\tau_{h'}^W)_{\bob''} \Big) \ket{\hat{\psi}} \cdot \E_u \mathbf{1}[(g' - g_{h'})(u) = (g -g_h)(u)] \\
	&\leq md/q.
\end{align*}
The last inequality follows from the Schwartz-Zippel lemma: two distinct polynomials of total degree at most $md$ cannot agree on more than $md/q$ fraction of points. 

		

\end{proof}


We now give a proof of the following Lemma, which was used in several derivations in \Cref{lem:qld-construct-the-paulis}. 
\begin{lemma}
\label{lem:qld-constructing-the-paulis-helper}
	For $W \in \{X,Z\}$, on average over $u \in \F_q^m$, we have that
	\begin{gather}
		\label{eq:qld-sg-cons}
		\sum_g (\hat{S}^W_g)_{\alice \alice'} \otimes (\hat{M}^{(\Point,W),u}_{g(u)})_{\bob \alice''} \approx_{\delta_S} \Id
	\end{gather}
	and
	\begin{gather}
		\label{eq:qld-sg-cons2}
		(\hat{S}^W_g \cdot (\Id - \hat{M}^{(\Point, W), u}_{g(u)}))_{\alice \alice'} \approx_{\delta_S} 0
	\end{gather}
	where the answer summation is over $g \in \ideg_{d,m}(\F_q)$.
      Furthermore, all symmetric equivalents of these approximations
      also hold.
\end{lemma}
\begin{proof}
We first prove \Cref{eq:qld-sg-cons}. From \Cref{eq:qld-s-point-con-alice} of \Cref{lem:qld-4-7} we have that on average over $u \in \F_q^m$,
\begin{equation}
\label{eq:qld-constructing-the-paulis-helper0}
	  \bra{\hat{\psi}} \sum_g (\hat{S}^W_g)_{\alice \alice'} \otimes (\hat{M}^{(\Point, W), u}_{g(u)})_{\bob \alice''} \ket{\hat{\psi}} \geq 1 - \delta_S
\end{equation}
which is equivalent to $\sum_g (\hat{S}^W_g)_{\alice \alice'} \otimes (\hat{M}^{(\Point, W), u}_{g(u)})_{\bob \alice''}  \simeq_{\delta_S} \Id$. By \Cref{fact:agreement}, we have that
\[
\sum_g (\hat{S}^W_g)_{\alice \alice'} \otimes (\hat{M}^{(\Point, W), u}_{g(u)})_{\bob \alice''} \approx_{\delta_S} \Id.
\]

We now prove \Cref{eq:qld-sg-cons2}. \Cref{eq:qld-constructing-the-paulis-helper0} and \Cref{fact:agreement} shows that
\[ (\hat{S}^W_{[g(u)=a]})_{\alice \alice'} \approx_{\delta_S} (\hat{M}^{(\Point, W), u}_{a})_{\bob \alice''}\]
where the answer summation is over $a \in \F_q$. Left-multiplying this expression by $(\hat{S}^W_{[g(u)=a]})_{\alice \alice'}$ and using \Cref{fact:add-a-proj} we get
\begin{align*}
(\hat{S}^W_{[g(u)=a]})_{\alice \alice'} &\approx_{\delta_S} (\hat{S}^W_{[g(u)=a]})_{\alice \alice'} \otimes (\hat{M}^{(\Point, W), u}_{a})_{\bob \alice''} \\
&\approx_\eps (\hat{S}^W_{[g(u)=a]} \cdot \hat{M}^{(\Point, W), u}_{a})_{\alice \alice'}
\end{align*}
where the second approximation follows from the following calculation:
\begin{align}
	& \E_u \sum_a \left \| (\hat{S}^W_{[g(u)=a]})_{\alice \alice'} \cdot \Big ( (\hat{M}^{(\Point, W), u}_{a})_{\alice \alice'} - (\hat{M}^{(\Point, W), u}_{a})_{\bob \alice''} \Big ) \ket{\hat{\psi}} \right \|^2  \\
	\leq & \E_u \sum_a \bra{\hat{\psi}} \Big ( (\hat{M}^{(\Point, W), u}_{a})_{\alice \alice'} - (\hat{M}^{(\Point, W), u}_{a})_{\bob \alice''} \Big )^2 \ket{\hat{\psi}} \\
	= & \E_u \sum_a \left \| (\hat{M}^{(\Point, W), u}_{a})_{\alice \alice'} - (\hat{M}^{(\Point, W), u}_{a})_{\bob \alice''} \ket{\hat{\psi}} \right \|^2 \\
	\leq & \eps
\end{align}
The first inequality follows from the fact that $\hat{S}^W_{[g(u)=a]} \leq \Id$ for all $a$, and the second inequality follows from the self-consistency of $\hat{M}^{(\Point, W), u}_{a}$ (established in \Cref{lem:qld-comm-cons}). Thus
\[
	(\hat{S}^W_{[g(u)=a]})_{\alice \alice'} \approx_{O(\delta_S + \eps)} (\hat{S}^W_{[g(u)=a]} \cdot \hat{M}^{(\Point, W), u}_{a})_{\alice \alice'}~.
\]
On the other hand, notice that
\begin{align*}
	&\E_u \sum_g \left \| (\hat{S}^W_{g} \cdot (\Id - \hat{M}^{(\Point, W), u}_{g(u)}))_{\alice \alice'} \ket{\hat{\psi}} \right \|^2 \\
	=& \E_u \sum_g \bra{\hat{\psi}} ((\Id - \hat{M}^{(\Point, W), u}_{g(u)}) \cdot \hat{S}^W_{g} \cdot (\Id - \hat{M}^{(\Point, W), u}_{g(u)}))_{\alice \alice'} \ket{\hat{\psi}} \\
	=& \E_u \sum_a \bra{\hat{\psi}} ((\Id - \hat{M}^{(\Point, W), u}_{a}) \cdot \hat{S}^W_{[g(u)=a]} \cdot (\Id - \hat{M}^{(\Point, W), u}_{a}))_{\alice \alice'} \ket{\hat{\psi}} \\
	\leq& O(\delta_S + \eps).
\end{align*}
\Cref{eq:qld-sg-cons2} is thus obtained by absorbing $\eps$ into the function $\delta_S$.  
\end{proof}
	

 The next lemma shows that we can construct local unitaries $V_A,V_B$ from the measurements $\hat{S}$  such that the shared state $\ket{\hat{\psi}}$ is close to $\ket{\aux} \otimes \ket{\EPR_q}^{\otimes \nqubits}$ for some state $\ket{\aux}$, and the observables $\tilde{W}^j(\tilde{u})$ are mapped to Pauli observables $\tau^W(e_j \tilde{u})$ acting on $\ket{\EPR_q}^{\otimes \nqubits}$. Using the consistency between the $\{\tilde{M}^{W,\tilde{u}}_a\}_{a \in \F_q}$ measurements and the points measurements $\{M^{(\Point,W),u}_a\}_{a \in \F_q}$ established in \Cref{lem:qld-construct-the-paulis}, and the consistency between the points measurements and the ``total Pauli measurements'' $\{M^{(\Pauli,W)}_h\}_{h \in \F_q^\nqubits}$ established in \Cref{lem:qld-win-implications}, we deduce that the total Pauli measurements must be close (up to the local unitaries $V_A,V_B$) to the Pauli projectors $\{\tau^W_h\}_{h \in \F_q^\nqubits}$ acting on $\ket{\EPR_q}^{\otimes \nqubits}$.


\begin{lemma}
\label{lem:qld-unitary}
There exists a function $\delta_\qld(\eps,m,d,q) = O(\delta_S^{1/4} + md/q)$ (where $\delta_S$ is defined in \Cref{lem:qld-4-7}) and unitaries $V_\alice$ acting on registers $\alice \alice' \alice''$ and $V_\bob$ acting on $\bob \bob' \bob''$ such that 
\begin{enumerate}
	\item There exists a state $\ket{\aux}$ on registers $\alice \alice' \bob \bob'$ such that
	\[
		\left \| V_\alice \otimes V_\bob \ket{\hat{\psi}}_{\alice \alice' \alice'' \bob \bob' \bob''} - \ket{\aux}_{\alice \alice' \bob \bob'} \otimes \ket{\EPR_q}^{\otimes \nqubits}_{\alice'' \bob''} \right \|^2 \leq \delta_\qld.
	\]
	
	\item For all $W \in \{X,Z\}, h \in \F_q^\nqubits$
	\begin{align*}
		V_\alice \, \Big(M^{(\Pauli,W)}_h \Big)_{\alice} \, V_\alice^\dagger &\approx_{\delta_\qld} \Id_{\alice \alice'} \otimes (\tau^W_h)_{\alice''} \\
		 V_\bob \, \Big(M^{(\Pauli,W)}_h  \Big)_{\bob}\, V_\bob^\dagger &\approx_{\delta_\qld} \Id_{\bob \bob'} \otimes (\tau^W_h)_{\bob''}
	\end{align*}
	where the $\approx$ statements hold with respect to the state $\ket{\aux}_{\alice \alice' \bob \bob'} \otimes \ket{\EPR_q}^{\otimes \nqubits}_{\alice'' \bob''}$. 
\end{enumerate}
\end{lemma}

\begin{proof}

Define the linear map
\[
	V_\alice = \sum_{g_X,g_Z} (\hat{S}_{g_X,g_Z})_{\alice \alice'} \otimes (\tau^X(\coded(g_Z)) \tau^Z(\coded(g_X)))_{\alice ''}\;,
\]
where the $\tau^X(\cdot),\tau^Z(\cdot)$ act on $(\C^q)^{\otimes \nqubits}$. (Also note that the argument of $\tau^X$ is $\coded(g_Z)$ and the argument of $\tau^Z$ is $\coded(g_X)$.) We verify the unitarity of $V_\alice$:
\begin{align*}
	V_\alice V_\alice^\dagger &= \sum_{g_X,g_Z} (\hat{S}_{g_X,g_Z})_{\alice \alice'} \otimes (\tau^X(\coded(g_Z)) \tau^Z(\coded(g_X)) \tau^Z(\coded(g_X)) \tau^X(\coded(g_Z)) )_{\alice ''} \\
	&= \sum_{g_X,g_Z} (\hat{S}_{g_X,g_Z})_{\alice \alice'} \otimes \Id_{\alice ''} \\
	&= \Id_{\alice \alice' \alice''}
\end{align*}
where in the first line we used that the measurement $\{ \hat{S}_{g_\xpt,g_\zpt} \}$  is projective and in the second line we used that the $\tau^Z$ and $\tau^X$ observables are self-inverse. The unitary $V_\bob$ is defined analogously.

Let $\{e_j\}_{j \in \{1,\ldots,t\}}$ denote the self-dual basis for $\F_q$ over $\F_2$ specified in \Cref{sec:subfields}, and let $\tilde{W}^j$ be the observables introduced in~\eqref{eq:def-tildewj}. Conjugating the observable $\tilde{W}^j$ acting on the registers $\alice \alice' \alice''$ by the unitary $V_\alice$ we get
\begin{align*}
	&V_\alice \, \tilde{W}^j(\tilde{u}) \, V_\alice^\dagger \\
  =& \sum_{g_X, g_Z} (-1)^{\tr( e_j(\coded(g_W) \cdot \tilde{u}))} \Big( \hat{S}_{g_X,g_Z}  \Big)_{\alice \alice'} \otimes \Big ( \tau^X(\coded(g_Z)) \tau^Z(\coded(g_X)) \tau^W(e_j \tilde{u}) \\
  &\hspace{24em}\tau^Z(\coded(g_X))^\dagger \tau^X(\coded(g_Z))^\dagger \Big)_{\alice ''} \\
	 =& \sum_{g_X, g_Z} \Big( \hat{S}_{g_X,g_Z}  \Big)_{\alice \alice'} \otimes \Big ( \tau^W(e_j \tilde{u}) \Big)_{\alice ''} \\
	 =& \Id_{\alice \alice'} \otimes \Big( \tau^W(e_j \tilde{u}) \Big)_{\alice ''}\;,
\end{align*}
where in the third line we used the (anti-)commutation relation
\[
	\tau^X(\coded(g_Z)) \tau^Z(\coded(g_X)) \tau^W(e_j \tilde{u}) = (-1)^{\tr( e_j(\coded(g_W) \cdot \tilde{u}))} \tau^W(e_j \tilde{u}) \tau^X(\coded(g_Z)) \tau^Z(\coded(g_X)) \;.
\]
An entirely analogous calculation shows that $V_\bob \, \tilde{W}^j(\tilde{u}) \, V_\bob^\dagger = \Id_{\bob \bob'} \otimes \Big( \tau^W(e_j \tilde{u}) \Big)_{\bob ''}$.

Next, from the self-consistency property of the observables $\tilde{W}^j$ specified in \Cref{lem:qld-construct-the-paulis} we get that for $W \in \{X,Z\}$,
\[
	\E_{j \in \{1,\ldots,t\}} \E_{\tilde{u} \in \F_q^\nqubits} \bra{\hat{\psi}} \left ( \tilde{W}^j(\tilde{u}) \otimes \Id - \Id \otimes \tilde{W}^j(\tilde{u}) \right)^\dagger \left ( \tilde{W}^j(\tilde{u}) \otimes \Id - \Id \otimes \tilde{W}^j(\tilde{u}) \right) \ket{\hat{\psi}} \leq \delta_S\;,
\]
which, since $\tilde{W}^j$ is Hermitian, is equivalent to
\begin{equation}
\label{eq:qld-unitary-1}
\E_{j \in \{1,\ldots,t\}} \E_{\tilde{u} \in \F_q^\nqubits} \bra{\hat{\psi}} \tilde{W}^j(\tilde{u}) \otimes \tilde{W}^j(\tilde{u}) \ket{\hat{\psi}} \geq 1 - \delta_S/2\;.
\end{equation}
Let $\ket{\vartheta} = V_\alice \otimes V_\bob \ket{\hat{\psi}}$, and for $W \in \{X,Z\}$ define the operator
\begin{align*}
	H_W &= \E_{j \in \{1,\ldots,t\}} \E_{\tilde{u} \in \F_q^\nqubits} \tau^W(e_j \tilde{u}) \otimes \tau^W(e_j \tilde{u}) \\
		&= \E_{j \in \{1,\ldots,t\}} \Big ( \E_{s \in \F_q} \tau^W(e_j s) \otimes \tau^W(e_j s) \Big)^{\otimes \nqubits} \\
		&= \Big ( \E_{s \in \F_q} \tau^W(s) \otimes \tau^W(s) \Big)^{\otimes \nqubits}
\end{align*}
where in the second line we used that for all $\tilde{u} \in \F_q^\nqubits$, $\tau^W(\tilde{u}) = \bigotimes_{i = 1}^\nqubits \tau^W(\tilde{u}_i)$ with $\tau^W(\tilde{u}_i)$ being the generalized Pauli $W$ observable acting on the $i$-th qudit. In the third line, we used the fact that for every fixed $j$, we have $\E_{s \in \F_q} \tau^W(e_j s) \otimes \tau^W(e_j s) = \E_{s \in \F_q} \tau^W(s) \otimes \tau^W(s)$. \Cref{eq:qld-unitary-1} is then equivalent to 
\[
	\bra{\vartheta} H_W \ket{\vartheta} \geq 1 - \delta_S/2\;,
\]
for $W \in \{X,Z\}$. 

Observe that
\begin{align*}
	\E_{s,s'} \tau^X(s) \tau^Z(s') \otimes \tau^X(s) \tau^Z(s') &= \E_{s,s'} \sum_{a,b \in \F_q} (-1)^{\tr(as')} \ketbra{a+s}{a} \otimes (-1)^{\tr(bs')} \ketbra{b+s}{b} \\
	&= \E_{s} \sum_{a\in \F_q} \ketbra{a+s}{a} \otimes \ketbra{a+s}{a} \\
	&= \proj{\EPR_q}
\end{align*}
where to go from the first to the second line we used \Cref{lem:cancellation}. This implies that $H_X H_Z = \Big ( \E_{s,s' \in \F_q} \tau^X(s)\tau^Z(s') \otimes \tau^X(s)\tau^Z(s') \Big)^{\otimes \nqubits} = \ketbra{\EPR_q}{\EPR_q}^{\otimes \nqubits}$.
By the triangle inequality, we have
\[
	\left \| H_X \ket{\vartheta} - H_Z \ket{\vartheta} \right\| \leq \left \| \ket{\vartheta} - H_X \ket{\vartheta} \right\| + \left \| \ket{\vartheta} - H_Z \ket{\vartheta} \right\|\;.
\]
Note that for $W \in \{X,Z\}$,
\[
	\left \| \ket{\vartheta} - H_W \ket{\vartheta} \right\|^2 = 1 - 2\bra{\vartheta} H_W \ket{\vartheta} + \bra{\vartheta} H_W^2 \ket{\vartheta} \leq 2 - 2\bra{\vartheta} H_W \ket{\vartheta} \leq \delta_S\;,
\]
where in the inequality we used that $H_W \leq \Id$. Furthermore,
\begin{align*}
	\bra{\vartheta} H_W^2 \ket{\vartheta} &=\| H_W \ket{\vartheta} \|^2\\
	 &= \| \ket{\vartheta} + (H_W - \Id) \ket{\vartheta} \|^2 \\
	 &\geq (\|\ket{\vartheta}\| - \| (H_W - \Id) \ket{\vartheta} \|)^2 \\
	 &\geq (1 - \delta_S)^2 \\
	 &\geq 1 - 2\delta_S\;. 
\end{align*}
Therefore 
\begin{align*}
2\sqrt{\delta_S} &\geq \left \| H_X \ket{\vartheta} - H_Z \ket{\vartheta} \right\|^2  \\
&= \bra{\vartheta} H_X^2 \ket{\vartheta} + \bra{\vartheta} H_Z^2 \ket{\vartheta} - 2\bra{\vartheta} H_X H_Z \ket{\vartheta} \\
&\geq 2 - 4\delta_S - 2\bra{\vartheta} H_X H_Z \ket{\vartheta} \;,
\end{align*}
which implies that 
\[
	\bra{\vartheta} (\Id_{\alice \alice' \bob \bob'} \otimes (\proj{\EPR_q}^{\otimes \nqubits})_{\alice'' \bob''}) \ket{\vartheta} \geq 1 - 2\delta_S - \sqrt{\delta_S} \;.
\]
Define the unnormalized state $\ket{\aux_0}$ on the registers $\alice \alice' \bob \bob'$ as
\[
	\ket{\aux_0}_{\alice \alice' \bob \bob'} = (\bra{\EPR_q}^{\otimes \nqubits})_{\alice'' \bob''} \cdot \ket{\vartheta}_{\alice \alice' \alice'' \bob \bob' \bob''}
\]
and define the normalized state $\ket{\aux} = \frac{1}{\| \, \ket{\aux_0} \, \|} \ket{\aux_0}$. 
The inner product between $\ket{\vartheta}$ and $\ket{\aux} \otimes \ket{\EPR_q}^{\otimes \nqubits}$ can thus be evaluated as
\begin{align*}
	\bra{\vartheta} (\ket{\aux} \otimes \ket{\EPR_q}^{\otimes \nqubits}) &= \frac{1}{\| \, \ket{\aux_0} \, \|} \cdot \bra{\vartheta} (\ket{\aux_0} \otimes \ket{\EPR_q}^{\otimes \nqubits}) \\
	&= \frac{1}{\| \, \ket{\aux_0} \, \|} \cdot \bra{\vartheta} (\Id_{\alice \alice' \bob \bob'} \otimes (\proj{\EPR_q}^{\otimes \nqubits})_{\alice'' \bob''}) \ket{\vartheta} \\
	&= \frac{1}{\| \, \ket{\aux_0} \, \|} \cdot \| \, \ket{\aux_0} \, \|^2 \\
	&= \| \, \ket{\aux_0} \, \| \\
	&\geq \sqrt{1 - 2\delta_S - \sqrt{\delta_S}} \\
	&\geq 1 - O(\sqrt{\delta_S})\;.
\end{align*}
This implies $\| \ket{\vartheta} - \ket{\aux} \otimes \ket{\EPR_q}^{\otimes \nqubits} \|^2 \leq O(\sqrt{\delta_S})$, which establishes the first item of the lemma.

We now establish the second item of the lemma. In what follows, all $\simeq$ and $\approx$ statements hold with respect to the state $\ket{\hat{\psi}}$. From \Cref{lem:qld-win-implications} we have the following implications: for all $W \in \{X,Z\}$ and on average over $u \in \F_q^m$,
	\begin{align}
	\label{eq:qld-unitary-2}
	M^{(\Pauli,W)}_{[g_h(u)=a]} \otimes \Id_{\bob \bob' \bob''} &\simeq_\eps \Id_{\alice \alice' \alice''} \otimes M^{(\Point,W),u}_a\;, \\
	\label{eq:qld-unitary-3}
	\Id_{\alice \alice' \alice''} \otimes M^{(\Point,W),u}_a &\simeq_\eps  M^{(\Point,W),u}_a \otimes \Id_{\bob \bob' \bob''}\;,
	\end{align}
	where $M^{(\Pauli,W)}_{[g_h(u)=a]} = \sum_{h : g_h(u) = a} M^{(\Pauli,W)}_h$. 
From \Cref{lem:qld-construct-the-paulis} we get that
	\begin{equation}
	\label{eq:qld-unitary-4}
		M^{(\Point,W),u}_a \otimes \Id_{\bob \bob' \bob''} \simeq_{\delta_S} \Id_{\alice \alice' \alice''}  \otimes \tilde{M}^{W,\ind_m(u)}_a\;.
	\end{equation}
Using \Cref{fact:triangle-for-simeq} with Equations \eqref{eq:qld-unitary-2}, \eqref{eq:qld-unitary-3}, and \eqref{eq:qld-unitary-4}, we get 
	\begin{equation}
		\label{eq:qld-unitary-5}
		M^{(\Pauli,W)}_{[g_h(u)=a]} \otimes \Id_{\bob \bob' \bob''} \simeq_{\sqrt{\delta_S}} \Id_{\alice \alice' \alice''}  \otimes \tilde{M}^{W,\ind_m(u)}_a\;,
	\end{equation}
	where we used that $\delta_S \geq \eps$. 
	
	Next, we have for all $u \in \F_q^m$ 
	\begin{align}
		&V_\bob \, \tilde{M}^{W,\ind_m(u)}_a \, V_\bob^\dagger \notag \\
		= &\sum_{g_X, g_Z} \Big( \hat{S}_{g_X,g_Z}  \Big)_{\bob \bob'} \otimes \Big ( \tau^X(\coded(g_Z)) \tau^Z(\coded(g_X)) \tau^W_{g_W(u) - a}(\ind_m(u))    \tau^Z(\coded(g_X))^\dagger \tau^X(\coded(g_Z))^\dagger \Big)_{\bob ''} \notag \\
		= &\sum_{g_X, g_Z} \Big( \hat{S}_{g_X,g_Z}  \Big)_{\bob \bob'} \otimes \sum_{h : g_h(u) = g_W(u) - a} \Big ( \tau^X(\coded(g_Z)) \tau^Z(\coded(g_X)) \tau^W_h   \tau^Z(\coded(g_X))^\dagger \tau^X(\coded(g_Z))^\dagger \Big)_{\bob ''} \notag \\
		= &\sum_{g_X, g_Z} \Big( \hat{S}_{g_X,g_Z}  \Big)_{\bob \bob'} \otimes \sum_{h : g_h(u) = g_W(u) - a} \Big ( \tau^W_{h + \coded(g_W)}   \Big)_{\bob ''} \notag \\
	= &\sum_{g_X, g_Z} \Big( \hat{S}_{g_X,g_Z}  \Big)_{\bob \bob'} \otimes \sum_{h : g_h(u) = a} \Big ( \tau^W_{h}   \Big)_{\bob ''} \notag \\
= &\Id_{\bob \bob'} \otimes (\tau^W_{[g_h(u)=a]})_{\bob''}\;,		\label{eq:qld-unitary-6}
	\end{align}
	where in the second line we expanded out the definitions of $V_\bob$ and $\tilde{M}^{W,\ind_m(u)}_a$ (defined in \cref{eq:tilde_M}); in the third line we expanded out the definition of $\tau^W_{g_W(u) - a}(\ind_m(u))$ (defined in \cref{eq:def-tauwu}) and used the fact that for all $h \in \F_q^M$, the inner product $h \cdot \ind_m(u)$ is equal to $g_h(u)$; in the fourth line we used the identities
	\[
		\tau^X(s) \tau^Z_h \tau^X(s)^\dagger = \tau^Z_{h + s} \qquad \text{and} \qquad \tau^Z(s) \tau^X_h \tau^Z(s)^\dagger = \tau^X_{h + s}~;
	\]
	and in the fifth line is due to the following short calculation: for all fixed $g_W$ we have
	\begin{align*}
		\sum_{h : g_h(u) = g_W(u) - a}  \tau^W_{h + \coded(g_W)} = \sum_{h: g_W(u) - g_h(u) = a} \tau^W_{h + \coded(g_W)} = \sum_{h: g_W(u) + g_h(u) = a} \tau^W_{h + \coded(g_W)} = \sum_{h': g_{h'}(u) = a} \tau^W_{h'}~.
	\end{align*}
	The second equality from the fact that $\F_q$ is a field of characteristic $2$ (so subtraction is the same as addition), and the third equality follows from the fact that letting $h' = h + \coded(g_W)$, we have $g_{h'} = g_h + g_W$ (in other words, the low-degree encoding $h \mapsto g_h$ is linear in $h$). 
	
	


Let $\ket{\Delta} = \ket{\aux} \otimes \ket{\EPR_q}^{\otimes \nqubits}$. We show that $V_\alice \, \Big(M^{(\Pauli,W)}_h \Big)_{\alice} \, V_\alice^\dagger \approx_{\delta_\qld} \Id_{\alice \alice'} \otimes (\tau^W_h)_{\alice''}$ (the argument for the operators acting on registers $\bob \bob' \bob''$ proceeds identically). 
	\begin{align}
		&\sum_h \bra{\Delta} \left ( V_\alice \, M^{(\Pauli,W)}_h \, V_\alice^\dagger - (\tau^W_h)_{\alice''}  \right)^\dagger \left ( V_\alice \, M^{(\Pauli,W)}_h \, V_\alice^\dagger - (\tau^W_h)_{\alice''} \right) \ket{\Delta} \notag \\
		&= 2 - 2 \Re \, \sum_h \bra{\Delta} V_\alice\, M^{(\Pauli,W)}_h \, V_\alice^\dagger  \, (\tau^W_h)_{\alice''} \ket{\Delta} \notag \\
		&= 2 - 2 \, \sum_h \bra{\Delta} V_\alice \, M^{(\Pauli,W)}_h \, V_\alice^\dagger  \, \otimes (\tau^W_h)_{\bob''} \ket{\Delta} \label{eq:qld-unitary-7}
\end{align}
	where in the last line we used that $(\tau^W_h)_{\alice''} \ket{\EPR_q}^{\otimes \nqubits} = (\tau^W_h)_{\bob''} \ket{\EPR_q}^{\otimes \nqubits}$.
	We now claim that
	\begin{align}
	&  \, \sum_h \bra{\Delta} V_\alice \, M^{(\Pauli,W)}_h \, V_\alice^\dagger  \, \otimes (\tau^W_h)_{\bob''} \ket{\Delta}  \notag \\
	&\approx_{md/q}  \E_u \sum_{\substack{h, h' : \\ g_h(u) = g_{h'}(u)}} \bra{\Delta} V_\alice \, M^{(\Pauli,W)}_h \, V_\alice^\dagger  \, \otimes (\tau^W_{h'})_{\bob''} \ket{\Delta} \label{eq:qld-unitary-8} \\
	&=  \E_u \sum_{a \in \F_q} \bra{\Delta} V_\alice \, M^{(\Pauli,W)}_{[g_h(u)=a]} \, V_\alice^\dagger \, \otimes (\tau^W_{[g_h(u)=a]})_{\bob''} \ket{\Delta} \notag \\
	&=  \E_u \sum_{a \in \F_q} \bra{\Delta} V_\alice \, M^{(\Pauli,W)}_{[g_h(u)=a]} \, V_\alice^\dagger \, \otimes V_\bob  \, \tilde{M}^{W,\ind_m(u)}_a \, V_\bob^\dagger \ket{\Delta} \label{eq:qld-unitary-9}
	\end{align}
	where the equality in \Cref{eq:qld-unitary-9} follows from the identity in \Cref{eq:qld-unitary-6}, and the approximation in \Cref{eq:qld-unitary-8} follows from the following calculation:
	\begin{align*}
		& \E_u \sum_{\substack{h \neq h' : \\ g_h(u) = g_{h'}(u)}} \bra{\Delta} V_\alice \, M^{(\Pauli,W)}_h \, V_\alice^\dagger  \, \otimes (\tau^W_{h'})_{\bob''} \ket{\Delta}  \\
		&= \sum_{h \neq h'} \left( \E_u \mathbf{1}[g_h(u) = g_{h'}(u)] \right) \bra{\Delta} V_\alice \, M^{(\Pauli,W)}_h \, V_\alice^\dagger  \, \otimes (\tau^W_{h'})_{\bob''} \ket{\Delta} \\
		&\leq \frac{md}{q}\;.
	\end{align*}
	In the last line we used that for for distinct $h \neq h'$, the polynomials $g_h \neq g_{h'}$ and thus by the Schwartz-Zippel lemma can only agree on at most $md/q$ fraction of points $u \in \F_q^m$. 
	
	Continuing from \Cref{eq:qld-unitary-9}, we use the first item of the Lemma which shows that $\| \ket{\Delta} - \ket{\vartheta} \|^2 \leq O(\sqrt{\delta_S})$ and equivalently $\| V_\alice^\dagger \otimes V_\bob^\dagger \ket{\Delta} - \ket{\hat{\psi}} \|^2 \leq O(\sqrt{\delta_S})$, so 
	\begin{align*}
		\text{\cref{eq:qld-unitary-9}} &\approx_{O(\delta_S^{1/4})} \E_u \sum_{a \in \F_q} \bra{\hat{\psi}} M^{(\Pauli,W)}_{[g_h(u)=a]} \otimes \tilde{M}^{W,\ind_m(u)}_a \ket{\hat{\psi}} \\
		&\geq 1 - O(\delta_S^{1/2})\;,
	\end{align*}
        where the last line comes from \eqref{eq:qld-unitary-5}.
	Together with \Cref{eq:qld-unitary-7}, 
	\[
	\text{\cref{eq:qld-unitary-7}} \leq O(\delta_S^{1/4} + md/q)\;.	
	\]
	This concludes the proof of the second item of the lemma.
\end{proof}

We now prove \Cref{thm:pauli-appendix}. 
\begin{proof}[Proof of \Cref{thm:pauli-appendix}]
	Let $\strategy = (\psi,M)$ be a projective strategy for $\game^\pauli_\qldparams$ that succeeds with probability at least $1 - \eps$, where $\ket{\psi}$ is a bipartite state on registers $\alice \bob$. As mentioned at the beginning of \Cref{sec:commutation}, we assume that the state $\ket{\psi}$ is padded with sufficiently many ancilla $\ket{0}$ qubits. Let $\phi_\alice$ be the following isometry mapping the register $\alice$ to the registers $\alice \alice' \alice''$ where $\alice'$ and $\alice''$ are isomorphic to $(\C^q)^{\otimes \nqubits}$:
	\[
		\phi_\alice: \ket{\theta}_\alice \mapsto V_\alice (\ket{\theta}_\alice \otimes (\ket{\EPR_q}^{\otimes \nqubits})_{\alice' \alice''})
	\]
	where $V_\alice$ is the unitary acting on $\alice \alice' \alice''$ from \Cref{lem:qld-unitary}. Define the isometry $\phi_\bob$ analogously. \Cref{lem:qld-unitary} then implies that there exists a state $\ket{\aux}$ on registers $\alice \alice' \bob \bob'$ such that
	\[
		\left \| \phi_\alice \otimes \phi_\bob \ket{\psi} - \ket{\aux} \otimes \ket{\EPR_q}^{\otimes \nqubits} \right \|^2 \leq \delta_\qld
	\]
	and for $W \in \{X,Z\}$,
	\[
		\phi_\alice \, (M^{(\Pauli,W)}_h)_\alice \, \phi_\alice^\dagger \approx_{\delta_\qld} (\tau^W_h)_{\alice''} \qquad \text{and} \qquad \phi_\bob \, (M^{(\Pauli,W)}_h)_\bob \, \phi_\bob^\dagger \approx_{\delta_\qld} (\tau^W_h)_{\bob''}
	\]
	where the $\approx$ statement holds with respect to the state $\ket{\aux} \otimes \ket{\EPR_q}^{\otimes \nqubits}$. This shows that $\game^\pauli_\qldparams$ is a self-test for the strategy $\strategy^\pauli$ with robustness $\delta_\qld(\eps,m,d,q)$. We conclude the proof by making the identification of registers $\alice \alice'$ with $\alice'$ and $\bob \bob'$ with $\bob'$ to be consistent with the statement of \Cref{thm:pauli-appendix}. We also note that, when unpacking the dependence on parameters $\eps, m, d, q$, the function $\delta_\qld(\eps,m,d,q) = O(\delta_S^{1/4} + md/q)$ has the form $a(md)^a(\eps^b + q^{-b} + 2^{-bmd})$ for some universal constants $a > 1, 0 < b < 1$, as desired.
\end{proof}
 

 










		
